\documentclass[12pt,leqno,a4paper]{article}
% Minimalbeispiel für Diplom- oder Studienarbeit in Latex
% OHNE GEWAEHR
% Antje Schweitzer, Juli 2011 & Juli 2014 
% folgende Dateien gehoeren zu diesem Beispiel:
% - thesis.tex (der eigentliche Text)
% - references.bib (die "Datenbank" mit den Literaturverweisen im Bibtex-Format)
% - example.pdf (eine Beispielgrafik)
% damit kann man dann nach der Anleitung unten folgendes Dokument erzeugen:
% - thesis.pdf 

% Minimal example for student papers in LaTeX, can be used as a template
% without warranty
% (English "translation" by Nils Reiter, July 2014)
% You need multiple files for this example to work:
% - thesis.tex (this file)
% - references.bib (a BibTeX file containing the bibliographic entries)
% - example.pdf (an example figure)

% ANLEITUNG -- INSTRUCTIONS
% To create a PDF, please run the following steps:
% - pdflatex thesis.tex
% - bibtex thesis
% - pdflatex thesis.tex
% - pdflatex thesis.tex
% (yes, some steps need to be run multiple times, for reasons)

% Antje Schweitzer, Oct. 2016 - updated statement of authorship
% 
% Antje Schweitzer, Dec 2020 
% - updated translation of statement of authorship ;)
% - made title fit into the window prescribed for Computer Science theses
% please note that for the Computer Science students, the statements
% have slightly different wordings and need to be on the last page
% please check requirements by C.S. department in that case
% (also no guarantee that the window is correct -
% the version I compiled and printed does fit)
% - changed to UTF-8 encoded .tex file, included inputenc utf8, 
% so Umlauts can be typed as usual

\usepackage[hidelinks]{hyperref}
\providecommand{\doi}[1]{\url{https://doi.org/#1}}
\usepackage{natbib}
\usepackage{epsfig}
\usepackage{booktabs}
\usepackage{float}
\usepackage{graphicx}
\usepackage{float}
\usepackage{caption}
\captionsetup[longtable]{font=normalsize}
\usepackage{longtable}
\usepackage{tabularx}
\usepackage{array}
\usepackage{capt-of}
\usepackage{enumitem}
\usepackage{xcolor}
\usepackage{xcolor}
\definecolor{darkgreen}{RGB}{0,100,0}
\definecolor{softgreen}{RGB}{35,120,70}
\definecolor{softorange}{RGB}{190,95,20}
\newcommand{\pres}[1]{\textcolor{softgreen}{#1}}
\newcommand{\change}[1]{\textcolor{softorange}{#1}}
\usepackage{xurl}
\usepackage{amsmath} 
\urlstyle{same}
\usepackage{makecell}
\usepackage[paper=a4paper,left=3cm,right=3cm]{geometry}% http://ctan.org/pkg/geometry
\usepackage[T1]{fontenc}        
\usepackage{alphabeta}
\usepackage[utf8]{inputenc} 
\usepackage[main=english,german,greek]{babel}
\languageattribute{greek}{polutoniko}
\usepackage{tikz}
\usepackage{tikzpagenodes}
\usetikzlibrary{calc}
\usepackage{subcaption}

\renewcommand{\baselinestretch}{1.3}
\parskip = \medskipamount
\frenchspacing
\bibpunct[; ]{(}{)}{;}{a}{,}{;}


\newcommand{\Titel}{Beyond EER: Speaker-Level Evaluation and
Performance Assessment of Voice Anonymization
with Alternative Metrics and Phonetic Features
Analysis}

\begin{document}

% larger left margin for title page to get it centered into C.S. template

\begin{titlepage}
  % different margin so titlepage  into template
  \newgeometry{left=3.5cm,right=2cm}
  \large
   \begin{center}
    Institut für Maschinelle Sprachverarbeitung\\
    Universität Stuttgart\\
    Pfaffenwaldring 5B\\
    D-70569 Stuttgart\\    
    
    \vspace{2.5cm}
    Master Thesis \\
   % \vfill
    {\LARGE \bf \Titel} \\
    \vspace{2cm}
    Ioanna Karagianni\\
       \vfill
    \begin{tabular}[t]{lr}
    Studiengang: & M.Sc. Computational Linguistics \\ % oder: B.Sc. Maschinelle Sprachverarbeitung, M.Sc. Informatik, ...\\
    \\
    \\
    {Prüfer*innen:} & Prof. Dr. Ngoc Thang Vu\\
                    & Dr. Antje Schweitzer\\
     
    {Betreuer:} & Prof. Dr. Ngoc Thang Vu \\ 
    \\
    \\
    {Beginn der Arbeit:} & 22.01.2026\\
    {Ende der Arbeit:} & 22.07.2026\\
    \end{tabular}
  \end{center}
\setlength{\hoffset}{0cm}

  \normalsize
\end{titlepage}

% back to original geometry
\newgeometry{left=3cm,right=3cm}

\newpage
\thispagestyle{empty}

% note that for Computer Science theses (not Maschinelle Sprachverarbeitung or Computational Linguistics)
% this statement has a different wording and needs to be on the last rather than the first page. 
% Please adapt accordingly. 

\begin{otherlanguage}
{german}
\noindent\textbf{Erklärung (Statement of Authorship)}\\


\noindent \noindent Hiermit erkläre ich, dass ich die vorliegende Arbeit selbstständig verfasst und dabei keine anderen als die angegebenen 
Quellen und Hilfsmittel benutzt habe. Sämtliche wörtlichen oder sinngemäßen Übernahmen und Zitate aus anderen Werken sind kenntlich gemacht 
und nachgewiesen. Ich versichere, dass ich keine Hilfsmittel verwendet habe, deren Nutzung die Prüferin oder der Prüfer explizit ausgeschlossen 
hat. Im Anhang „Nutzung von KI-Tools“ habe ich dokumentiert, welche KI-Tools zur Erstellung wesentlicher Teile dieser Arbeit in welcher Art 
und Weise verwendet wurden. Mit Abgabe der Arbeit übernehme ich die volle Verantwortung für alle Inhalte. Die Grundsätze zur KI-Nutzung für 
Studierende an der Universität Stuttgart habe ich zur Kenntnis genommen und befolgt. Weder diese Arbeit noch wesentliche Teile daraus waren 
bisher Gegenstand eines anderen Prüfungsverfahrens. Das elektronische Exemplar stimmt mit allen eingereichten Exemplaren überein.\footnote{Non-binding translation for convenience: I hereby declare that the work presented in this thesis is entirely my own. I did not use any other sources and aids than the listed ones. I have marked all direct or indirect transfers from other sources contained therein as referenced quotations. I assure that I have not used any aids that were explicitly ruled out by the examiner. In the Appendix „Use of AI Tools“ I have documented which AI tools were used and how they contributed in the creation of this work. By submitting this work, I assume full responsibility for all content. I have taken notice of and I comply with the Principles for the Use of Artificial Intelligence by Students at the University of Stuttgart. Neither this work nor significant parts of it were part of another examination procedure. The electronic copy is consistent with all submitted hard copies.}

%% delete the following for theses in German; 
%% delete or keep for English theses. 
%% The German version above must be retained 
%% and signed even if the thesis is written
%% in another language than German.
\end{otherlanguage}
\vspace{2cm}


\noindent Ioanna Karagianni, Stuttgart, July 2026


\newpage
\noindent\textbf{Acknowledgments}\\
\noindent Many people have assisted me and accompanied me during this Master's degree. First of all I would like to thank the IMS for providing me with the opportunity to study in this master's and for creating a warm and inspiring environment for me to develop my technical skills and research interests. I found my calling through this program and developed as a person, researcher, and professional at the same time. Coming from a completely theoretical background, IMS inspired me and gave me the time, material, and space to cultivate my skills and pursue a technical career after exploring many different aspects of natural language and speech processing. I would first like to thank Professor Vu, Thang and Dr. Meyer, Sarina for introducing me to this topic that accompanied me for more than half a year, for their help and guidance through the months of me writing this thesis. Also, the writing of this thesis would not be possible without the support of my family, friends and loved ones. First of all, I would like to deeply thank my parents and my uncle who supported me during my studies, always believed in me and always filled my most difficult moments with their optimism and emotional support. Nothing would have been possible without their continuous unconditional love and encouragement. 

 \begin{otherlanguage}
{greek}
Μαμά, μπαμπά, θείε σας αγαπώ και σας ευχαριστώ για όλα.
\end{otherlanguage}

My deepest gratitude goes to my partner Nikolas, whose love, support, and unwavering faith in me made this journey possible. Thank you for standing by me in every step, reminding me to rest, have a positive mindset, and keep going when things felt impossible.
Also, I would like to thank my friends in Greece and the little family I made in Stuttgart for the beautiful moments we spent together as well as for their support. Thank you all for always hyping me up and believing in me, thank you for offering me a second eye when I needed it the most and most importantly thank you for being my safe space for the last (almost) three years.
Each of you has played a unique role in this journey and I feel deeply grateful for every memory and for having you in my life.

%ABSTRACT -------------------------------------------------------------
\newpage
\thispagestyle{empty} % optional: removes page number/header
 % adjust this to move the block up/down

\begin{center}
\textbf{Abstract}\\[1em] % title in normal size and bold
\begin{minipage}{0.85\textwidth}
\noindent
Voice anonymization systems are conventionally evaluated using the Equal Error Rate (EER), a single population-level threshold that has been shown to obscure substantial variability in privacy protection across individual speakers. 
This thesis reproduces three baseline systems from the 2024 Voice Privacy Challenge (VPC) i.e. STTTS, NAC, and ASRBN and evaluates them not only with global and speaker-level EER, but also with two alternative frameworks: The evidence-based ZEBRA framework (dECE and worst-case LLR computed at the speaker-level) and the legally motivated Linkability \& Singling Out metrics. 
In addition, a feature-level analysis examines whether phonetic, prosodic, acoustic, and linguistic properties preserved after anonymization are associated with individual speaker vulnerability. Results show that while ASRBN achieves the strongest population-level privacy protection under both EER and ZEBRA, a subset of speakers remains consistently re-identifiable across all three architectures, confirming that global metrics can mask persistent individual-level failures.
Linkability \& Singling Out reveal that system rankings shift once evidence is aggregated across utterances: NAC, ranked second under EER and ZEBRA, shows the highest residual Linkability at longer conversation lengths, showing that low verification error does not guarantee resistance to re-identification. The feature analysis provides exploratory evidence that rhythm, phone-duration, and pitch/formant trajectories retain speaker-discriminative structure and show negative descriptive
associations with speaker-level EER in several conditions, while linguistic similarity adds complementary predictive power. These findings support combining population-level, speaker-level, and real-life grounded metrics over global thresholds alone, and point to residual rhythmic, durational, acoustic, and linguistic cues as targets for more robust anonymization.
\end{minipage}
\end{center}


\newpage
\tableofcontents
\cleardoublepage

\listoftables
\addcontentsline{toc}{section}{List of Tables}
\cleardoublepage

\listoffigures
\addcontentsline{toc}{section}{List of Figures}
\cleardoublepage


\begin{tikzpicture}[remember picture,overlay]
\node[
  anchor=south east,
  inner sep=0pt,
  yshift=20mm
] (epigraph) at (current page text area.south east) {%
  \parbox{\textwidth}{\small
    \begin{otherlanguage*}{greek}
    “...λόγον δὲ μόνον ἄνθρωπος ἔχει τῶν ζῴων· ἡ μὲν οὖν φωνὴ τοῦ λυπηροῦ καὶ ἡδέος ἐστὶ σημεῖον... ὁ δὲ λόγος ἐπὶ τῷ δηλοῦν ἐστι τὸ συμφέρον καὶ τὸ βλαβερόν, ὥστε καὶ τὸ δίκαιον καὶ τὸ ἄδικον·.”
    \end{otherlanguage*}

    \medskip
    “…Of all the animals, only the human being possesses \textit{logos} (speech). For \textit{voice} (\textit{phōnē}) is a sign of pleasure and pain ... \textbf{But \textit{logos} exists in order to make clear what is beneficial and what is harmful, and therefore also what is just and what is unjust.} For it is distinctive of human beings, in comparison with the other animals, that they alone have perception of good and evil, of the just and the unjust...”

    \hfill\textit{(Aristotle, \textit{Politics} 1253a10–19)}
  }%
};
\end{tikzpicture}

\clearpage
\newpage
\thispagestyle{empty}
\mbox{This page was intentionally left blank}

\newpage
% INTRODUCTION --------------------------------------------------------
\section{Introduction} In the last years, speech data seem to be in high demand, due to the magnitude of applications they can be utilized for. 
From voice assistants and smart home systems, psychotherapy and medical consultations, to voice communication in IoT \citep{meyer2025usecasesvoiceanonymization} 
the need for big amounts of speech data to train Speech Synthesis and Speech Recognition Systems is high. 
This need also brings some privacy concerns in the way these data are collected, stored and utilized as speech recordings contain personally 
identifiable information and not only explicit identity markers such as the speaker’s low level descriptors (pitch, loudness, jitter etc.), 
but also demographic, biometric, and emotional attributes that can be inferred from the signal itself \citep{nautsch20_interspeech}.

Consequently, the use of raw speech data raises concerns regarding re-identification and unintended inference, 
especially under data protection frameworks such as the \textbf{EU General Data Protection Regulation (GDPR)}, which defines anonymous information as: 

\begin{quote}
``\textit{... anonymous information, namely information which does not relate to an identified or identifiable natural person or to personal data rendered anonymous in such a manner that the data subject is not or no longer identifiable.}'' \citep{noauthor_recital_nodate}
\end{quote}

This anonymization does not only require the alteration of the speaker's voice, in order to make them non-identifiable, 
but also the removal or the change of specific words in an utterance and sounds in the background \citep{tomashenko_voiceprivacy_2022-1}. 
Therefore, to address these risks and to develop systems that comply with the privacy laws, the \textbf{Voice Privacy Challenge (VPC)} series (2020–2024) has emerged \citep{tomashenko_voiceprivacy_2022-1, tomashenko_voiceprivacy_nodate, tomashenko2024voiceprivacy2024challengeevaluation}. 
\textbf{Voice Privacy} comprises the set of transformations applied to a speech signal to remove, mask, or obfuscate speaker-specific and sensitive 
cues (such as gender, accent, emotion, or health signals) that enable re-identification or inference of the speaker, while maintaining utility of 
the speech (e.g. paralinguistic, linguistic content) for downstream tasks \citep{champion_anonymizing_2024}. The goal of the challenge is to support the participants to advance and evaluate techniques for voice anonymization \citep{tomashenko_voiceprivacy_2022-1, VoicePrivacy2024}. 
The term \textbf{Voice Anonymization} refers to the transformation of a speech signal such that the speaker's identity is concealed while at the 
same time the linguistic content and the paralinguistic attributes of the utterance are preserved (the so-called: \textbf{Privacy-Utility tradeoff}) \citep{tomashenko2024voiceprivacy2024challengeevaluation}. 
% The three basic components of the VPC are: the proposed \textbf{Baseline Systems}, the \textbf{Attacker Models} (Automatic Speaker Verification - ASV systems), which aim to de-anonymize speech and identify the speakers, and lastly, \textbf{Evaluation Metrics} which measure the privacy and the utility of the anonymized speech. 

During the last years, there have been remarkable developments in the anonymization systems developed by the VPC organisers and the participants \citep{panariello_voiceprivacy_2024, VoicePrivacy2024, tomashenko_voiceprivacy_2022-1} 
and a transition from traditional models in the first two challenges \citep{mawalim_x-vector_2020, gupta_design_nodate} to neural, more advanced ones 
in the latest VPC of 2024 \citep{leang_exploring_nodate, hua_emotional_2024}. 
However, the development of these systems underlines that anonymization is not only a technical challenge but also a question of how privacy 
should be quantified, measured and safeguarded in real-life scenarios. Ever since the beginning of the Challenge in 2020, doubts are openly 
expressed regarding the credibility of the standard privacy metric utilized consistently across all the VPCs, which is 
the \textbf{Equal Error Rate} (\textbf{EER}). These doubts were centered around the inability of this metric to capture and express the 
variability of privacy protection across the individual speakers, as it is noticed that the anonymization performance of some individual speakers 
remains very low across fundamentally different anonymization systems \citep{gaznepoglu2025sayexploitinglinguisticcontent, williams_anonymizing_2024,nautsch20_interspeech}. 
These limitations of the EER have motivated speech engineers to explore alternative or complementary evaluation metrics, 
like \textbf{Linkability \& Singling Out} \citep{vauquier25_interspeech} grounded in data protection laws and 
the \textbf{Zero Evidence Biometric Recognition (ZEBRA) framework} inspired by forensics \citep{nautsch20_interspeech, 10.1016/j.csl.2021.101299}.

Regardless of the need for different evaluation metrics, the phenomenon of some speakers performing low in the EER evaluation, 
even though different anonymization systems have been utilized, persists and requires further attention and investigation 
\citep{gaznepoglu2025sayexploitinglinguisticcontent, williams_anonymizing_2024}. 
More specifically, there seems to exist the hypothesis that there are speakers that are easier to anonymize (resulting in higher EER scores) and 
speakers that are harder to anonymize (resulting in very low EER scores). This trend can be attributed to a number of factors, like the linguistic 
content of the speaker's speech and the temporal dynamics of their speech \citep{tomashenko_exploiting_2025, tomashenko_analysis_2025}. 
Despite these first explanation efforts on the analysis of the individual speakers' performance, a collective targeted work providing an in-depth study of the prosodic, acoustic and phonetic characteristics is yet 
to be performed.

In this thesis, the primary focus falls on the examination of alternative evaluation metrics for assessing privacy in voice anonymization systems, as proposed in recent literature and on
providing a deeper understanding of how these metrics capture privacy preservation compared to the traditional EER used in the VPC 2024. 
Moreover, an effort is made to investigate the phonetic, acoustic and prosodic factors and their combination with linguistic factors at the speaker level, 
that may affect negatively the effectiveness of anonymization. 
To address these initial hypotheses, a series of experiments is conducted using the baseline systems, the attacker model and evaluation pipeline provided by VPC 2024 \citep{tomashenko2024voiceprivacy2024challengeevaluation}, in combination with alternative evaluation frameworks and metrics introduced by the earlier and latest research in the field.

Therefore, a gap can be identified from which four main research questions (RQs) seem to emerge, offering a ground for further experimentation and research:
\begin{enumerate}[label=\textbf{RQ\arabic*:}, leftmargin=*]
    \item \textbf{How well or poorly do the three best-performing anonymization systems in the Voice Privacy Challenge of 2024 (B3, B4, B5) preserve privacy across individual speakers, and where are their differences noticed?}

    \item \textbf{How do other existing metrics (Linkability \& Singling Out, ZEBRA) seem to perform after evaluating the same baseline systems, for the same speakers?}

    \item \textbf{Do some other factors, such as gender, phoneme duration per speaker, pitch, rhythm, formant trajectories, and the linguistic content of speech, seem to affect the anonymization systems' performance?}

    \item \textbf{What could be some weaknesses of the latest Voice Privacy Challenge anonymization systems, based on the findings of the evaluation metrics used and the feature analysis conducted, and what should be taken into consideration to make them more robust?}
\end{enumerate}

To answer these questions, this thesis reproduces the STTTS, NAC, and ASRBN baselines of VPC 2024 and evaluates them using global and speaker-level EER (RQ1), 
the ZEBRA framework, and the Linkability \& Singling Out metrics (RQ2). In addition, phonetic, prosodic, acoustic, embedding-based, and linguistic features as well as the role of gender (RQ3) are analyzed to 
examine which preservation properties remain speaker-discriminative after anonymization and whether they are associated with speaker vulnerability. Lastly, the findings of the evaluation metrics and the feature 
analysis are discussed to identify weaknesses of the anonymization systems and provide recommendations for future implementations (RQ4).

The results demonstrate the importance of evaluating voice anonymization from multiple perspectives. 
Although the systems provide meaningful privacy protection at the population level, some speakers remain consistently vulnerable across different 
anonymization architectures. Moreover, the system rankings depend on the evaluated threat, particularly when several utterances from the same speaker are 
aggregated. The feature analysis further indicates that rhythm, phone-duration patterns, pitch and formant trajectories, as well as linguistic similarity with the enrollment utterances may 
preserve residual speaker information, although their contribution varies across systems and trial partitions. 
These findings support the use of alternative population-level and speaker-level frameworks that underline different aspects of leakage rather than relying on a global threshold-based metric alone.

The remainder of this thesis is organized as follows: Section \ref{sec:background_work} introduces the VPC and its components. Section \ref{sec:related_work} reviews related work on the limitations of EER, alternative privacy metrics, and residual phonetic and linguistic identity cues. 
Section \ref{sec:methodology} describes the experimental methodology, followed by the results in Section \ref{sec:results} and their discussion in Section \ref{sec:discussion}. Section \ref{sec:conclusions} concludes the thesis and outlines directions for future research. The code used to reproduce the results
reported in this thesis is available in the following GitHub repository:
\par\noindent
\href{https://github.com/joannakarayianni/Voice-Privacy-Challenge-2024-Thesis}{\texttt{https://github.com/joannakarayianni/Voice-Privacy-Challenge-2024-Thesis}}

%  \section{Background Work}

% \subsection{The Voice Privacy Challenge Initiative}
% The integration of speech systems in numerous services and industries, like telecommunications, banking and public services, healthcare, the workplace (for meeting recordings), language learning applications, and social media, has made speech data a valuable form of information needed to enhance and train more reliable speech systems in order to elevate the user experience \citep{meyer2025usecasesvoiceanonymization}. Yet, these speech data are also a sensitive source of information, as they inherently encode individually identifiable traits like sensitive cues about the speaker's identity, their gender, age and health indicators. The regulatory framework, GDPR which emerged in 2016 \citep{noauthor_regulation_2016} recognises biometric data used for the purpose of uniquely identifying a natural person as a special category of sensitive personal data (Article 9.1), whose processing is in principle prohibited unless specific legal exceptions apply (according to Article 9.2). Speech data inherently falls within the scope of sensitive data because, as previously mentioned, they reflect a combination of biological and behavioural characteristics. On a technical side, the risk of privacy intrusion arises from the very moment speech is captured and persists throughout its processing and storage lifecycle. Frontend and backend speech processing systems can derive linguistic, paralinguistic, and extralinguistic information from speech representations, many of which are sensitive and personal in nature. The multifunctionality of speech data increases the risk of "function creep", where data collected for one purpose, e.g. speech recognition, may later be repurposed for identity recognition or profiling without appropriate legal basis \citep{nautsch_gdpr_2019}. 


% Therefore, these challenges underline the need for approaches that limit the exposure of sensitive speaker information. In this context, speaker anonymization emerged as a key technical strategy aimed at reducing the identifiability of speakers while allowing speech data to remain usable for downstream tasks (Article 5.1 - "data minimization"), motivating the initiative of Voice Privacy Challenge in 2020. The goal of the initiative is to bring together researchers, speech engineers and privacy professionals, and provide them with the means to develop their methodologies in building state of the art systems for voice anonymization, as well as investigating evaluation metrics that would be meaningful in the comparison of the different models \citep{tomashenko_voiceprivacy_2022-1}. 


% The organisers of the challenge are describing the privacy preservation as a "game" between the users who share their data to be used for legitimate downstream tasks (e.g. human communication, automated processing, model training) and an attacker who attempts to infer the speaker's identity from the shared speech data. Within this framework, an anonymization system is applied to the speech signal, while the Automatic Speaker Verification (ASV) system that serves as the attacker is assumed to have information about the anonymization mechanism (depending on the different attacker setup defined per each challenge version) and attempts to re-identify the speaker. The organising committee of the challenge, in each of the challenge's versions, provide the participants with common baseline speech anonymization systems (with most of the times fundamentally different architectures) to be enhanced, fixed training, development and evaluation datasets, as well as standard evaluation protocols to allow for a fair comparing manner. As the goal is the development of the best anonymization system, the systems are evaluated by jointly considering privacy as well as the utility of the speech as the linguistic and paralinguistic value of the data in terms of clear speech or emotion should be preserved for usage on the training of various speech systems. The different anonymization tasks per challenge version are presented in more detail in Table~\ref{tab:vpc_anonymization_tasks} according to their official corresponding Evaluation Plans by \cite{tomashenko_voiceprivacy_nodate, tomashenko2022voiceprivacy2022challengeevaluation, tomashenko2024voiceprivacy2024challengeevaluation}. All in all, through this setup, the VoicePrivacy Challenge has established a shared experimental framework that advances the development of voice anonymization techniques while aligning technical progress with real-world privacy requirements.


% \small
% \renewcommand{\arraystretch}{1.12}
% \setlength{\tabcolsep}{6pt}
% \setlength{\parskip}{0pt}

% \begin{longtable}{p{2.2cm} p{3.5cm} p{7.3cm}}

% \toprule
% \parbox[t]{2.2cm}{\centering\textbf{VPC}\par\textbf{Edition}} & \textbf{Anonymization Scope} & \textbf{Task Requirements and Evaluation} \\
% \midrule
% \endfirsthead

% \toprule
% \parbox[t]{2.2cm}{\centering\textbf{VPC}\par\textbf{Edition}} & \textbf{Anonymization Scope} & \textbf{Task Requirements and Evaluation} \\
% \midrule
% \endhead

% \bottomrule
% \endfoot

% \textbf{VPC 2020}
% & Speaker-level anonymization
% &
% 1. Output a speech waveform. \par
% 2. Conceal speaker identity. \par
% 3. Preserve speech characteristics. \par
% 4. Ensure all trial utterances from the same speaker are mapped to the same pseudo-speaker. \par
% 5. Utility assessed via ASR on clean data. \par
% 6. Voice distinctiveness assessed in post-evaluation. \par
% 7. Subjective intelligibility and naturalness included.
% \\

% \textbf{VPC 2022}
% & Speaker-level anonymization
% &
% 1. Output a speech waveform. \par
% 2. Conceal speaker identity. \par
% 3. Preserve linguistic content and paralinguistic attributes. \par
% 4. Ensure consistent pseudo-speaker assignment per speaker across utterances. \par
% 5. Utility assessed via ASR on anonymized data. \par 6. Pitch preservation assessed using $\rho_{F_0}$. \par 7. Voice distinctiveness assessed using GVD. \par
% 8. Subjective intelligibility and naturalness included.
% \\

% \textbf{VPC 2024}
% & Utterance-level anonymization
% &
% 1. Output a speech waveform. \par
% 2. Conceal speaker identity independently for each utterance. \par
% 3. Preserve linguistic, paralinguistic and emotional content. \par
% 4. Utility assessed via ASR using WER. \par
% 5. Utility assessed via SER using UAR. \par
% \\
% \end{longtable}
% \setcounter{table}{0}
% \captionof{table}{Comparison of anonymization task specifications across the Voice Privacy Challenge editions (2020--2024).}
% \label{tab:vpc_anonymization_tasks}

% \normalsize

\section{Background Work}
\label{sec:background_work}

\subsection{The Voice Privacy Challenge Initiative}
The integration of speech systems into numerous services and industries has made speech data a valuable resource for 
training reliable speech systems \citep{meyer2025usecasesvoiceanonymization}. Yet speech data is also sensitive, as it inherently encodes 
cues about a speaker's identity (e.g. gender, age, health indicators). The GDPR, which emerged in 2016 \citep{noauthor_recital_nodate}, 
recognises biometric data used to uniquely identify a natural person as a special category of sensitive personal data (Article 9.1), 
whose processing is in principle prohibited unless specific legal exceptions apply (Article 9.2). For speech data the risk of privacy intrusion arises from the moment speech is captured 
and persists throughout its processing and storage lifecycle. Moreover, the multifunctionality of speech data increases the risk of "function creep," 
where data collected for one purpose (e.g. speech recognition) may later be used for identity recognition or profiling without
legal basis \citep{nautsch_gdpr_2019}.

Speaker anonymization emerged from these challenges as a key technical strategy to reduce speaker identifiability while keeping speech data usable for downstream tasks (Article 5.1 - "data minimization"), 
motivating the VPC initiative in 2020. Its goal is to bring together researchers, speech engineers, and privacy professionals to develop 
state-of-the-art anonymization systems \citep{tomashenko_voiceprivacy_2022-1}.

The organisers describe privacy preservation as a "game" between users who share data for legitimate downstream tasks and an attacker who tries 
to infer speaker identity from that data. An anonymization system is applied to the speech signal, while an Automatic Speaker Verification (ASV) 
system acts as the attacker, assumed to have knowledge of the anonymization mechanism depending on the attacker setup defined per challenge version. 
The organising committee provides common baseline systems, fixed training/development/evaluation datasets, and standard evaluation protocols for 
fair comparison. Systems are evaluated jointly on privacy and utility, since the linguistic and paralinguistic value of the data must be preserved. 
VPC has therefore established a shared experimental framework in an effort to align technical progress with real-world privacy requirements.

% As previously mentioned the integration of speech systems in numerous services and industries, 
% has made speech data a valuable form of information needed to enhance and train more reliable speech systems 
% \citep{meyer2025usecasesvoiceanonymization}. Yet, these speech data are 
% also a sensitive source of information, as they inherently encode sensitive cues about the speaker's identity (like gender, age and health indicators). 
% The regulatory framework, GDPR which emerged in 2016 \citep{noauthor_recital_nodate} recognises biometric data used for the purpose of uniquely 
% identifying a natural person as a special category of sensitive personal data (Article 9.1), whose processing is in principle prohibited unless specific legal exceptions apply (Article 9.2). 
% Speech data falls within the scope of sensitive data as it reflects a combination of 
% biological and behavioural characteristics, and the risk of privacy 
% intrusion arises from the very moment speech is captured and persists 
% throughout its processing and storage lifecycle. Also, the multifunctionality 
% of speech data further increases the risk of "function creep", where 
% data collected for one purpose, e.g. speech recognition, may later be 
% repurposed for identity recognition or profiling without appropriate 
% legal basis \citep{nautsch_gdpr_2019}.

% These challenges underline the need for approaches that limit the 
% exposure of sensitive speaker information. Speaker anonymization emerged 
% as a key technical strategy aiming to reduce the identifiability of 
% speakers while allowing speech data to remain usable for downstream 
% tasks (Article 5.1 - "data minimization"), motivating the VPC initiative in 2020. The goal of the initiative is to bring 
% together researchers, speech engineers and privacy professionals to 
% develop state-of-the-art anonymization systems and investigate 
% meaningful evaluation metrics for their comparison 
% \citep{tomashenko_voiceprivacy_2022-1}.

% The organisers describe privacy preservation as a "game" between users 
% who share their data for legitimate downstream tasks and an attacker who 
% attempts to infer the speaker's identity from the shared speech data. An 
% anonymization system is applied to the speech signal, while an Automatic 
% Speaker Verification (ASV) system serves as the attacker, assumed to 
% have information about the anonymization mechanism depending on the 
% attacker setup defined per challenge version. The organising committee 
% provides participants with common baseline anonymization systems, fixed 
% training, development and evaluation datasets, and standard evaluation 
% protocols allowing for fair comparison. Systems are evaluated by jointly 
% considering privacy and utility, as the linguistic and paralinguistic 
% value of the data should be preserved for downstream tasks. 
% VPC has therefore established a shared experimental 
% framework that aims to advance voice anonymization techniques while aligning 
% technical progress with real-world privacy requirements.


% anonymization tasks per challenge version are presented in 
% Table~\ref{tab:vpc_anonymization_tasks} according to their official 
% Evaluation Plans \citep{tomashenko_voiceprivacy_nodate, 
% tomashenko2022voiceprivacy2022challengeevaluation, 
% tomashenko2024voiceprivacy2024challengeevaluation}. Through this setup, 


% \subsection{Baseline Anonymization Systems}
% As previously mentioned, the organisers of the challenge provide the participant teams with baseline anonymization systems, which have fundamentally different architectures among them. Based on these baselines the teams are expected to enhance their algorithms in order to optimize the privacy-utility tradeoff under the common evaluation protocol per challenge. To understand better the nature of the baselines and the role they play in system development, a review of the baseline systems per challenge follows. 

% \subsubsection{Voice Privacy Challenge 2020}
% In the first edition of the challenge in 2020 the baselines where 2, namely Baseline-1 (B1) and Baseline-2 (B2). 

% \textbf{B1} is based on x-vector anonymization with neural speech synthesis. This entitles a three stage pipeline including: Feature Extraction, X-vector Anonymization and Speech Synthesis, inspired from neural TTS systems.  For the first step, the speaker identities were represented as 512-dimansional x-vectors and were extracted using a Time Delay Neural Network (TDNN) speaker embedding model trained on VoxCeleb 1 and 2. The linguistic content was represented as 256-dimensional bottleneck (BN) features extracted through a Factorized Time Delay Neural Network (TDNN-F) ASR acoustic model which was trained on LibriSpeech train-clean-100 and train-other-500. Lastly, F0 was extracted with a standard pitch tracker. Then, for each source speaker a new anonymized x-vector is computed from averaging a set of candidate x-vectors from a speaker pool, using Probabilistic Discriminant Analysis (PLDA). The resulting vector defines the corresponding pseudo-speaker of the source speaker. The Speech Synthesis step, includes the acoustic model which synthesizes the Mel-spectrograms from the anonymized x-vector, the original BN features and lastly the original F0 \citep{ tomashenko_voiceprivacy_nodate}.

% \textbf{B2} is based on traditional signal processing techniques and employs the McAdams coefficient. Initially, the speech signal is decomposed using Linear Predictive Coding (LPC) into LPC coefficients and a residual signal, with the residual left unchanged. The LPC coefficients are then converted into pole positions via polynomial root-finding. Subsequently, the McAdams transformation is applied to the angles of the complex-conjugate poles, while real poles remain unchanged. Finally, resynthesis is performed by converting the modified poles back to LPC coefficients and combining them with the original residual signal to generate a new speech frame in the time domain \citep{ tomashenko_voiceprivacy_nodate}.

% In this first version of the challenge from the 25 teams who submitted, 16 submissions were considered successful. Most of the submitted systems retained the B1 pipeline with modifications on its modules. From the results of this challenge, it can be deduced that neural speech synthesis with x-vectors offered better anonymization results across the different attacker models utilized in this challenge and the note that results on EER seem to be consistent with the results in \( C_{\mathrm{llr}} \).However, the privacy gain that is noticed, the anonymization system's performance is accompanied by a degradation in utility, as WERs increased significantly depending on the dataset and system configuration, giving a first important indication that privacy levels are sensitive to the attacker model \citep{ tomashenko_voiceprivacy_pdf_nodate, tomashenko_voiceprivacy_2022-1}.

% \subsubsection{Voice Privacy Challenge 2022}
% The second VoicePrivacy Challenge in 2022, retained the anonymization paradigm introduced in 2020 and enhanced the already existing baseline systems, which now are \textbf{B1.a} and \textbf{B1.b} for the anonymization using x-vectors and neural waveform models and \textbf{B2} for the McAdams coefficient, which was left unchanged in its core architecture \citep{tomashenko2022voiceprivacy2022challengeevaluation}.

% \textbf{B1.a} keeps the VPC 2020 design in the speech synthesis backend (with the acoustic model predicting the mel-filterbanks, which then get passed on the Neural Source-Filter -NSF vocoder), while the \textbf{B1.b } version converts the
% BN, F0, and x-vector features directly using an NSF model and adopts a HiFi-GAN vocoder \citep{tomashenko2022voiceprivacy2022challengeevaluation}.

% In the second version of the challenge from the 43 teams that participated, the teams that submitted results were 6 and the valid submitted systems were 16. Once again 14 systems adopted B1 as their core system architecture, while 2 systems utilized B2. The teams in this challenge seem to perform a change to almost all the components of B1. According to the results reports,     most submitted systems outperformed the baselines achieving higher EER  scores than unprotected speech under the semi-informed attacker model, while maintaining competitive WERs. Nevertheless, no system achieved consistently optimal performance across all evaluation metrics. For example, systems that showcased stronger privacy protection (i.e.\ EERs $> 30\%$) tended to show increased degradation in utility and voice distinctiveness, as reflected by higher WER and more negative gain of voice distinctiveness (GVD) values. On the other hand, systems that better preserved intelligibility and naturalness generally provided weaker anonymization. \citep{ panariello_voiceprivacy_2024, tomashenko_voiceprivacy_2022}

% \subsubsection{Voice Privacy Challenge 2024}
% Lastly, in the latest version of the VPC 2024 the baseline systems are counting up to 6. More specifically, \textbf{B1} of VPC 2024 adopts the same system architecture as B1.a from VPC 2022, and \textbf{B2} of VPC 2024, retains the same system components as B2 from 2022.

% A new addition \textbf{B3}, represents a conceptual shift from voice conversion toward ASR+TTS-style anonymization. Speech is factorized into discrete linguistic units (phonetic transcriptions) and explicit prosodic attributes and as a second step a Wasserstein GAN generates artificial pseudo-speaker embeddings that are explicitly dissimilar from the original speaker embedding via a cosine threshold. In parallel, pitch and energy are randomly perturbed at the phone level, which aims to suppress fine-grained speaker specific prosodic patterns while preserving global prosodic structure \citep{tomashenko2024voiceprivacy2024challengeevaluation}.

% \textbf{B4} represents an anonymization system using Neural Audio Codec (NAC) language modeling and treats the anonymization task as a conditional sequence generation problem, rather than a transformation of speaker representations. B4 assumes that speaker identity is implicitly encoded in acoustic token sequences, and anonymization is achieved by conditioning generation on semantic tokens from the source utterance and acoustic tokens sampled from a pseudo-speaker pool. This systems main contribution is that anonymization is performed in token space, using a GPT-like decoder to generate new acoustic tokens that satisfy both semantic consistency and pseudo-speaker style constraints \citep{tomashenko2024voiceprivacy2024challengeevaluation}.

% Baselines \textbf{B5} and \textbf{B6} represent anonymization using ASR-BN with Vector Quantization (VQ), aiming at reducing speaker leakage at the feature level and forcing the quantized bottleneck representations to lie in a finite codebook. By quantizing BN features, B5 and B6 aim to improve disentanglement between linguistic content and speaker characteristics.  The distinction between the two baselines lies in the ASR acoustic model used for feature extraction, as B5 relies on a pretrained wav2vec2 model followed by TDNN-F layers, whereas B6 uses a fully TDNN-F-based architecture trained from scratch. This change tries to assess whether self-supervised pretraining introduces stronger speaker leakage compared to supervised acoustic modeling \citep{tomashenko2024voiceprivacy2024challengeevaluation}.

% In this version of the challenge, 40 teams were registered, with 13 ultimately submitting results, leading to a total of 36 submitted anonymization systems, with the challenge results offering some influential contributions on the state of the art systems of voice anonymization.\citep{tomashenko_third_2026}. More specifically,  \cite{yao2025npuntuvoiceprivacy2024} adopted a disentangled neural codec architecture combined with a step-by-step disentanglement strategy to explicitly separate speaker identity, time-variant linguistic content, and emotional state. On the testing set, this system achieved an EER of 36.99\%, while maintaining a WER of 1.85\% and an UAR of 64.48\%, significantly improving the privacy-utility-emotion trade-off compared to earlier baseline and participant systems. Additionally, \cite{tan_system_nodate} propose a system with explicit sentiment transfer and feature interpolation, aiming to jointly address speaker privacy preservation and emotional utility. Their system utilizes WavLM to extract raw speech sequence features, and combines these with randomly selected features from an anonymization pool through a feature interpolation mechanism that controls the degree of speaker identity suppression while preserving emotion related characteristics. Sentiment embeddings are explicitly incorporated during the speech synthesis process to ensure that emotional cues are being retained after anonymization. Their results on the testing setup mark EER scores of 34.12\% for female speakers, 36.08\% for male speakers, a 2.37\% WER score, and a 60.95\% UAR, achieving in that way the best trade-off of privacy protection and linguistic content preservation in the 1st, 2nd and 3rd privacy categories \citep{VoicePrivacy2024,tomashenko_third_2026}.

\subsection{Baseline Anonymization Systems}
Teams are provided with baseline anonymization systems of various architectures to enhance for optimizing the privacy-utility 
tradeoff under the common evaluation protocol. 
The following sections review these baselines, focusing on VPC 2024.

\subsubsection{Voice Privacy Challenge 2020}
The first edition introduced two baselines. \textbf{B1} is based on x-vector 
anonymization with neural speech synthesis, following a three stage 
pipeline of feature extraction, x-vector anonymization, and speech 
synthesis. Speaker identities are represented as x-vectors extracted via a TDNN model, linguistic content as 
bottleneck (BN) features from a TDNN-F ASR model, 
and F0 via a standard pitch tracker. The new anonymized x-vector gets 
computed by averaging candidate x-vectors from a speaker pool using 
Probabilistic Discriminant Analysis (PLDA), defining a pseudospeaker, after which an acoustic model 
synthesized Mel-spectrograms from the anonymized x-vector, original 
BN features and original F0 \citep{tomashenko_voiceprivacy_nodate}. 
\textbf{B2} uses traditional signal processing via the McAdams coefficient, 
decomposing speech with Linear Predictive Coding (LPC) and modifying the angles of complex-conjugate poles before resynthesis 
\citep{tomashenko_voiceprivacy_nodate}. Of the 25 submitted systems, 
16 were successful, with B1 offering better anonymization 
results, though at the cost of utility degradation \citep{panariello_voiceprivacy_2024, tomashenko_voiceprivacy_2022-1}.

\subsubsection{Voice Privacy Challenge 2022}
VPC 2022 retains and enhances the 2020 paradigm. \textbf{B1.a} kept the 
original synthesis backend with a Neural Source Filter (NSF) vocoder, while \textbf{B1.b} converted 
BN, F0, and x-vector features directly using an NSF model with a 
HiFi-GAN vocoder \citep{tomashenko2022voiceprivacy2022challengeevaluation}. 
\textbf{B2} remained unchanged. Among 16 valid submissions from 6 teams, most 
outperform the baselines under the semi-informed attacker, but no 
system achieved consistently optimal performance as stronger privacy 
protection consistently came at the cost of utility and voice 
distinctiveness \citep{panariello_voiceprivacy_2024, 
tomashenko_voiceprivacy_2022}.

\subsubsection{Voice Privacy Challenge 2024}
VPC 2024 introduced six baselines, with \textbf{B1} and \textbf{B2} being retained from 2022. 
The three new baselines evaluated in this thesis represent 
fundamentally different anonymization paradigms. \textbf{B3} \citep{meyer_prosody_2023} marks a 
shift toward ASR+TTS style anonymization: Speech is factorized into 
discrete phonetic transcriptions and prosodic attributes, a Wasserstein 
GAN generates pseudospeaker embeddings explicitly dissimilar from the 
original via a cosine threshold, and pitch and energy are randomly 
perturbed at the phone level to suppress fine-grained speaker-specific 
prosodic patterns while preserving the global prosodic structure 
\citep{tomashenko2024voiceprivacy2024challengeevaluation}.
Then, \textbf{B4} treats anonymization as a conditional sequence generation problem in token 
space: Semantic tokens from the source utterance are concatenated with 
acoustic tokens from a pseudospeaker pool, and a GPT-like decoder 
generates new acoustic tokens satisfying both semantic consistency and 
pseudospeaker style constraints 
\citep{tomashenko2024voiceprivacy2024challengeevaluation}. Lastly, \textbf{B5} and \textbf{B6} 
apply vector quantization (VQ) to Automatic Speaker Recognition (ASR) bottleneck features to reduce 
speaker leakage at the feature level, with B5 using a pretrained 
wav2vec2 model followed by Time Delay Neural Network (TDNN-F) layers and B6 using a fully 
supervised TDNN-F architecture, allowing assessment of whether 
self-supervised pretraining introduces stronger speaker leakage 
\citep{tomashenko2024voiceprivacy2024challengeevaluation}.

Of the 40 registered teams, 13 submitted a total of 36 systems, 
yielding influential contributions
\citep{tomashenko_third_2026}. \cite{yao2025npuntuvoiceprivacy2024} adopted a 
disentangled neural codec with step-by-step disentanglement of speaker 
identity, linguistic content, and emotional state, achieving an EER of 
36.99\%, WER (Word Error Rate) of 1.85\%, and UAR (Unweighted Average Recall) of 64.48\% on the test set. 
\cite{tan_system_nodate} proposed a system combining WavLM feature 
extraction with feature interpolation and explicit sentiment transfer, 
achieving EER scores of 34.12\% (female) and 36.08\% (male), a WER of 
2.37\%, and a UAR of 60.95\%, obtaining the best privacy-utility-emotion 
tradeoff across the top three privacy categories 
\citep{VoicePrivacy2024, tomashenko_third_2026}.

%%%%%%%%%%%% ATTACKER SYSTEMS %%%%%%%%%%%%%%%
\subsection{Attacker Systems}
The second fundamental component of the VPC is the attacker system, 
in the form of an ASV model, which constitutes the threat scenario 
under which anonymization systems are evaluated. Its goal is to 
re-identify the original enrollment speaker from the anonymized trials.

VPC 2020 defines three attacker models depending on the adversary's 
level of knowledge \citep{tomashenko_voiceprivacy_nodate, 
tomashenko_voiceprivacy_2022-1}: the \textbf{ignorant attacker}, 
unaware that anonymization was applied, who verifies anonymized trial 
utterances against original enrollment utterances using an ASV system 
trained on original data. Then the \textbf{lazy-informed attacker}, who 
knows the anonymization method but not the pseudospeaker mapping, 
and operates on anonymized enrollment and trial utterances with an 
ASV system trained on original data, and lastly the \textbf{semi-informed 
attacker}.

The semi-informed attacker constitutes the strongest threat model and 
was adopted as the default setting in VPC 2022 and VPC 2024 
\citep{tomashenko_voiceprivacy_pdf_nodate, 
tomashenko_voiceprivacy_2022-1}. In both editions, the attacker has 
full knowledge of the anonymization system and access to multiple 
enrollment utterances per speaker, anonymizes enrollment data using 
the same pipeline as the trial data, and retrains the ASV system 
accordingly. The only difference in VPC 2024 is that anonymization 
operates at utterance level, meaning pseudospeaker assignment is 
performed independently per utterance rather than per speaker 
\citep{tomashenko2024voiceprivacy2024challengeevaluation}.

% %%%%%%%%%%%% EVALUATION METRICS %%%%%%%%%%%%%%%
% \subsection{Evaluation Metrics for the Voice Privacy Challenge}
% In this section the evaluation metrics used so far in the assessment of the performance of the voice anonymization systems will be presented and discussed thoroughly.

% % As mentioned in the previous section the Voice Anonymization Systems aim to balance \textbf{privacy} (meaning the concealment of the speaker utterance) and \textbf{utility} (meaning the linguistic content presentation for using the anonymized speech for training speech systems in downstream tasks). 

% \subsubsection{Voice Privacy Challenge 2020} \label{subsec:vpc2020}
% Starting off, from the first Voice Privacy Challenge in 2020 \citep{tomashenko_introducing_2020, tomashenko_voiceprivacy_2022, tomashenko_voiceprivacy_nodate}, where the evaluation of the models (like also  in the following challenges) is assessed by using ASV based metrics, because the system should be able to make the attacker incompetent to re-identify the anonymized utterance's enrollment speaker. The challenge followed an extensive evaluation framework that combined \textbf{objective} and \textbf{subjective} metrics. On one hand, the objective metrics were derived from ASV and ASR systems, while on the other hand, the subjective metrics were based on human listening and perception tests.

% \textbf{Objective Metrics}

% Objective \textbf{privacy} evaluation was performed using an ASV system based on x-vector embeddings and probabilistic linear discriminant analysis (PLDA) scoring. Three speaker verification metrics were used: the \textbf{Equal Error Rate (EER)} and the \textbf{log-likelihood-ratio (LLR) cost functions} $C_{\mathrm{llr}}$ and $C_{\mathrm{llr}}^{\min}$.

% \paragraph{Equal Error Rate (EER).}
% Originally a biometric and authentication system metric \citep{noauthor_equal_nodate}, the EER is uded to evaluate the accuracy of recognition and authentication algorithms. In the evaluation benchmark  of the VPC 2020 it was used to measure verifiability after the anonymization. In its computation a threshold ~$\theta$ determines whether the two utterances are classified as belonging to the same speaker. Two error types are possible to occur, a \textit{false alarm}, when two different speakers are incorrectly accepted as the same person (False Acceptance Rate),  and a \textit{miss}, when two utterances from the same speaker are incorrectly rejected (False Rejection Rate).
% \begin{equation}
% P_{\mathrm{fa}}(\theta) = 
% \frac{\#\{\text{impostor trials with score} > \theta\}}
%      {\#\{\text{total impostor trials}\}},
% \qquad
% P_{\mathrm{miss}}(\theta) =
% \frac{\#\{\text{target trials with score} \le \theta\}}
%      {\#\{\text{total target trials}\}}.
% \end{equation}
% The Equal Error Rate corresponds to this point where both error rates are equal:
% \begin{equation}
% \mathrm{EER} = P_{\mathrm{fa}}(\theta_{\mathrm{EER}}) = P_{\mathrm{miss}}(\theta_{\mathrm{EER}}).
% \end{equation}
% A low EER indicates that the ASV system can recognize the speaker accurately, implying weak anonymization. On the otehr hand, a high EER (above 50\%) means that the system performs no better than random guessing, which reflects stronger anonymization, therefore higher privacy.

% \paragraph{Log-likelihood-ratio cost functions.}
% As mentioned in \cite{mingote21_interspeech}, the \(C_{\mathrm{llr}}\) "\textit{...is an application-independent metric that measures the cost of soft detection decisions over all the operating points.}" It is a metric that corresponds to the assessment of ASV and to understand it better, given a set of target and impostor LLR scores, the log-likelihood ratio cost function is defined as:

% \begin{equation}
% C_{\mathrm{llr}} = 
% \frac{1}{2}
% \left[
% \frac{1}{N_{\mathrm{tar}}}
% \sum_{i \in \text{targets}}
% \log_{2}\!\left( 1 + e^{-\mathrm{LLR}_{i}} \right)
% +
% \frac{1}{N_{\mathrm{imp}}}
% \sum_{j \in \text{impostors}}
% \log_{2}\!\left( 1 + e^{\mathrm{LLR}_{j}} \right)
% \right],
% \end{equation}
% where \(N_{\mathrm{tar}}\) and \(N_{\mathrm{imp}}\) are the number of target and impostor trials, respectively. The first term penalizes target trials whose scores are not sufficiently high, while the second penalizes impostor trials whose scores are too high. A perfectly discriminative and calibrated system achieves \(C_{\mathrm{llr}} = 0\), while \(C_{\mathrm{llr}} = 1\) corresponds to a completely random choice by the attacker system and is also what equals to a perfect anonymization system, with the posterior matching the prior of 0.5 \citep{10.1016/j.csl.2021.101299}.

% According to \cite{tomashenko_voiceprivacy_2022}, \(C_{\mathrm{llr}}\) is further broken down to a \textbf{discrimination loss}, \(C_{\mathrm{llr}}^{\min}\), which reflects how well the system  can separate target from impostor trials under optimal/oracle \citep{10.1016/j.csl.2021.101299} calibration (meaning the transformation of raw LLR scores into values that reflect accurately the strength of evidence for or against the same speaker), and a \textbf{calibration loss}, which measures how well the raw scores are already calibrated. In these cost functions high values mean that anonymization has degraded the ASV system's discriminative ability, corresponding to stronger privacy protection.

% \paragraph{Voice Similarity Matrices: De-identification (DeID) \& Gain of Voice Distinctiveness (GVD).}
% These two metrics were introduced by \cite{10.1016/j.csl.2021.101299}, and were used in the first Voice Privacy Challenge \citep{tomashenko_introducing_2020, tomashenko_voiceprivacy_2022-1, tomashenko_voiceprivacy_nodate} as post evaluation in order to provide an interpretable vizualization of anonymization performance across speakers. The goal is to show how well the ASV system would handle the inter-speaker and intra-speaker similarities. First a \emph{voice similarity matrix} $M = (M(i,j))_{1 \leq i,j \leq N}$ is constructed for a set of $N$ speakers using log-likelihood ratio (LLR) scores between segments of their speech:
% \begin{equation}
% M(i, j) = \sigma \!\left(
%   \frac{1}{n_i n_j}
%   \sum_{k=1}^{n_i} \sum_{l=1}^{n_j}
%   \mathrm{LLR}(x_k^{(i)}, x_l^{(j)})
% \right),
% \end{equation}
% where $\sigma(\cdot)$ denotes the sigmoid function, $x_k^{(i)}$ is the $k$-th utterance($x$) of speaker~$i$ (enrollment spear), and $n_i$ is the number of utterances available for that speaker. Three matrices are computed: $M_{oo}$ on original-original data, $M_{aa}$ on anonymized-anonymized data, and $M_{oa}$ between original-anonymized data (original for speaker ~$i$, anonymized for speaker~$j$). 

% Then, the diagonal dominance of a matrix~$M$ measures how well speakers can be distinguished from one another:
% \begin{equation}
% D_{\mathrm{diag}}(M) =
% \left|
%   \frac{1}{N} \sum_{i=1}^{N} M(i,i)
%   - \frac{1}{N(N-1)} \sum_{\substack{j,k=1\\ j \neq k}}^{N} M(j,k)
% \right|.
% \end{equation} 
% It is the absolute difference between the mean similarity of same-speaker pairs (diagonal entries) and that of different-speaker pairs (off-diagonal entries). A higher $D_{\mathrm{diag}}(M)$ indicates clearer separation between speakers, while lower values indicate that the speaker identities are less distinguishable. This measure equals~0 for a uniform matrix (no distinction between speakers, which is what is needed for a successful anonymization system on the $M_{oa}$) and reaches~1 for a perfectly discriminative matrix.

% Using these definitions, the \textbf{de-identification} metric is expressed as
% \begin{equation}
% \mathrm{DeID} = 1 - 
% \frac{D_{\mathrm{diag}}(M_{oa})}{D_{\mathrm{diag}}(M_{oo})},
% \end{equation}
% and is reported in percent. A value of $\mathrm{DeID}=100\%$ corresponds to \emph{perfect de-identification}, where the anonymized voices have no resemblance to the originals, while $\mathrm{DeID}=0\%$ means no de-identification, where anonymized voices can be distinguished. The \textbf{gain of voice distinctiveness} (GVD) evaluates whether the anonymized voices remain distinguishable from one another. It is defined as
% \begin{equation}
% \mathrm{GVD} = 10 \log_{10}
% \!\left(
%   \frac{D_{\mathrm{diag}}(M_{aa})}{D_{\mathrm{diag}}(M_{oo})}
% \right).
% \end{equation}
% A $\mathrm{GVD}$ value of 0 indicates that supposedly, the overall voice distinctiveness has been preserved after anonymization, while positive and negative values correspond respectively to an increase or a loss of voice distinctiveness.

% \paragraph{Word Error Rate - WER.}
% For the \textbf{Utility} of the anonymized audio data, WER was used to assess the performance of the Automatic Speech Recognition System (ASR).

% \textbf{Subjective Metrics}

% In a similar manner as the objective metrics, the subjective evaluation considered once again metrics that assessed both the privacy as well as the utility of the anonymized speech \citep{tomashenko_voiceprivacy_2022-1, tomashenko_voiceprivacy_nodate}.

% For \textbf{privacy}, \textbf{Speaker Verifiability} and \textbf{Speaker Linkability} were utilized. The latter assesses the way humans perceive the distinctive voice characteristics of different speakers, resulting in them being effectively (or non effectively) distinguished. Speaker Verifiability, was assessed between two different utterances (original and anonymized), where the subjects were asked to evaluate the speaker similarity from a scale of 1(different speaker)-10(same speaker).

% On the other hand, the \textbf{utility} of anonyzmized speech was subjectively evaluated with \textbf{Speaker Intelligibility} and \textbf{Speech Naturalness}. The first metric appraises the naturalness of the anonymized utterance through a large scale crowd sourcing evaluation pipeline, in a scale from 1 (unintelligible speech) - 10 (clear and intelligible speech). Lastly, Speech Naturalness is yet again evaluated in a scale from 1 (unnatural speech) - 10 (completely natural speech).

% \subsubsection{Voice Privacy Challenge 2022}
% Many of the metrics that were used in the Voice Privacy Challenge 2020, were used also in VPC 2022, with some additional refinements and newly introduced metrics \citep{tomashenko2022voiceprivacy2022challengeevaluation, tomashenko_voiceprivacy_2022, panariello_voiceprivacy_2024}. The differentiation once again between privacy and utility assertion remains, as well as the objective and subjective evaluation metrics for each. The novelty of this version's challenge was the introduction of a \textbf{Ranking Policy}. In more detail, this policy suggests that the evaluation framework is based on a set of minimum target privacy requirements, instead of ranking the submitted anonymization systems globally using a single privacy-utility tradeoff. The organizers suggest different evaluation conditions with each one of them corresponding to a specific EER threshold for the systems to be satisfied. if these thresholds are satisfied, then the submitted anonymization system would be qualified for ranking. More specifically, there were 4 different EER condition target scores (15\%, 20\%, 25\% and 30\%) that needed to be met by the submitted anonymization systems. For every condition c, only systems greater than the target EERc were retained. Then, those valid submissions systems were ranked based on their average WER. This approach allowed for the comparison of the different systems within comparable privacy bands. Lastly this evaluation approach, allows for the participants to tune their systems to different privacy-utility operating points. From this it was noticed that the higher the EER percentage of an anonymization system was(i.e. the more privacy a system has achieved), the larger WER this system would demonstrate (i.e. utility was lower) \citep{ panariello_voiceprivacy_2024}. 

% \textbf{Objective Metrics}

% Once again, \textbf{EER} was employed as the only metric to measure the level of \textbf{privacy} of each of the anonymization systems. (For additional information on computation of EER see Subsection:\ref{subsec:vpc2020})


% \textbf{WER} was used as the main metric responsible for the assessment of \textbf{utility} (For additional information on computation of WER see Subsection:\ref{subsec:vpc2020}), through two specialized ASR evaluation scenarios: One ASR system decodes the original trial data and the other decodes anonymized data and is trained on an anonymized dataset itself. The second approach seemed more robust and decreased the WER of the anonymized data in comparison to the first ASR system. 

% Moreover, \textbf{GVD} was used in the same manner as the VPC 2020 to evaluate the need for distinctiveness of the anonymized speaker's voices (For additional information on computation of GVD see Subsection:\ref{subsec:vpc2020}).

% \paragraph{Pitch Correlation Metric - $\rho^{F_0}$.}
% This new secondary utility metric was introduced in VPC in 2022 and its goal was to indicate if the anonymized system was able to preserve the intonation (F0 or Fundamental Frequency) of the trial utterance. For its computation, initially the pitch contour is extracted, then because the anonymized and trial utterance durations may differ, the shorter sequence between them gets linearly interpolated so that they both have the same length. Additionally, to find the best alignment between the two pitch trajectories a temporal lag alignment is applied and lastly the Pearson correlation between the two pitch sequences (over only voiced segments) is computed. This gives as a result a $\rho^{F_0}$ per trial-anonymized utterance pair, which are later averaged per dataset.  The rules of the challenge indicated for this metric that all submissions should achieve an average $\rho^{F_0} > 0.3$ for each dataset and condition. 

% \textbf{Subjective Metrics}

% For once again the subjective metrics, similarly to the 2020 VPC include \textbf{speaker verifiability}, \textbf{speech intelligibility}, and \textbf{speech naturalness}. In speaker verifiability pairs of enrollment utterance - original/anonymized trial utterance are assessed by evaluators and rated for similarity on a 1-10 scale. While for speech intelligibility and naturalness evaluators rate single original or anonymized trial utterances on a 1 to 10 scale. In those cases, higher scores indicate more natural sounding speech or higher intelligibility, respectively.

% \subsubsection{Voice Privacy Challenge 2024}
% In VPC 2024, the only objective privacy metric that was kept was the EER, while the subjective metrics were abandoned alltogether. In more detail, the GVD and $\rho^{F_0}$ were discontinued as the requirements for preserving voice distinctiveness and intonation were eliminated in line with the considered application scenarios, which were the preservation of emotional state as a key paralinguistic attribute in many real case scenarios \citep{tomashenko2024voiceprivacy2024challengeevaluation}. Now, all anonymization is performed exclusively on utterance level, meaning that each utterance is assigned a pseudospeaker independently from other utterances. The ranking policy from VPC 2022 was retained, with minimum target EER thresholds (adjusted to 10\%, 20\%, 30\%, and 40\%) defining separate evaluation conditions. Systems meeting each threshold are ranked independently based on WER and UAR (Unweighted Average Recall), assessing the privacy-utility tradeoff for both linguistic and emotional content preservation \citep{tomashenko2024voiceprivacy2024challengeevaluation}.

% The \textbf{EER} remains the sole \textbf{privacy} metric, computed using an ECAPA-TDNN based ASV system following a semi-informed attacker model. The attacker has access to both enrollment utterances and the anonymization system itself, allowing them to anonymize enrollment data and retrain the ASV system on anonymized LibriSpeech-train-clean-360 data to reduce mismatch with anonymized trial data \citep{tomashenko2024voiceprivacy2024challengeevaluation}.

% For \textbf{utility} assessment, two objective metrics were employed: \textbf{WER} and \textbf{UAR (Unweighted Average Recall)}. Unlike VPC 2022, the ASR evaluation model used for WER computation is now fixed and trained exclusively on original LibriSpeech-train-960 data \citep{tomashenko2024voiceprivacy2024challengeevaluation}. The introduction of UAR represents a significant expansion of utility evaluation to include emotional content preservation. The SER (Speech Emotion Recognition) system is trained on IEMOCAP data using a leave-one-conversation-out cross-validation protocol across five folds, evaluating four emotion classes: neutral, sadness, anger, and happiness, The introduction of UAR represents a significant expansion of utility evaluation to include emotional content preservation. The SER (Speech Emotion Recognition) system is trained on IEMOCAP data using a leave-one-conversation-out cross-validation protocol across five folds, evaluating four emotion classes: neutral, sadness, anger, and happiness \citep{tomashenko2024voiceprivacy2024challengeevaluation}. The UAR metric is calculated as the average of class-wise recalls:
% \begin{equation}
% \mathrm{UAR} = \frac{\sum_{i=1}^{N_{\mathrm{class}}} R_i}{N_{\mathrm{class}}},
% \end{equation}
% where $N_{\mathrm{class}}$ is the number of emotion classes (four in this case) and $R_i$ is the recall for class $i$, computed as the number of true positives divided by the total number of samples in that class. It treats all emotion classes equally, preventing bias toward majority classes in imbalanced datasets. The higher the UAR, the better the system preserves emotional content across all emotion categories. \citep{tomashenko2024voiceprivacy2024challengeevaluation}.

\subsection{Evaluation Metrics}
In this section the evaluation metrics used in the assessment of voice 
anonymization systems are presented, organized by challenge edition.

\subsubsection{Voice Privacy Challenge 2020} \label{subsec:vpc2020}
VPC 2020 \citep{tomashenko_introducing_2020, tomashenko_voiceprivacy_2022-1, 
tomashenko_voiceprivacy_nodate} follows an extensive evaluation framework 
combining \textbf{objective} and \textbf{subjective} metrics, with 
objective metrics derived from ASV and ASR systems and subjective metrics 
based on human listening tests.

\paragraph{Equal Error Rate (EER).}
Originally a biometric authentication metric \citep{noauthor_equal_nodate}, 
the EER measures verifiability after anonymization. A threshold $\theta$ 
determines whether two utterances are classified as belonging to the same 
speaker, with two possible error types: a \textit{false alarm} (two 
different speakers incorrectly accepted) and a \textit{miss} (two 
utterances from the same speaker incorrectly rejected):
\begin{equation}
\begin{aligned}
P_{\mathrm{fa}}(\theta)
  &= \frac{\#\{\text{impostor trials with score} > \theta\}}
           {\#\{\text{total impostor trials}\}}, \\[0.75em]
P_{\mathrm{miss}}(\theta)
  &= \frac{\#\{\text{target trials with score} \le \theta\}}
           {\#\{\text{total target trials}\}}.
\end{aligned}
\end{equation}
The EER corresponds to the point where both error rates are equal:
\begin{equation}
\mathrm{EER} = P_{\mathrm{fa}}(\theta_{\mathrm{EER}}) = 
P_{\mathrm{miss}}(\theta_{\mathrm{EER}}).
\end{equation}
A low EER indicates accurate speaker recognition and therefore weak 
anonymization, while a high EER (above 30\%) reflects stronger 
anonymization and higher privacy.

\paragraph{Log-likelihood-ratio cost functions.}
As noted by \cite{mingote21_interspeech}, $C_{\mathrm{llr}}$ is an 
application-independent metric measuring the cost of soft detection 
decisions over all operating points:
\begin{equation}
C_{\mathrm{llr}} = 
\frac{1}{2}
\left[
\frac{1}{N_{\mathrm{tar}}}
\sum_{i \in \text{targets}}
\log_{2}\!\left( 1 + e^{-\mathrm{LLR}_{i}} \right)
+
\frac{1}{N_{\mathrm{imp}}}
\sum_{j \in \text{impostors}}
\log_{2}\!\left( 1 + e^{\mathrm{LLR}_{j}} \right)
\right],
\end{equation}
where $N_{\mathrm{tar}}$ and $N_{\mathrm{imp}}$ are the number of target 
and impostor trials respectively. A perfectly discriminative system 
achieves $C_{\mathrm{llr}} = 0$, while $C_{\mathrm{llr}} = 1$ 
corresponds to random attacker performance and perfect anonymization 
\citep{10.1016/j.csl.2021.101299}. $C_{\mathrm{llr}}$ decomposes into 
a \textbf{discrimination loss} $C_{\mathrm{llr}}^{\min}$, reflecting 
target-impostor separation under optimal calibration, and a 
\textbf{calibration loss}, measuring how well raw scores are already 
calibrated \citep{tomashenko_voiceprivacy_2022-1}. Higher values indicate 
stronger privacy protection.

\paragraph{Voice Similarity Matrices: DeID \& GVD.}
Introduced by \cite{10.1016/j.csl.2021.101299} and used as post-evaluation 
in VPC 2020 \citep{tomashenko_introducing_2020, tomashenko_voiceprivacy_2022-1, 
tomashenko_voiceprivacy_nodate}, these metrics visualize anonymization 
performance across speakers via a voice similarity matrix 
$M = (M(i,j))_{1 \leq i,j \leq N}$ built from LLR scores:
\begin{equation}
M(i, j) = \sigma \!\left(
  \frac{1}{n_i n_j}
  \sum_{k=1}^{n_i} \sum_{l=1}^{n_j}
  \mathrm{LLR}(x_k^{(i)}, x_l^{(j)})
\right),
\end{equation}
with three variants computed: $M_{oo}$ (original-original), $M_{aa}$ 
(anonymized-anonymized), and $M_{oa}$ (original-anonymized). Speaker 
separability is measured via diagonal dominance:
\begin{equation}
D_{\mathrm{diag}}(M) =
\left|
  \frac{1}{N} \sum_{i=1}^{N} M(i,i)
  - \frac{1}{N(N-1)} \sum_{\substack{j,k=1\\ j \neq k}}^{N} M(j,k)
\right|.
\end{equation}
\textbf{De-identification} (DeID) measures how well anonymized voices are 
dissociated from their originals:
\begin{equation}
\mathrm{DeID} = 1 - 
\frac{D_{\mathrm{diag}}(M_{oa})}{D_{\mathrm{diag}}(M_{oo})},
\end{equation}
where DeID$=100\%$ indicates perfect de-identification. \textbf{Gain of 
Voice Distinctiveness} (GVD) evaluates whether anonymized voices remain 
mutually distinguishable:
\begin{equation}
\mathrm{GVD} = 10 \log_{10}
\!\left(
  \frac{D_{\mathrm{diag}}(M_{aa})}{D_{\mathrm{diag}}(M_{oo})}
\right),
\end{equation}
where GVD$=0$ indicates preserved distinctiveness, with positive and 
negative values corresponding to a gain or loss respectively.

\paragraph{Word Error Rate (WER).}
WER assessed utility via ASR performance on anonymized speech and was 
retained as the primary utility metric across all three challenge editions.

Subjective evaluation assessed both privacy and utility 
\citep{tomashenko_voiceprivacy_2022-1, tomashenko_voiceprivacy_nodate}. 
For privacy, \textbf{Speaker Verifiability} rated original-anonymized 
utterance similarity on a 1-10 scale, while \textbf{Speaker Linkability} 
assessed how distinctively human listeners perceived different speakers' 
voices. For utility, \textbf{Speech Intelligibility} and \textbf{Speech 
Naturalness} were each rated on a 1-10 scale via crowdsourced evaluation. 

\subsubsection{Voice Privacy Challenge 2022}
VPC 2022 retained the 2020 evaluation framework with refinements 
\citep{tomashenko2022voiceprivacy2022challengeevaluation, 
tomashenko_voiceprivacy_2022, panariello_voiceprivacy_2024}. The key 
novelty was a \textbf{Ranking Policy}: rather than ranking systems 
globally, four minimum EER thresholds (15\%, 20\%, 25\%, 30\%) defined 
separate evaluation conditions, with systems meeting each threshold ranked 
by average WER. This allowed comparison within comparable privacy bands 
and confirmed that higher privacy consistently came at the cost of utility 
\citep{panariello_voiceprivacy_2024}. EER, WER, and GVD were used 
unchanged from VPC 2020 (see Section~\ref{subsec:vpc2020}), with one 
new metric introduced.

\paragraph{Pitch Correlation Metric - $\rho^{F_0}$.}
This new utility metric assessed whether anonymization preserved 
the intonation contour of the original utterance. Pitch contours are 
extracted for both utterances, the shorter sequence is linearly 
interpolated to match lengths, temporal lag alignment is applied, and 
the Pearson correlation over voiced segments is computed. All submissions 
were required to achieve $\rho^{F_0} > 0.3$.

\subsubsection{Voice Privacy Challenge 2024}
VPC 2024 retained only objective metrics \citep{tomashenko2024voiceprivacy2024challengeevaluation}. 
Here, the ranking policy from VPC 2022 was retained with adjusted EER thresholds (10\%, 20\%, 30\%, 40\%), 
with systems meeting each threshold ranked jointly by WER and UAR.

\textbf{EER} remains the sole \textbf{privacy} metric (see 
Section~\ref{subsec:vpc2020}), computed using an \textbf{Emphasized Channel Attention Propagation and Aggregation Time Delay Neural Network (ECAPA-TDNN)} ASV system 
under the semi-informed attacker model \citep{tomashenko2024voiceprivacy2024challengeevaluation}.

For \textbf{utility}, \textbf{WER} is computed using a fixed ASR model 
trained on original \texttt{LibriSpeech-train-960} (see Section~\ref{subsec:vpc2020}). 
The newly introduced \textbf{UAR (Unweighted Average Recall)} expands 
utility evaluation to emotional content preservation, computed by a SER (Speaker Emotion Recognition)
system:
\begin{equation}
\mathrm{UAR} = \frac{\sum_{i=1}^{N_{\mathrm{class}}} R_i}{N_{\mathrm{class}}},
\end{equation}
where $N_{\mathrm{class}}$ is the number of emotion classes and $R_i$ 
is the per-class recall. Higher values indicate
better emotional content preservation 
\citep{tomashenko2024voiceprivacy2024challengeevaluation}.

 \section{Related Work}
 \label{sec:related_work}
 
 \subsection{The Problem with EER}
 Recent studies have shown that the EER does not accurately capture the variability of privacy protection across speakers. 
 It has been reported that anonymization performance for an individual can still remain very low or very high even though fundamentally 
 different anonymization architectures are employed. In other words, the global EER values often mask substantial inconsistencies in the speaker 
 level and may give a misleading impression of the actual anonymization system's performance.

\cite{williams_anonymizing_2024} argue that anonymization systems exhibit strong speaker-dependent behaviour and identify speakers who are consistently
either easy to imitate or difficult to recognize, regardless of the anonymization method used. Their analysis reveals that many speakers maintain 
similar privacy characteristics across all systems, and accordingly they are being categorized in a rank from easy to difficult to recognise. 
This means that some speakers remain vulnerable across different anonymization methods, showcasing weaknesses in evaluation datasets and in the 
anonymization process itself.

In \cite{gaznepoglu2025sayexploitinglinguisticcontent}, their analysis begins by running the B3, B4, and B5 anonymization systems from the VPC 2024 
and computing individual per speaker EER scores under the semi-informed attack model. This reveals that while mean EER values fall between 20-35 \%, 
several speakers across both genders and across the three baselines score extremely low EERs (even 2\%), indicating that anonymization repeatedly 
fails for them across all systems. Motivated by these findings, the authors wanted to check if the vulnerabilities of some speakers arise from the linguistic patterns they follow. 
They test this by adapting a BERT model as a text-based ASV attacker applied directly to the transcripts of anonymized utterances, with this resulting in EERs as low as 1–2\% for the same vulnerable speakers. 
When trying to investigate the linguistic patterns used by these individuals they found that characteristic vocabulary and thematic patterns, such as religious terms, were able to create intra-speaker linguistic similarity that attackers can exploit.
That means that, some speakers can be re-identified purely from textual information even when their acoustic profiles have been anonymized. 
Following this study, \cite{franzreb_content_2026} demonstrate that the used evaluation corpus in VPC 2024, LibriSpeech contains significant content 
leakage, as the speakers' utterances are constituted by them reading books with distinct lexical and phonetic distributions. 
As a result, speaker verification systems are shown to once again partially rely on differences in phone inventories (meaning specific word usage) 
rather than purely on biometric voice characteristics. From their experiments using representations based solely on phone identities and phone durations, 
they show that a non-trivial re-identification performance can still be achieved even when spectral information is removed or altered through a TTS-ASR anonymization system. 
Also, they report a strong correlation between speaker-specific phonetic distributions and EER (a Pearson correlation of 0.59), indicating that content variability itself acts as an identity cue.

Also, EER's unreliability may not stem only from speaker-dependent variability, but also from confounds built into the evaluation protocol itself.
\cite{franzreb2026improvingspeakeranonymizationevaluations} show that the standard VPC privacy evaluation overestimates privacy when a same-gender 
target selection algorithm assigns pseudospeakers, because the attacker's recognizer learns to encode target-speaker information rather than source-speaker identity alone.
When correcting this with an adversarial, gradient-reversal classifier they noted lower, more realistic EER estimates. 
This suggests that how pseudospeakers are assigned during attacker training can itself distort EER-based privacy claims.

In another related study, \cite{Tsaprazlis_2026} argue that apart from the fact that EER is average-based (which is already mentioned), it can obscure adversaries who 
succeed with high confidence against only a subset of speakers. 
Reframing speaker privacy as a membership inference problem, they show that anonymization systems with near-identical EER can differ by orders of magnitude in true positive rate 
once the adversary is evaluated at very low false-positive rates.
This means that a system can appear equally private under EER while in fact permitting high-confidence re-identification of specific individuals. 
When applying this analysis directly to the VPC 2024 ASRBN baseline, they also find that simple learned attackers recover speaker gender 
and accent with near to perfect accuracy even after anonymization, suggesting that EER evaluation can mask the leakage of sensitive demographic attributes that are not captured by EER.

All in all, these findings reveal a fundamental limitation of global EER: As a single aggregating score it cannot expose the individuals who become outliers, performing either very poorly, nor the sensitive attributes that may remain identifiable even when speaker identity itself is well protected. 

% % % % % % ALTERNATIVE METRICS % % % % % % % % % % % % 
\subsection{Alternative Metrics Focusing on Individual Speakers}
These limitations of the EER have motivated speech researchers to explore alternative evaluation metrics, like \textbf{Linkability \& Singling Out} 
\citep{vauquier25_interspeech} which are grounded in data protection laws, and \textbf{Zero Evidence Biometric Recognition (ZEBRA)} which is inspired by forensics
and aims to understand the variation in privacy preservation provided to a population \citep{nautsch20_interspeech}. 

The ZEBRA framework measures the gap in the attacker system's prior belief and posterior belief after it has observed the anonymized utterances. 
It provides two complementary metrics, the first being the \textbf{Expected Privacy Disclosure} (calculated through Empirical Cross Entropy - ECE) 
which reflects the amount of identity evidence remaining across the whole population \citep{10.1016/j.csl.2021.101299, nautsch20_interspeech}. 
As it is calculated in bits, 0 equals a perfect anonymization. To be more precise, it is computed as the area between the ideal ECE curve and the 
observed curve integrated across all the priors. A perfect anonymization equals a likelihood rate (LR) of 1 (meaning no information gain for the attacker system).

The other metric in the framework is the \textbf{Worst-case Privacy Disclosure} which captures the strongest piece of evidence for an 
individual speaker that was able to survive after the anonymization has been performed. It is the maximum absolute log-likelihood ratio (LLR) observed
when testing the system, which is then mapped into an A-F categorical ranking ("A" showing best privacy). If LLR equals 0 then that means perfect anonymization. 
All in all, the ZEBRA metric seems ideal for evaluation, as it is not based on a single global threshold that ignores outliers, but also provides information about the individual levels of achieved privacy.

In a similar manner, \textbf{Linkability} measures the probability that an anonymized trial embedding is more similar to its true enrollment embedding than to any other 
speaker’s embedding from the pool of enrollment speaker embeddings. High Linkability means therefore that although the speaker identity might not be 
directly revealed, an attacker can aggregate multiple anonymized samples and reconstruct a stable voice profile for an individual speaker. 
This is again considered a failed anonymization because the privacy law says that anonymized data should not allow records of same speakers to be 
connected across different contexts.

\textbf{Singling Out}, estimates the probability that an attacker can design a binary predicate that "fires" on exactly one speaker’s anonymized 
representation. Even if the attacker cannot identify who the speaker is, just isolating a single entry can already be thought of as a failed attempt 
to keep privacy. It is also noteworthy that Linkability \& Singling Out are suggested by the organisers of the VPC of 2024 as privacy metrics that should be integrated in the evaluation pipeline for future versions of the challenge \citep{tomashenko_third_2026}.

\subsection{Prosodic Features as Indicators of Speaker Identification}

\subsubsection{Prosody and Pitch Trajectories}

It is well known that prosody in speech plays a major role in distinguishing speakers from one another in perceptual experiments. 
Prosody is not a single acoustic characteristic, rather a combination of F0 (fundamental frequency or pitch) contours, 
the durations of segments and pauses in speech (i.e. speech tempo) and the specific intonation or linguistic syllable stress.
All of them are features that convey an underlying emotion or specific message \citep{helander_importance_2007, mary_significance_2019}. 
An individual's pitch characteristics have been proven to contain cues about their identity, as they are based on physiological 
properties such as the shape of the vocal tract and the larynx. Therefore, speech scientists in the field of voice anonymization have developed 
systems that target the concealment and modification of the original speaker F0 trajectories. 
In many voice anonymization systems, speaker identity is represented through x-vector embeddings extracted from the source speech, while pitch 
trajectories, were often preserved or modified separately during synthesis \citep{gaznepoglu2022voiceprivacy2022descriptionspeaker}. 
In this context, \cite{Tavi_2022} tested a speaker de-identification approach that explicitly targets dynamic F0 trajectories using Functional 
Data Analysis (FDA) and showcased that pitch variation can be decomposed into principal components that differentiate speaking styles and speaker 
groups. In particular, the first principal component captures global pitch elevation patterns, separating adult from child speech and distinguishing 
male from female speakers, while secondary components encode contour tilt and finer intonational shape differences. 
Their key result was that when dynamic F0 reconstruction is combined with systematic formant shifts (F1–F3), de-identification performance improves 
significantly, yielding relative EER increases of up to 25\% compared to formant-only anonymization and reinforcing the view that speaker identity is also "hidden" in the temporal pitch trajectories.
 
\subsubsection{Phoneme Durations and Phone Diversity}
The importance of analyzing the phoneme durations as a potential privacy leakage channel is demonstrated through a number of papers in speaker verification systems. 
As early as in the work of \cite{bulgakova_speaker_2015}, speaker-specific segmental duration per phone are put into the spotlight as an 
underinvestigated factor of verifying speaker identity. In their method the authors utilize a verification system based on phone durations that relies on 
forced phonetic alignment and computation of average durations for each phone class. Their results show that segmental durational characteristics 
alone provide meaningful discriminative power and that certain phones exhibit particularly strong speaker-discriminative ability. 
More specifically, vowels in stressed and pre-stressed positions, as well as sonorant consonants, are found to carry especially high identity 
information, with sex-specific variations. In more recent years, \cite{tomashenko_analysis_2025} design an attacker that relies exclusively on 
phoneme duration statistics, and show that anonymization systems that preserve phoneme timing patterns remain vulnerable in terms of 
privacy. In more detail, they compute speaker-specific mean duration vectors and apply basic distance metrics (cosine) which indicate that speakers can be re-identified with extremely 
low EERs on original speech and on anonymization systems not using phoneme duration modifications, like B3 on the VPC of 2024. 
On the other hand, when an anonymization system regenerates speech and alters durations, it performs much higher in terms of EER, showcasing a better privacy preservation. 
In a later adaptation \cite{tomashenko_exploiting_2025} move to a more advanced neural attacker that learns context dependent duration embeddings. 
Tested on B3 of VPC 2024 (which modifies utterance pitch and energy), this new approach elevates the ASV performance achieving EERs as low as 1.8\%. 
This approach captures along the individual speaking rate, rhythmic patterns conditioned on the phonetic context, as well as the sequential structure of the phoneme. 
Complementary evidence that phoneme durations encode speaker information is provided by \cite{FUJITA_2024} and even though their work focuses on multi-speaker speech 
synthesis rather than voice privacy, they derive speaker embeddings solely from phoneme identities and their durations
and use them as input to a speaker identification model, to check if it would be able to remap to the original speakers. 
Their experiments achieve a moderate speaker discrimination performance (15.2\% EER), indicating once again that phoneme timing patterns reflect 
partially the original speaker identity. 

Noteworthy in understanding how an attacker re-identifies a speaker is also the contribution of individual phonetic variation of the individual. 
\cite{franzreb_private_2025} in their investigation on how prosodic and phonetic factors contribute to identity leakage in Voice Conversion systems (VC), 
focus on manipulating phone durations and phonetic variation. From their results it is acquired that when phone durations are progressively modified in the output speech, 
EER increases once again confirming that duration patterns encode speaker-specific information. However, it is proved that by reducing phonetic variation (in terms of adopting fewer clusters to quantize the target features of each phone) 
stronger privacy gains are achieved.

A more recent study by \cite{franzreb_content_2026} demonstrates that speaker identity can be partially inferred from speaker-specific distributions of phones and vocabulary, rather than purely from phone duration patterns. This finding highlights the need to 
disentangle temporal dynamics as biometric cues from linguistic content as a confounding factor in privacy evaluation because in real-life scenarios, 
speakers are unlikely to share uniform contextual content across their utterances, and their speech is typically spontaneous and characterized by 
greater linguistic variability. Lastly, also for phonetic diversity, \cite{klein_phonetic_2024} define phonetic richness as the 
count of distinct phonemes present in an utterance and show that it correlates positively with ASV scores, an effect that is especially seen 
for utterances with repetitive content. Their results indicate that not only duration but also the sheer inventory of phonemes realized in an 
utterance carries speaker-discriminative information.

\subsubsection{Rhythm}

Extending the investigation on rhythm-based speaker representations, \cite{mehlman_rhythm_2025} propose a transformer-based speaker identification 
model trained solely on frame-aligned character sequences derived from ASR, no energy features, no pitch and no spectral features. 
Evaluated on LibriSpeech and VoxCeleb1, their rhythm only system achieves an above "chance" accuracy on read speech (0.39 on LibriSpeech), 
but performs markedly worse on spontaneous speech (0.032 on VoxCeleb1). These findings suggest that rhythmic structure does constitute an 
identifiable speaker-specific cue and its discrimination is influenced by speech context and intra-speaker variability. 
From a privacy perspective, this indicates that rhythm-based identity leakage may be dependent on the corpus and therefore suggests a strong re-identification cue for the adversary in the VPC 2024. 
There has also been extensive research on how linguistic rhythm can be captured through simple duration based statistics derived from the 
segmentation of speech in consonants and vowels and how this information can be speaker-specific \citep{ramus_correlates_1999}. 
Specifically, linguists have utilized measures such as the \textbf{Proportion of Vocalic Intervals (\%V)}, the \textbf{Standard Deviation of Vocalic Intervals (ΔV)} and also 
the \textbf{Standard Deviation of Consonantal Intervals (ΔC)}, demonstrating that these features reliably distinguish between rhythm classes and suggesting that rhythm is encoded in the 
distributional properties of segmental timing. Under this theory similar rhythmic properties  may encode speaker-specific "signatures" that persist 
after anonymization, thereby constituting a potential privacy leakage channel. Another metric that looks like it is able to distinguish differences 
in stress- and syllable-timing is again based on vowels and is the \textbf{Pairwise Variability Index (PVI)} \citep{ling_q_2000}. 
This metric captures the normalized durational variability between successive vowels in a sentence and has been shown to reliably distinguish 
rhythmic typologies. Although originally applied at the language level, it could theoretically provide a grounded way of quantifying local temporal 
alternation patterns in speaker anonymization and offer an additional testing cue on whether speakers exhibit stable rhythmic signatures that may 
persist after anonymization. At last, in terms of consonants another proposed rhythm metric is the \textbf{Variation Coefficient of Consonantal Intervals (varcoΔC)}, 
defined as the standard deviation of consonantal interval durations divided by their mean duration and scaled by 100 \citep{dellwo_rhythm_nodate}. 
By normalizing ΔC with respect to the mean consonantal duration, varcoΔC captures relative variability and disentangles rhythmic structure from global tempo effects. 
In the context of voice anonymization, this rate normalized metric could allow for an assessment of whether attackers exploit rhythmic structure or just primarily global temporal characteristics.
 
 \subsection{Correlation of Prosodic and Linguistic Features}
 
Another hypothesis is that the combination of \textit{what} a speaker says and \textit{how} it is said may jointly constitute a cue exploitable by 
an attacker model for re-identification. Supporting this from a syntactic and semantic perspective, \cite{he_layer-wise_2025} investigate the extent 
to which deep speech encoders capture grammatical distinctions directly from acoustic input. Using a minimal pair probing framework, they construct 
controlled speech stimuli differing in specific syntactic properties and apply layer-wise linear probes to representations from multiple pretrained 
speech models. Their experiments show that syntactic distinctions are linearly recoverable from acoustic representations across several architectures 
and layers and that time-resolved probing reveals grammatical information becoming detectable early in the acoustic signal. 
Such evidence supports the hypothesis that temporal dynamics and other acoustic cues may correlate with syntactic and lexical choices and form identity cues when combined with speaker-specific 
phonetic patterns.

\section{Methodology}
\label{sec:methodology}

% This chapter describes the methodological pipeline used to evaluate speaker anonymization performance and to investigate speaker-specific vulnerability after anonymization. The methodology consists of three main parts. The first part consists of the reproduction of the VPC 2024 baseline anonymization systems and the computation of ASV-based privacy scores, including both mean EER and speaker-level EER. The second part introduces complementary privacy metrics from literature, namely the ZEBRA framework, Linkability, and Singling Out. The third part focuses on feature-level analyses designed to examine whether differences in speaker vulnerability can be linked to preserved phonetic/acoustic or linguistic information.

% \subsection{Reproduction of the VPC 2024 Baseline Anonymization Systems \& per speaker EER calculation }

% To be able to apply the additional proposed metrics and then evaluate the phonetic features that might influence the speaker re-identification, the 3 official baseline anonymization systems from VPC 2024 were reproduced: B3, B4, and B5, representing 3 different anonymization paradigms: phonetic transcription-based speech synthesis (B3), neural audio codec language modeling (B4), and ASR bottleneck-based resynthesis with vector quantization (B5). The official VPC 2024 repository was used for the execution of the anonymization pipelines\footnote{\url{https://github.com/Voice-Privacy-Challenge/Voice-Privacy-Challenge-2024}}. The official challenge's recipe and configuration files provided by the organizers in this mentioned repo were used for anonymization. Specifically, the following baseline configurations were used:
% \path{configs/anon_sttts.yaml} for B3, 
% \path{configs/anon_nac.yaml} for B4, and 
% \path{configs/anon_asrbn.yaml} for B5. Moreover, the VPC 2024 LibriSpeech \citep{librispeech} development and test partitions, including the enrollment and trial subsets. For each baseline, the anonymized waveforms were generated while preserving the original directory structure and utterance ids. Also, like the basic recipe, anonymization happened at the utterance level.

% \subsubsection{B3: Phonetic Transcriptions and GAN-based Speaker Anonymization}

% Baseline B3 (or \textit{STTTS} as it will be also used for easier identification from now on) as mentioned before, is a speech synthesis-based anonymization system that preserves linguistic and general prosodic information while replacing the speaker embedding. In the first stage, the system extracts a speaker embedding, phonetic transcription, fundamental frequency ($F0$), energy, and phone durations from the input waveform. The speaker embedding is then replaced by an artificial pseudo-speaker embedding generated by a Wasserstein GAN, while $F0$ and energy are modified at the phone level using random multiplicative factors. The generated pseudo-speaker embedding is accepted only if it is sufficiently dissimilar from the original embedding according to a cosine-distance criterion. Finally, the anonymized waveform is synthesized using a FastSpeech2-based speech synthesis model and a HiFi-GAN vocoder.

% \subsubsection{B4: Neural Audio Codec Language Modeling}

% Baseline B4 (or \textit{NAC} as it will be also used for easier identification from now on) is based on neural audio codec language modeling. The system first extracts semantic tokens from the source utterance, which are intended to represent the linguistic content. In parallel, acoustic tokens are extracted from a pool of pseudo-speaker prompts using a neural audio codec. The semantic tokens of the source utterance are concatenated with a randomly selected acoustic prompt, and a decoder Transformer generates a continuation of acoustic tokens. These generated acoustic tokens are then decoded back into a waveform using the neural audio codec decoder. The goal is to preserve the semantic content of the original utterance while producing speech associated with a different pseudo-speaker.

% \subsubsection{B5: ASR Bottleneck Features with Vector Quantization}

% Baseline B5 (or \textit{ASRBN} as it will be also used for easier identification from now on) is based on ASR bottleneck features with vector quantization. The system extracts fundamental frequency ($F0$) using torch implementation and acoustic vector-quantized bottleneck features from an ASR acoustic model. In B5, this acoustic model combines a pretrained wav2vec2 model with three additional TDNN-F layers. Vector quantization is applied to the inner bottleneck representation of the final TDNN-F layer in order to reduce the amount of speaker information encoded in the bottleneck features. The VQ-BN features, $F0$, and a designated pseudo-speaker represented by an one-hot vector are then used to synthesize the anonymized waveform with a HiFi-GAN network.

This chapter describes the methodological pipeline used to evaluate speaker anonymization performance and investigate speaker-specific vulnerability 
after anonymization. The methodology consists of three main parts: Reproduction of the VPC 2024 baseline anonymization systems and computation of global and speaker-level EER, evaluation using the 
ZEBRA framework, Linkability, Singling Out, and a feature-level analysis examining whether speaker vulnerability can be linked to preserved phonetic, acoustic, or linguistic information.

\subsection{Experimental Setup: Reproduction of the VPC 2024 Baseline Systems}
\label{subsec:reproduction_vpc}

Three official VPC 2024 baseline systems are reproduced: B3, B4, and 
B5, representing fundamentally different anonymization paradigms. The 
official VPC 2024 repository\footnote{\url{https://github.com/Voice-Privacy-Challenge/Voice-Privacy-Challenge-2024}} 
is used as the codebase throughout this thesis, with the official configuration files 
(\path{configs/anon_sttts.yaml}, \path{configs/anon_nac.yaml}, and 
\path{configs/anon_asrbn.yaml} for B3, B4, and B5 respectively) modified as needed when necessary. The semi-informed attacker adopted for this challenge, is trained on the anonymized versions of LibriSpeech \texttt{train-clean-360} \citep{librispeech} produced by each baseline as proposed by the protocol. 
Then, anonymization is applied at utterance level to the VPC 2024 LibriSpeech development and test partitions (four partitions in total: \texttt{dev-f}, \texttt{dev-m}, \texttt{test-f}, \texttt{test-m}).

\subsubsection{B3: Phonetic Transcriptions and GAN-based Speaker 
Anonymization}

Baseline B3 (from now on: \textit{STTTS}) is a speech synthesis based 
system that preserves linguistic and prosodic information while replacing 
the speaker embedding. It extracts: A speaker embedding, phonetic 
transcription, $F0$, energy, and phone durations from the input waveform. 
Then, the speaker embedding gets replaced by a pseudospeaker 
embedding generated from a Wasserstein GAN, accepted only if sufficiently 
dissimilar from the original via a cosine-distance criterion. Pitch 
and energy are randomly perturbed at the phone level. As a last step, the anonymized 
waveform is synthesized using a FastSpeech2 acoustic model and a 
HiFi-GAN vocoder \citep{tomashenko2024voiceprivacy2024challengeevaluation}.

\subsubsection{B4: Neural Audio Codec Language Modeling}

Baseline B4 (from now on: \textit{NAC}) treats anonymization as a sequence generation problem in token space. Semantic tokens 
representing the linguistic content of the source utterance are 
concatenated with acoustic tokens sampled from a pool of pseudospeakers, 
and a decoder transformer generates new acoustic tokens satisfying both 
semantic consistency and constraints of the pseudospeaker's style. The generated 
tokens are then decoded back into a waveform via a neural audio codec 
decoder \citep{tomashenko2024voiceprivacy2024challengeevaluation}.

\subsubsection{B5: ASR Bottleneck Features with Vector Quantization}

Baseline B5 (from now on: \textit{ASRBN}) reduces speaker ID leakage at the 
feature level through vector quantization of ASR bottleneck features. 
The system combines a pretrained wav2vec2 model with three TDNN-F layers, 
applying vector quantization to the inner bottleneck representation of 
the final layer in order to suppress speaker information. The resulting VQ-BN 
features, pitch, and a one-hot pseudospeaker vector are passed to a 
HiFi-GAN network for the last step of waveform synthesis.

\subsubsection{ASV Evaluation with Global EER and Speaker-Level EER}
\label{subsubsec:global_individual_eer}

For each of the three reproduced baselines, the official VPC 2024 ASV evaluation 
pipeline is applied, by performing the needed code and configuration changes, in order to exclude the utility evaluation. Under the
semi-informed attacker scenario, the ECAPA-TDNN attacker
is trained for each anonymization condition using the corresponding 
anonymized version of LibriSpeech \texttt{train-clean-360}, letting the attacker adjust to the same embedding space as the anonymized speech. The official global EER is 
then computed separately for each condition over the LibriSpeech 
enrollment-trial pairs (\texttt{dev-f}, \texttt{dev-m}, \texttt{test-f}, \texttt{test-m}), as the pipeline suggests.

The first evaluation modification comes from the acquisition of speaker-level EER scores by adapting the official ASV evaluation code, similarly to 
\cite{gaznepoglu2025sayexploitinglinguisticcontent}. After trial-level 
scores are computed, the score table is grouped by enrollment speaker 
identifier, and for each speaker $s$ with available target scores 
$\mathcal{T}_s^{+}$ and non-target scores $\mathcal{T}_s^{-}$.
Speakers for whom only one score type is available are excluded. 
The reason for this choice is that in EER calculation, the miss rate computation needs target trials, while the false alarm rate needs the non-target trials.

Therefore, the speaker-level EER gets defined as:
\begin{equation}
\mathrm{EER}_s =
\mathrm{EER}\left(\mathcal{T}_s^{+}, \mathcal{T}_s^{-}\right).
\end{equation} 
The global EER produced by the ASV pipeline is kept unchanged, with the speaker-level analysis 
serving as an additional diagnostic tool for getting an insight of the individual's behaviour rather than a suggested replacement.

\subsection{Alternative Metrics Evaluation}
\label{subsec:alternative_metrics_evaluation}

\subsubsection{ZEBRA Framework}
\label{subsubsec:zebra_framework}

The first set of metrics that are compared to EER, is taken from the ZEBRA framework 
\citep{nautsch20_interspeech}, which evaluates the privacy disclosure from 
an evidence-based perspective using Empirical Cross Entropy (ECE) 
profiles and calibrated log-likelihood ratios (LLRs), unlike EER which identifies a single threshold. The framework includes two complementary measures. 
On the one hand, \texttt{dECE} measures the difference between the ECE of a system with zero evidence 
(meaning perfect privacy) and the ECE profile of the evaluated ASV system's scores over 
the adversary's prior beliefs. A value of 0 in this metric indicates that the
ASV scores (which in the current setup correspond to cosine-based x-vector scoring) provide no identity evidence beyond the adversary's prior, after it has "heard" the anonymized speech, while higher values indicate residual speaker-discriminative information. 
On the other hand, \texttt{max\_abs\_LLR} captures the strongest individual evidence 
observed after ASV score-to-LLR calibration, serving as a worst-case 
disclosure measure. The notion here is that an anonymization system may provide acceptable 
average population protection while still leaving individual speakers highly exposed.

The core ZEBRA computations rely on the existing ZEBRA functions in the 
VPC evaluation framework for privacy assessment. The additional implementation work consists 
of adapting the acquired VPC ASV score outputs to the input format required by the
ZEBRA framework. In more detail, the original VPC score files contain four columns:
\begin{equation}
\left[
\texttt{enroll\_id},\ 
\texttt{trial\_id},\ 
\texttt{label},\ 
\texttt{score}
\right].
\end{equation}
The added preprocessing step generates the two framework-required ZEBRA input files: a 
\path{scores_zebra} file retaining enrollment ID, trial ID, and score:
\begin{equation}
\left[
\texttt{enroll\_id},\ 
\texttt{trial\_id},\ 
\texttt{score}
\right],
\end{equation}
and a \path{key_zebra} file converting the ASV binary labels to textual trial 
types:
\begin{equation}
1 \rightarrow \texttt{target},
\qquad
0 \rightarrow \texttt{nontarget},
\end{equation}
yielding:
\begin{equation}
\left[
\texttt{enroll\_id},\ 
\texttt{trial\_id},\ 
\texttt{trial\_type}
\right].
\end{equation}

Since the ASV scores are raw cosine similarity scores rather than calibrated LLR 
values (which the framework accepts), they have to be transformed using the Pool Adjacent Violators (PAV) calibration, which estimates a monotonic mapping from raw scores to 
posterior probabilities and converts them into LLRs by removing the 
empirical prior log-odds. Laplace smoothing is then applied to avoid the infinite 
LLR values at the score distribution extremes. Letting $\mathcal{S}^{+}$ 
and $\mathcal{S}^{-}$ denote target and non-target scores respectively, PAV calibration produces the calibrated LLRs associated with the target ($\mathcal{L}^{+}$) and non-target ($\mathcal{L}^{-}$) scores:
\begin{equation}
\mathcal{L}^{+}, \mathcal{L}^{-}
=
\mathrm{Calibrate}_{\mathrm{PAV}}
\left(
\mathcal{S}^{+}, \mathcal{S}^{-}
\right).
\end{equation}
The global ZEBRA metrics are then computed as:
\begin{equation}
\mathrm{dECE}
=
\frac{
\mathrm{ECE}(\mathcal{L}^{+})
+
\mathrm{ECE}(-\mathcal{L}^{-})
}{
\log(2)
},
\end{equation}
\begin{equation}
\mathrm{max\_abs\_LLR}
=
\frac{
\max \left(
\left|
\mathcal{L}^{+} \cup \mathcal{L}^{-}
\right|
\right)
}{
\log(10)
},
\end{equation}
where dividing by $\log(10)$ expresses the worst-case evidence in base 10 
LLR units. The \texttt{max\_abs\_LLR} value is mapped to an evidence 
tag from \texttt{A} to \texttt{F}, with lower tags indicating weaker 
evidence and stronger privacy, in the same way the framework defines. At the end, for each evaluated condition, the implementation outputs:
 \begin{itemize}
     \item A calibrated LLR score file,
     \item A global ZEBRA summary,
     \item An ECE plot over the attacker's priors comparing the system curve against the zero-evidence reference (as seen in the Appendix \ref{app:zebra-curves}).
 \end{itemize}
 %add appendix

As an additional diagnostic analysis, the ZEBRA computation is extended 
to the speaker-level to examine whether speakers with low EER also show 
higher evidence leakage across the majority of their anonymized utterances (\texttt{dECE}). For each enrollment speaker $s$, trials are extracted and 
divided into target and non-target subsets, excluding speakers for whom 
only one trial type was available. This filtering step is necessary because \texttt{dECE} is computed from calibrated LLRs estimated from both target and non-target score distributions. PAV calibration is then applied locally per speaker:
\begin{equation}
\mathcal{L}^{+}_{s}, \mathcal{L}^{-}_{s}
=
\mathrm{Calibrate}_{\mathrm{PAV}}
\left(
\mathcal{S}^{+}_{s}, \mathcal{S}^{-}_{s}
\right).
\end{equation}

The speaker-level \texttt{dECE} is computed as:
\begin{equation}
\mathrm{dECE}_{s}
=
\frac{
\mathrm{ECE}(\mathcal{L}^{+}_{s})
+
\mathrm{ECE}(-\mathcal{L}^{-}_{s})
}{
\log(2)
},
\end{equation}
and the speaker-level worst-case evidence as:
\begin{equation}
\mathrm{max\_abs\_LLR}_{s}
=
\frac{
\max \left(
\left|
\mathcal{L}^{+}_{s} \cup \mathcal{L}^{-}_{s}
\right|
\right)
}{
\log(10)
}.
\end{equation}
This pipeline produces a per-speaker ZEBRA table for every trial partition of the form:
\begin{equation}
\left[
\texttt{speaker},\ 
\texttt{n\_trials},\ 
\texttt{n\_tar},\ 
\texttt{n\_non},\ 
\texttt{dECE},\ 
\texttt{max\_abs\_LLR},\ 
\texttt{tag}
\right],
\end{equation}
where \texttt{speaker} denotes the enrollment speaker ID, \texttt{n\_trials} is the total number of trials for that speaker, \texttt{n\_tar} and \texttt{n\_non} denote the number of target and non-target trials respectively, \texttt{dECE} measures the speaker-level deviation from the zero-evidence privacy condition, and \texttt{max\_abs\_LLR} captures the strongest calibrated evidence observed among that speaker's trials. 
The \texttt{tag} field provides the leakage A-F ranking based on \texttt{max\_abs\_LLR}.

%This is important because an anonymization system may appear effective on average while still leaving particular speakers more vulnerable to re-identification. The speaker-level ZEBRA analysis also enables a direct comparison with the individual speaker-level EER scores, allowing the thesis to examine whether speakers with low EER are also the ones with high residual calibrated evidence, or whether the two metrics capture different aspects of speaker-specific privacy risk.

\subsubsection{Linkability \& Singling Out}
\label{subseubsec:linkabilitysinglingout}
The next two metrics are the legally inspired 
Linkability \& Singling Out proposed by 
\cite{vauquier25_interspeech}. First, Linkability measures whether an attacker can correctly associate a trial embedding 
with the corresponding enrollment speaker among the trial candidates, 
while Singling Out measures whether the attacker can \textit{isolate} 
exactly one individual from the set of anonymized speech samples. In a different manner from the EER and ZEBRA initializations, Linkability \& Singling Out are computed using directly the speaker embeddings, rather than the ASV scores.

These metrics are implemented based on the official
repository of the paper \footnote{\url{https://github.com/brijmohan/voice-anonymization-legal-eval}}, but adapted to the VPC 2024 setup using the x-vector embeddings already produced from the anonymization pipelines. Unlike the original paper, which uses Common Voice A and B as disjoint subsets, the current experiments use the LibriSpeech enrollment and trial partitions already met in the VPC 2024 pipeline. The needed enrollment embeddings are loaded from the pooled development and test enrollment directories, namely
(\texttt{libri\_dev\_enrolls} and \texttt{libri\_test\_enrolls}), and the
trial embeddings are loaded from the female and male development and test trial
directories per anonymization method. For this experiment, pooling across the development and test sets is necessary in order to obtain a sufficient number of speakers and utterances for evaluation across multiple population sizes and conversation lengths. This setup adaptation follows the logic of the original Linkability \& Singling Out framework, where enrollment and test material are disjoint at the utterance level but contain overlapping speakers. Thus, the same-speaker identities are used to define the enrollment references and the trial population, but the actual enrollment and trial embeddings are derived from separate utterance partitions, preserving the required separation between enrollment and trial material. Lastly, speaker representations are acquired from
\texttt{speaker\_vectors.pt} files per anonymization method. There, the enrollment data are represented
by a single speaker-level embedding and trial data are grouped by speaker
using \texttt{id2idx} and \texttt{idx2spk} mappings.

Linkage is successful when the cosine similarity between the test embedding
and the correct enrollment embedding exceeds that of all the other candidates:
\begin{equation}
s(x_i^{\mathrm{test}}, x_i^{\mathrm{enroll}})
> \max_{j \neq i} s(x_i^{\mathrm{test}}, x_j^{\mathrm{enroll}}).
\end{equation}
With Linkability defined as the probability of a successful linkage:
\begin{equation}
\pi_{\mathrm{link}} = \Pr_{x_i^{\mathrm{test}}}
\left[
s(x_i^{\mathrm{test}}, x_i^{\mathrm{enroll}})
> \max_{j \neq i} s(x_i^{\mathrm{test}}, x_j^{\mathrm{enroll}})
\right].
\end{equation}
The evaluation is conducted across the following enrollment population sizes:
\begin{equation}
N' \in \{5, 10, 15, 20, 25, 30, 40, 50\},
\end{equation}
and the following conversation lengths:
\begin{equation}
L \in \{1, 3, 10, 30\}.
\end{equation}
The smaller $N'$ range compared to
the original paper (which numbers up to 22,024) reflects the smaller number of common
speakers available in the pooled LibriSpeech trial sets. Moreover, $L=10$
is added as an intermediate setting beyond the original paper's
$\{1, 3, 30\}$ to observe the transition between short and long
conversation conditions more gradually. For each test speaker, $L$
utterance-level embeddings are sampled and averaged
into a single test embedding and speakers who have fewer than $L$
utterances are excluded because the $L$ condition must be constructed from the same number of utterances for every evaluated speaker. For each value of $N'$, a candidate enrollment set
is constructed from the correct speaker and $N'-1$ randomly sampled
impostors, like suggested in the paper. Finally, Linkability is estimated as:
\begin{equation}
\widehat{\pi}_{\mathrm{link}}
= \frac{1}{|\mathcal{T}_L|}
\sum_{i \in \mathcal{T}_L}
\mathbf{1}
\left[
\arg\max_{j \in \mathcal{C}_i}
s(x_i^{\mathrm{test}}, x_j^{\mathrm{enroll}})
= i
\right],
\end{equation}
where $\mathcal{T}_L$ is the set of test speakers with at least $L$
utterances and $\mathcal{C}_i$ the candidate enrollment set for speaker
$i$. The experiment is repeated over 20 random runs per $(L, N')$
pair (more than the original paper's 5 runs) to obtain more stable
estimates given the much smaller data samples.

For Singling Out the attacker defines a binary predicate that fires when similarity to
an enrollment embedding exceeds the calibrated threshold:
\begin{equation}
p(x^{\mathrm{test}}) = \mathbf{1}
\left[
s(x^{\mathrm{test}}, x^{\mathrm{enroll}})
> s^{\mathrm{thresh}}
\right],
\end{equation}
with Singling Out succeeding if exactly one test speaker satisfies
the calibrated predicate:
\begin{equation}
\sum_{i=1}^{N}
p(x_i^{\mathrm{test}})
= 1.
\end{equation}
The random predicate baseline is $\exp(-1) \approx 0.3679$. Evaluation
here is conducted over $N \in \{20, 30, 40, 50\}$ and
$L \in \{1, 3, 10, 30\}$. The smaller $N$ range compared to the
original paper (which numbers up to 22,024) again reflects the limited number of
eligible speakers in the VPC setup. For each speaker, utterance
embeddings are permuted and split into disjoint groups of size $L$,
averaged into grouped embeddings, with up to 10 groups per speaker.
Only speakers with at least three groups are retained (this design choice is stricter than the original paper's minimum of $K \geq 2$) in order to ensure calibration
stability, guaranteeing at least two calibration groups and one test
group per speaker.

To continue, the evaluation follows a cross-validation procedure: In each fold,
one grouped embedding serves as the test embedding and the remainder
as calibration. With $K$ grouped embeddings per speaker, the number
of calibration groups is $M = K-1$. The threshold is set as the
average of the $M$-th and $(M+1)$-th highest pooled calibration score:
\begin{equation}
s^{\mathrm{thresh}} = \frac{1}{2}
\left(
s_{(M)} + s_{(M+1)}
\right),
\end{equation}
generalizing the original paper's fixed $M=9$ rule to variable data
availability. The predicate is then applied to the $N$ test embeddings,
recording success if exactly one exceeds the predicate/threshold:
\begin{equation}
\sum_{i=1}^{N}
\mathbf{1}
\left[
s(x_i^{\mathrm{test}}, x^{\mathrm{enroll}})
> s^{\mathrm{thresh}}
\right]
= 1.
\end{equation}

For each random run, $N$ speaker identities are sampled from the set of speakers that are present in both the pooled enrollment and trial partitions.
The enrollment embedding used to define the predicate is loaded from the enrollment set, while the calibration and test embeddings are constructed from the corresponding anonymized trial utterances.
The full experiment is repeated over five random runs for each $(L,N)$ combination, following the sampling strategy of the original framework.

\subsection{Phonetic Features Analysis}
\label{subsec:phonetic_features_analysis}

Next, a feature-level analysis examines whether speakers who remain more easily re-identifiable (based on ASV similarity and individual EER scores) preserve specific phonetic/acoustic patterns in their anonymized speech. 
The motivation is that speaker identity may be reflected in interpretable speech properties (e.g. phoneme durations, rhythm, fundamental frequency, formant trajectories). 

Initially, phone-level alignments are generated for all trial data (original and anonymized) using the Montreal Forced Aligner (MFA) \citep{mcauliffe17_interspeech}. 
MFA takes as input the speech waveform and its corresponding transcription and produces a time-aligned TextGrid file containing phone-level intervals with start and end times,
from which the phone durations per utterance are extracted. The alignments are applied consistently across all three anonymization systems and both genders on the four trial partitions, preserving original utterance identifiers so that the extracted phone intervals can be matched with the speaker-level ASV results. 

The analysis is performed at utterance level and averaged to speaker level. Let $x_{u}^{\mathrm{orig}}$ denote the original utterance and $x_{u}^{\mathrm{anon}}$ the corresponding anonymized utterance for utterance $u$. For each pair, phonetic, prosodic, acoustic, and embedding-based representations are extracted and compared either 
directly between $x_{u}^{\mathrm{orig}}$ and $x_{u}^{\mathrm{anon}}$, 
or indirectly by comparing anonymized utterances against their original 
speaker centroids. 

\subsubsection{Rhythm Features}
The first feature group consists of rhythm metrics derived from the precalculated MFA alignments. For this and the following analyses, silence and non-speech labels are excluded, and each remaining phone gets assigned to a vowel or consonant class. For each utterance, the following rhythm features get computed separately for both the original and their corresponding anonymized utterances: 
\begin{itemize}
    \item Percentage of vocalic intervals \citep{ramus_correlates_1999},
    \item Percentage of consonantal intervals \citep{ramus_correlates_1999},
    \item Mean vowel duration,
    \item Mean consonant duration,
    \item Standard deviation of vowel and consonant durations ($\Delta V$ and $\Delta C$) \citep{ramus_correlates_1999},
    \item Varco-V and Varco-C, i.e. variation coefficients that normalize $\Delta V$ and $\Delta C$ by the corresponding mean vowel and consonant durations \citep{dellwo_rhythm_nodate},
    \item Raw and normalized pairwise variability index for vowels (rPVI-V and nPVI-V) \citep{ling_q_2000},
    \item Raw and normalized pairwise variability index for consonants (rPVI-C and nPVI-C) \citep{ling_q_2000}.
\end{itemize}

As a next step, two additional comparison measures are derived for each rhythm feature. These values are not independent rhythm metrics, but change indicators that describe how anonymization alters each rhythmic property. 
So, for each utterance $u$ and rhythm feature $f$, the first value, the \textbf{signed shift}, is defined as:
\begin{equation}
s_{u,f} = f(x_u^{anon}) - f(x_u^{orig}),
\end{equation}
while the corresponding \textbf{absolute shift} is defined as:
\begin{equation}
a_{u,f} = \left| f(x_u^{anon}) - f(x_u^{orig}) \right|.
\end{equation}

The signed shift captures the direction of the change caused by anonymization. If positive it means that the rhythm feature is higher than in the original utterance, while a negative value means that the feature is "more suppressed" after anonymization. 
The absolute shift captures the magnitude of the change, meaning how strongly a rhythmic feature is modified by anonymization, regardless of whether the feature's values increase or decrease.

All of these features are then averaged at the speaker level in order to summarize how much, and in which direction, each speaker's rhythmic profile is modified after anonymization.

\subsubsection{Phone-Duration Features}
\label{subsec:phone_duration_features}

% The second feature group focused on phone-duration preservation. Based on the MFA alignments, original and anonymized phone intervals were compared by position. For each utterance, phone-level rows were created containing the utterance identifier, speaker identifier, phone index, original phone label, anonymized phone label, original duration, anonymized duration, and whether the phone label was preserved at the same position.

% For phone $i$ in utterance $u$, the duration shift was defined as

% \begin{equation}
% d_{u,i}^{shift} = d_{u,i}^{anon} - d_{u,i}^{orig},
% \end{equation}

% and the duration ratio was defined as

% \begin{equation}
% r_{u,i} = \frac{d_{u,i}^{anon}}{d_{u,i}^{orig}}.
% \end{equation}

% These measures were included to capture both absolute and relative changes in segmental timing after anonymization. At the utterance level, the analysis computed the number of original and anonymized phones, whether the full phone sequence was preserved, total phone duration, mean phone duration, phone rate, phone rate shift, and Pearson and Spearman correlations between original and anonymized phone durations by position. The correlation measures are used to assess whether the relative duration pattern of the phone sequence is preserved. Speaker-level summaries are then obtained by averaging these utterance-level measures over all utterances belonging to the same speaker.

% In addition, speaker-level phone-duration summaries were computed. For each speaker, the analysis compared the average phone duration and phone rate before and after anonymization. The corresponding shift measures capture whether anonymization makes the speaker's segmental timing slower or faster overall. To assess whether the relative phone duration pattern was preserved, phone durations were also compared by aligned phone position using Pearson and Spearman correlation. These correlation-based measures indicate whether the anonymized utterances preserve the original temporal structure of the phone sequence, independently of global duration changes.

% Finally, a more detailed phone-duration analysis was performed by computing per-speaker and per-phoneme duration statistics, including mean duration, mean shift, mean absolute shift and original-anonymized duration correlation. Local transition features between adjacent phones were also extracted for V$\rightarrow$V, V$\rightarrow$C, C$\rightarrow$V and C$\rightarrow$C transitions, in order to examine whether local timing relations between consecutive phones were preserved after anonymization.
The second feature group focuses on checking if phone durations are preserved in the anonymized speech. 
Based on the MFA alignments, original and anonymized phone intervals are compared by position. 
For each phone $i$ in utterance $u$, the \textbf{duration shift} and \textbf{duration ratio} are defined as:

\begin{equation}
d_{u,i}^{\mathrm{shift}} = d_{u,i}^{\mathrm{anon}} - d_{u,i}^{\mathrm{orig}},
\qquad
r_{u,i} = \frac{d_{u,i}^{\mathrm{anon}}}{d_{u,i}^{\mathrm{orig}}},
\end{equation}
capturing the \textbf{absolute} and \textbf{relative} changes in segmental timing respectively. At the utterance level, the analysis computes individual phone counts, total and mean phone duration, phone rate and phone rate shift, as well as Pearson and Spearman correlations between original and anonymized phone durations by position. 
Additionally, local transition features are extracted for V$\rightarrow$V, V$\rightarrow$C, C$\rightarrow$V, and C$\rightarrow$C in order to examine whether local timing relations between consecutive phones were preserved after anonymization. 
Finally, speaker-level summaries of these relative measures are obtained by averaging the utterance-level measures per speaker.

\subsubsection{Acoustic and Prosodic Features}
\label{subsec:acoustic_and_prosodic_features}

% Thirdly, the acoustic and prosodic preservation between original and anonymized utterances was examined. Here, utterance-level acoustic functionals extracted with openSMILE \citep{10.1145/1873951.1874246} and time varying F0, and formant trajectories extracted with Parselmouth \citep{parselmouth}. The motivation for including these features was that anonymization may not only preserve segmental timing patterns, but also residual prosodic and spectral cues related to pitch and vocal tract characteristics.

% First, acoustic descriptors were extracted with openSMILE using the eGeMAPS feature set. For each original-anonymized utterance pair, the extracted functionals were stored with separate original and anonymized prefixes. The features included F0 related statistics, formant frequencies and bandwidths, jitter, shimmer, harmonics-to-noise ratio, spectral slope features, MFCCs, loudness, and voiced/unvoiced segment descriptors. For each speaker, preservation was summarized by computing Pearson correlation, Spearman correlation, mean signed shift and mean absolute shift between original and anonymized utterance-level feature values.

% Second, dynamic F0 and formant trajectory preservation was analyzed. F0 trajectories were extracted with a 10 ms time step, while F1, F2 and F3 trajectories were estimated using Parselmouth/Praat formant analysis. For each utterance $u$, the original trajectory $z_u^{orig}$ and the anonymized trajectory $z_u^{anon}$ were compared using Pearson correlation, Spearman correlation and dynamic time warping distance (DTW) \citep{stent_dynamic_2024}. Unlike direct correlation DTW allows local temporal realignment between the original and anonymized trajectories before computing their distance and this is useful because anonymization may slightly stretch or compress the timing of an utterance while still preserving the overall pitch or formant contour.
% The correlation-based preservation score was computed for F0, F1, F2, F3, averaged per speaker and defined as:

% \begin{equation}
% \rho_u = corr(z_u^{orig}, z_u^{anon}).
% \end{equation}

% %These resulting speaker-level acoustic and trajectory summaries are later merged with ASV-derived speaker-level measures, including per-speaker EER, target-score statistics, non-target-score statistics and target/non-target score separation, in order to examine whether preservation of pitch, formant structure, spectral characteristics and the prosodic/acoustic cues extracted from openSMILE are associated with speaker reidentification from the attacker after anonymization.

The third feature group examines preservation of acoustic and prosodic characteristics between original and anonymized utterances.

The acoustic descriptors get extracted using openSMILE  \citep{10.1145/1873951.1874246} through the eGeMAPS \citep{eyben_geneva_2016} feature set, 
covering among others, F0, formant frequencies and bandwidths, jitter, shimmer, MFCCs, loudness (energy). For each, speaker preservation is summarized via Pearson and Spearman correlations and mean signed and absolute shifts between the original and the anonymized feature values.

Pitch and formant trajectory preservation is analyzed using Parselmouth \citep{parselmouth}, with F0 getting extracted at a 10 ms time step and
 F1, F2, F3 getting estimated via Praat \citep{praat} formant analysis. For each utterance $u$, the original trajectory $z_u^{\mathrm{orig}}$ and anonymized trajectory $z_u^{\mathrm{anon}}$ are compared using 
Pearson and Spearman correlations and Dynamic Time Warping distance 
(DTW) \citep{stent_dynamic_2024}, which allows local temporal realignment before computing distance and is therefore suitable for cases where anonymization slightly stretches or compresses utterance timing while preserving the overall contour shape. 
The correlation-based preservation score is averaged per speaker and defined as:

\begin{equation}
\rho_u = \mathrm{corr}(z_u^{\mathrm{orig}}, z_u^{\mathrm{anon}}).
\end{equation}

\subsubsection{ASV Embedding Similarity}
\label{subsec:asv_embedding_similarity}

The fourth analysis step uses the ASV embedding similarity of the original and anonymized x-vector embeddings, loaded from the ASV evaluation outputs. For each utterance, the cosine similarity between the original and anonymized embedding is computed as:

\begin{equation}
cos(e_u^{orig}, e_u^{anon}) =
\frac{e_u^{orig} \cdot e_u^{anon}}
{\|e_u^{orig}\| \|e_u^{anon}\|}.
\end{equation}

The speaker-level embedding centroids are also computed by averaging utterance-level embeddings for each speaker. Then the similarity between the original centroid of a speaker and the anonymized centroid of the same speaker gets calculated, as well as the similarity to the anonymized centroids of other speakers. The motivation of this is to estimate whether the anonymized representation of a speaker remains closer to its own original representation than to other speakers.

\subsubsection{Linking Phonetic Features with ASV Scores}
\label{subsec:linking_phonetic_features_with_asv_scores}

For the feature-space similarity analysis (as reported in Subsection \ref{subsec:phonetic_acoustic_linguistic_factors}, 
Table \ref{tab:feature_space_target_non_separation} and Table \ref{tab:feature_space_similarity_asv}), the individual features within each of the extracted feature 
groups get combined into a multi-dimensional feature vector for each utterance. 
These vectors are standardized within each feature space, and original speaker centroids are computed by averaging the original utterance-level 
vectors for each speaker. Cosine similarity is then measured between each anonymized trial-utterance vector and the corresponding original 
enrollment-speaker centroid. Target trials correspond to same-speaker trials, while non-target trials correspond to different-speaker trials.

The target minus non-target feature-space separation is then computed as:
\begin{equation}
\Delta_F = \mu(S_F^{\mathrm{tar}}) - \mu(S_F^{\mathrm{non}}),
\end{equation}
where \(S_F^{\mathrm{tar}}\) and \(S_F^{\mathrm{non}}\) are the target and non-target cosine similarity distributions for feature space \(F\). Larger positive values indicate that same-speaker trials remain more similar than different-speaker trials after anonymization.

% For the speaker-level analysis, the same target-minus-non-target idea is also applied to ASV scores by subtracting the mean non-target ASV score from the mean target ASV score for each speaker. This ASV score separation is later used when correlating feature preservation measures with speaker-level ASV vulnerability.

\subsubsection{Correlation between Phonetic Features and EER Scores}
\label{subsec:feature_eer_scores_correlation}

The speaker-level relationship between a phonetic preservation feature (i.e. rhythm, F0, formant trajectory, and phone-duration) and EER is 
estimated as:

\begin{equation}
r = \mathrm{corr}(F_s, \mathrm{EER}_s).
\end{equation}

A negative correlation indicates that stronger preservation of a feature is associated with lower EER and therefore higher 
re-identification vulnerability, while a positive correlation indicates association with higher EER and stronger privacy. 
The preservation features reflect similarity between original and anonymized speech while the shift and distance features 
(signed shifts, absolute shifts, phone-duration shifts, phone-rate shifts, DTW distance) reflect the amount of change introduced by anonymization. 
For the cause of this thesis, which is to explore if preservation of phonetic/acoustic features helps the attacker re-identify an individual, 
only the results of the negative correlations are reported (as seen in Table \ref{tab:negative_preservation_feature_eer}).

% \subsubsection{Feature-space retrieval analysis}

% To further test whether phonetic information could support speaker re-identification, feature-space retrieval analyses were performed. For a given feature space, original speaker centroids were computed as

% \begin{equation}
% c_s^{orig} =
% \frac{1}{N_s}
% \sum_{u \in s} f(x_u^{orig}),
% \end{equation}

% where $N_s$ is the number of original utterances for speaker $s$. Each anonymized utterance or anonymized speaker centroid was then compared against all original speaker centroids using cosine similarity. The rank of the correct original speaker was recorded, together with top-$1$, top-$3$ and top-$5$ retrieval accuracy. A high top-$k$ retrieval accuracy indicates that the feature space still contains residual speaker-discriminative information after anonymization.

% A related classifier-based analysis was also performed. For selected feature groups, a logistic regression classifier was trained on original utterance-level features to predict speaker identity and then tested on anonymized utterance-level features. Features were standardized before classification, and class-balanced logistic regression was used to reduce the effect of class imbalance. Top-$1$, top-$3$ and top-$5$ classification accuracy were reported. Since the classifier was trained only on original features, above-chance performance on anonymized features would indicate that anonymized utterances still preserve speaker-specific structure in that feature space.


% \subsubsection{Linguistic, lexical and acoustic interaction features}

% In addition to the phonetic and acoustic analyses described above, a linguistic interaction analysis was conducted to examine whether lexical similarity contributes to residual speaker recognizability after anonymization. This analysis was motivated by the possibility that ASV scores may be influenced not only by acoustic or phonetic similarity, but also by overlap in linguistic content between enrollment and trial material.

% Trial and enrollment transcripts were read from the Kaldi-style text files and normalized by lowercasing and removing non-alphanumeric symbols. For each enrollment speaker, all enrollment utterance transcripts were concatenated into a speaker-level text profile. A TF-IDF vectorizer using word unigrams and bigrams was then fitted on the enrollment speaker profiles and the trial utterance transcripts. For each ASV trial pair, lexical similarity was computed as the cosine similarity between the TF-IDF vector of the trial utterance and the TF-IDF vector of the corresponding enrollment speaker profile:

% \begin{equation}
% sim_{text}(e,u) =
% cos(v_u^{text}, v_e^{text}),
% \end{equation}

% where $v_u^{text}$ is the TF-IDF representation of trial utterance $u$ and $v_e^{text}$ is the TF-IDF representation of enrollment speaker $e$.

% Additional lexical statistics were computed for each trial utterance, including word count, character count, type-token ratio, function-word ratio and average word length. These features were used to characterize the linguistic content of the trial utterance and to test whether lexical properties were associated with ASV score variation.

% The lexical features were then merged with the previously computed phonetic and acoustic feature-space similarities and with the ASV score table. The phonetic/acoustic similarity features included rhythm similarity, openSMILE acoustic similarity, trajectory-preservation similarity, phone-duration-preservation similarity and ASV embedding similarity. To examine whether lexical similarity modulated the effect of phonetic or acoustic similarity, interaction features were created by multiplying text similarity with each phonetic/acoustic similarity feature:

% \begin{equation}
% I_{u,k} =
% sim_{text}(e,u) \cdot sim_{k}(e,u),
% \end{equation}

% where $sim_{k}(e,u)$ denotes the similarity value from phonetic or acoustic feature space $k$. These interaction terms were used to test whether ASV scores were better explained when lexical similarity and phonetic/acoustic similarity were considered jointly.

% The merged trial-level table was analyzed using correlations and predictive models. First, Pearson and Spearman correlations were computed between ASV scores and lexical, phonetic/acoustic and interaction features. These correlations provide an initial view of whether individual features are linearly or monotonically associated with the ASV score. Ridge regression was then used to compare how well four feature groups predicted ASV scores: text-only features, phonetic/acoustic-only features, combined text and phonetic/acoustic features, and combined features with interaction terms. The Combined model includes both text-based and phonetic/acoustic features in the same regression model, allowing for testing whether phonetic/acoustic information adds explanatory power beyond the text-based variables alone. The interaction model further adds interaction terms between the feature groups, in order to examine whether the relationship between phonetic/acoustic information and ASV scores depends on the textual similarity of the trial. Ridge regression was selected because the feature sets contain correlated variables and interaction terms, which can make ordinary linear regression unstable. The ridge penalty regularizes the regression coefficients and therefore provides a more robust estimate of how much ASV-score variance can be explained by each feature group. Model performance was evaluated using cross-validated $R^2$. The $R^2$ score is defined as:

% \begin{equation}
% R^2 = 1 - \frac{\sum_i (y_i - \hat{y}_i)^2}{\sum_i (y_i - \bar{y})^2},
% \end{equation}

% where $y_i$ is the true ASV score, $\hat{y}_i$ is the predicted ASV score, and $\bar{y}$ is the mean true ASV score. A higher $R^2$ therefore means that the corresponding feature group captures more information related to the ASV scoring behavior, while values close to zero indicate limited explanatory power.

% Additionally, logistic regression was used to test whether the same feature groups could distinguish target from non-target trials. This analysis treats the trial label as a binary outcome, where target trials correspond to same-speaker comparisons and non-target trials correspond to different-speaker comparisons. Logistic regression estimates the probability that a trial belongs to the target class based on the selected feature group.

% The resulting predicted probabilities were evaluated using the receiver operating characteristic area under the ROC-AUC curve, which was used because it measures how well the model ranks target trials above non-target trials across all possible decision thresholds. It corresponds to the area under the ROC curve, where the true positive rate is plotted against the false positive rate. It can be interpreted as the probability that a randomly selected target trial receives a higher predicted target probability than a randomly selected non-target trial. A value of $0.5$ corresponds to chance-level discrimination, while values closer to $1.0$ indicate stronger separation.

% Model performance is evaluated using cross-validated $R^2$ for continuous ASV-score prediction and ROC-AUC for target/non-target classification. These two metrics allow the analysis to address two complementary questions:
% \begin{itemize}
%     \item How much of the variation in ASV scores can be explained by each feature group, namely text-only features, phonetic/acoustic-only features, combined features, and interaction-based features?
%     \item Do these feature groups preserve enough speaker-discriminative information to separate target trials from non-target trials, as reflected by the ROC-AUC score?
% \end{itemize}

\subsubsection{Linguistic and Acoustic Interaction Features}
\label{subsec:linguistic_acoustic_interaction}

At last, a linguistic interaction analysis is conducted to examine whether lexical similarity between enrollment and trial material 
contributes to residual speaker recognizability after anonymization, independently and in combination with phonetic and acoustic features.

Trial and enrollment transcripts are initially normalized, then for each enrollment speaker, all enrollment transcripts are concatenated into a speaker-level text profile, and a TF-IDF vectorizer using word unigrams and bigrams is fitted on both the enrollment profiles and trial transcripts. 
Lexical similarity for each ASV trial pair is then computed as:
\begin{equation}
\mathrm{sim}_{\mathrm{text}}(e,u) =
\cos(v_u^{\mathrm{text}}, v_e^{\mathrm{text}}),
\end{equation}
where $v_u^{\mathrm{text}}$ and $v_e^{\mathrm{text}}$ are the TF-IDF representations of trial utterance $u$ and enrollment speaker $e$ respectively. Additional lexical statistics are computed per trial utterance, namely: Word count, character count, type-token ratio, function-word ratio, and average word length.

Then, the lexical features are merged with the phonetic and acoustic feature space similarities computed earlier (rhythm, openSMILE, F0 preservation, trajectory preservation, phone-duration preservation, and ASV embedding similarity) and with the ASV score table. The interaction features are created by multiplying text similarity with each phonetic/acoustic similarity feature in the following manner:

\begin{equation}
I_{u,k} =
\mathrm{sim}_{\mathrm{text}}(e,u) \cdot \mathrm{sim}_{k}(e,u),
\end{equation}

where $\mathrm{sim}_{k}(e,u)$ denotes the similarity value from phonetic or acoustic feature space $k$. \textbf{Interaction} tests whether ASV scores are better explained when lexical and phonetic/acoustic similarity are considered jointly.

The merged table is analyzed using Pearson and Spearman correlations between ASV scores and all feature groups, followed by Ridge regression to compare the predictive power of four feature configurations: Text-only, phonetic/acoustic-only, combined, and combined with interaction terms. 
Ridge Regression is selected due to the correlated variables and interaction terms in the feature sets, with the ridge penalty providing more robust coefficient estimates. Model performance is evaluated using five-fold shuffled cross-validation at the ASV trial-pair level, with a fixed random seed. The folds are not speaker-grouped, so the reported cross-validated $R^2$ may be optimistic for unseen-speaker generalization:
\begin{equation}
R^2 = 1 - 
\frac{\sum_i (y_i - \hat{y}_i)^2}
{\sum_i (y_i - \bar{y})^2},
\end{equation}
where $y_i$ is the true ASV score, $\hat{y}_i$ the predicted score, and $\bar{y}$ the mean true score. 
Additionally, Logistic Regression is applied to test whether the same feature groups could distinguish target from non-target trials, evaluated using cross-validated ROC-AUC, 
interpreted as the probability that a randomly selected target trial receives a higher predicted score than a randomly selected non-target trial, with $0.5$ corresponding to chance and values closer to $1.0$ indicating stronger separation. 
Together, $R^2$ and ROC-AUC address two complementary questions: 

\begin{itemize}
    \item How much of the variation in ASV scores can be explained by each feature group (namely: Text-only features, phonetic/acoustic-only features, combined features, and interaction-based features)? (as demonstrated in Table \ref{tab:linguistic_model_comparison_all})
     \item Do these feature groups preserve enough speaker-discriminative information to separate target trials from non-target trials, as reflected by the ROC-AUC score? (as demonstrated in Table \ref{tab:linguistic_model_comparison_all})
\end{itemize}

\section{Results}
\label{sec:results}

In the current chapter, the results of the experiments and analyses described in Section~\ref{sec:methodology} are presented. First, the results from the trial partitions of LibriSpeech in VPC 2024, namely \texttt{dev-f}, \texttt{dev-m}, \texttt{test-f}, and \texttt{test-m}, are reported for global EER, global \texttt{dECE} according to the ZEBRA framework, speaker-level EER, and speaker-level \texttt{dECE}. Next, the results for Linkability \& Singling Out are presented and lastly, the results of the phonetic, acoustic, embedding-based, and linguistic feature analyses are reported. The analysis is presented through feature-space target/non-target separation, correlations with ASV similarity scores, predictive modeling of ASV-score variation, and speaker-level correlations with individual EER scores.

% Table~\ref{tab:dev_f_zebra_eer_ranks} reports the per-speaker ZEBRA dECE and EER results for the development female trial set. Overall, the global results indicate that B5 achieves the best privacy performance with the lowest global dECE value ($0.073$) and the highest global EER ($34.231\%$). B4 follows closely, with a global dECE of $0.088$ and an EER of $33.949\%$, while B3 performs worse overall, with higher dECE ($0.136$) and lower EER ($28.265\%$).

% At the speaker level, the results show substantial variability. Some speakers are consistently well protected, while others remain vulnerable. For example, speaker 84 is strongly protected for B4 and B5, with high EER values ($61.105\%$ and $69.142\%$) and very low dECE values, including $0.000$ for B4. In contrast, speaker 1673 remains vulnerable across systems, with low EER and high dECE values. 

% The table also shows that EER and dECE do not always rank speakers in the same way. For B4, speaker 6313 has a low EER of $19.949\%$, suggesting weak protection, but a relatively low dECE of $0.087$. Additionally, speaker 3536 has an EER of $30.278\%$, suggesting better protection, but a higher dECE of $0.201$. Speaker 2412 shows a similar mismatch, with an EER close to $30\%$ but a high dECE of $0.243$. For B3, speaker 6319 is also noteworthy, since their EER is close to the privacy-preserving range ($29.391\%$), while its dECE remains relatively moderate ($0.176$).

% \begin{table}[H]
% \centering
% \textbf{dev\_f}
% \scriptsize
% \setlength{\tabcolsep}{4pt}
% \begin{tabular}{lcccccccccccc}
% \toprule
% \textbf{Spk.}
% & \textbf{B3 dECE} & \textbf{DR} & \textbf{B3 EER} & \textbf{ER}
% & \textbf{B4 dECE} & \textbf{DR} & \textbf{B4 EER} & \textbf{ER}
% & \textbf{B5 dECE} & \textbf{DR} & \textbf{B5 EER} & \textbf{ER} \\
% \midrule
% 84   & 0.253 & 8  & 22.873 & 8  & 0.000 & 1  & 61.105 & 1  & 0.005 & 1  & 69.142 & 1  \\
% 1462 & 0.343 & 11 & 18.492 & 11 & 0.143 & 9  & 31.718 & 6  & 0.194 & 12 & 22.708 & 13 \\
% \textcolor{red}{1673} & 0.452 & 15 & \textcolor{red}{10.078} & 15 & 0.388 & 15 & \textcolor{red}{18.651} & 15 & 0.412 & 14 & \textcolor{red}{14.515} & 14 \\
% 1988 & 0.020 & 3  & 46.715 & 3  & 0.008 & 2  & 55.476 & 2  & 0.032 & 2  & 48.853 & 2  \\
% 1993 & 0.378 & 12 & 14.824 & 13 & 0.072 & 4  & 34.073 & 5  & 0.047 & 4  & 37.035 & 7  \\
% \textcolor{red}{2035} & 0.202 & 6  & \textcolor{red}{29.979} & 5  & 0.141 & 8  & \textcolor{red}{29.979} & 9  & 0.050 & 6  & \textcolor{red}{26.498} & 12 \\
% 2277 & 0.210 & 7  & 27.349 & 7  & 0.084 & 5  & 38.332 & 4  & 0.149 & 10 & 30.148 & 9  \\
% \textcolor{red}{2412} & 0.401 & 14 & \textcolor{red}{17.083} & 12 & 0.243 & 13 & \textcolor{red}{29.271} & 10 & 0.491 & 15 & \textcolor{red}{12.188} & 15 \\
% 3536 & 0.286 & 9  & 21.215 & 10 & 0.201 & 11 & 30.278 & 8  & 0.094 & 8  & 39.342 & 4  \\
% 5338 & 0.121 & 4  & 36.354 & 4  & 0.124 & 7  & 31.412 & 7  & 0.034 & 3  & 39.082 & 6  \\
% 5895 & 0.311 & 10 & 21.988 & 9  & 0.168 & 10 & 29.134 & 11 & 0.048 & 5  & 39.463 & 3  \\
% \textcolor{red}{6313} & 0.386 & 13 & \textcolor{red}{14.721} & 14 & 0.087 & 6  & \textcolor{red}{19.949} & 13 & 0.166 & 11 & \textcolor{red}{29.391} & 10 \\
% \textcolor{red}{6319} & 0.176 & 5  & \textcolor{red}{29.391} & 6  & 0.318 & 14 & \textcolor{red}{19.949} & 13 & 0.196 & 13 & \textcolor{red}{29.391} & 10 \\
% 6345 & 0.016 & 2  & 48.535 & 2  & 0.052 & 3  & 43.258 & 3  & 0.075 & 7  & 39.289 & 5  \\
% 8842 & 0.004 & 1  & 54.738 & 1  & 0.217 & 12 & 23.790 & 12 & 0.135 & 9  & 30.845 & 8  \\
% \midrule
% \textbf{Global}
% & \textbf{0.136} & -- & \textbf{28.265} & --
% & \textbf{0.088} & -- & \textbf{33.949} & --
% & \textbf{0.073} & -- & \textbf{34.231} & -- \\
% \bottomrule
% \end{tabular}
% \caption{Per-speaker ZEBRA dECE, EER, and corresponding ranks for B3, B4, and B5 on the development female trial set. DR corresponds to the dECE speaker rank, while ER corresponds to the EER speaker rank. Red marks speakers whose EER is below $30\%$ across all systems.}
% \label{tab:dev_f_zebra_eer_ranks}
% \end{table}


% For global dECE, Table~\ref{tab:dev_m_zebra_eer_ranks} shows that B4 and B5 perform similarly, with dECE values of $0.118$ for both systems and EER values of $30.588\%$ and $30.740\%$, respectively. B3 performs worse overall, with a higher dECE of $0.196$ and a lower EER of $23.788\%$. 

% Several speakers remain vulnerable across all three systems. Speakers 652, 2803, 2902, 5536, 7976, and 8297 have EER values below $30\%$ for B3, B4, and B5, meaning that anonymization does not provide sufficient protection under the EER criterion. Among them, speakers 652, 2803, and 5536 are particularly vulnerable, as they combine low EER values with high dECE values in at least two systems. For example, speaker 2803 has EER values of $9.174\%$, $18.946\%$, and $17.102\%$, together with high dECE values across all systems. However, EER and dECE do not always agree perfectly at the speaker level. For B4, speaker 6241 has a dECE value of $0.145$, indicating some leakage, while its EER is $31.659\%$, suggesting adequate protection. 

% \begin{table}[H]
% \centering
% \textbf{dev\_m}
% \scriptsize
% \setlength{\tabcolsep}{4pt}

% \begin{tabular}{lcccccccccccc}
% \toprule
% \textbf{Spk.}
% & \textbf{B3 dECE} & \textbf{DR} & \textbf{B3 EER} & \textbf{ER}
% & \textbf{B4 dECE} & \textbf{DR} & \textbf{B4 EER} & \textbf{ER}
% & \textbf{B5 dECE} & \textbf{DR} & \textbf{B5 EER} & \textbf{ER} \\
% \midrule
% 174  & 0.122 & 1  & 36.129 & 1  & 0.163 & 8  & 31.241 & 7  & 0.150 & 7  & 30.538 & 7  \\
% 251  & 0.124 & 2  & 33.333 & 3  & 0.118 & 6  & 31.528 & 6  & 0.105 & 5  & 37.274 & 4  \\
% \textcolor{red}{652}  & 0.581 & 14 & \textcolor{red}{8.011}  & 13 & 0.174 & 9  & \textcolor{red}{25.802} & 9  & 0.406 & 14 & \textcolor{red}{10.000} & 14 \\
% 1272 & 0.167 & 7  & 25.847 & 10 & 0.105 & 5  & 29.622 & 8  & 0.051 & 3  & 39.617 & 2  \\
% 2428 & 0.214 & 9  & 27.283 & 6  & 0.048 & 2  & 40.924 & 3  & 0.051 & 2  & 41.197 & 1  \\
% \textcolor{red}{2803} & 0.501 & 12 & \textcolor{red}{9.174} & 12 & 0.386 & 14 & \textcolor{red}{18.946} & 13 & 0.399 & 13 & \textcolor{red}{17.102} & 13 \\
% \textcolor{red}{2902} & 0.576 & 13 & \textcolor{red}{7.130}  & 14 & 0.190 & 10 & \textcolor{red}{24.596} & 10 & 0.236 & 10 & \textcolor{red}{29.393} & 8  \\
% 3752 & 0.127 & 3  & 29.719 & 5  & 0.061 & 3  & 35.722 & 4  & 0.065 & 4  & 38.668 & 3  \\
% \textcolor{red}{5536} & 0.359 & 11 & \textcolor{red}{15.401} & 11 & 0.385 & 13 & \textcolor{red}{12.816} & 14 & 0.177 & 9  & \textcolor{red}{28.543} & 10 \\
% 5694 & 0.135 & 4  & 30.182 & 4  & 0.063 & 4  & 42.385 & 2  & 0.110 & 6  & 35.444 & 6  \\
% 6241 & 0.137 & 5  & 33.868 & 2  & 0.145 & 7  & 31.659 & 5  & 0.311 & 12 & 22.209 & 12 \\
% 6295 & 0.145 & 6  & 27.283 & 6  & 0.036 & 1  & 45.435 & 1  & 0.036 & 1  & 36.521 & 5  \\
% \textcolor{red}{7976} & 0.204 & 8  & \textcolor{red}{26.645} & 8  & 0.203 & 11 & \textcolor{red}{23.300} & 11 & 0.175 & 8  & \textcolor{red}{29.331} & 9  \\
% \textcolor{red}{8297} & 0.216 & 10 & \textcolor{red}{26.019} & 9  & 0.300 & 12 & \textcolor{red}{19.651} & 12 & 0.254 & 11 & \textcolor{red}{25.526} & 11 \\
% \midrule
% \textbf{Global}
% & \textbf{0.196} & -- & \textbf{23.788} & --
% & \textbf{0.118} & -- & \textbf{30.588} & --
% & \textbf{0.118} & -- & \textbf{30.740} & -- \\
% \bottomrule
% \end{tabular}
% \vspace{-4pt}
% \caption{Per-speaker ZEBRA dECE, EER, and corresponding ranks for B3, B4, and B5 on the development male trial set. DR corresponds to the dECE speaker rank, while ER corresponds to the EER speaker rank. Red marks speakers whose EER is below $30\%$ across all systems.}
% \label{tab:dev_m_zebra_eer_ranks}
% \end{table}

% Table~\ref{tab:test_f_zebra_eer_ranks} shows that B5 performs best overall in global dECE, with the lowest global dECE value of $0.107$ and the highest global EER of $32.483\%$. B3 follows with a dECE of $0.125$ and an EER of $28.682\%$, while B4 has the highest dECE value ($0.142$) and its EER is slightly above the privacy threshold at $30.109\%$.

% In more detail, peakers 1995, 3570, and 1221 have EER values below $30\%$ for B3, B4, and B5, meaning that they remain insufficiently protected under the EER criterion. Speaker 3570 is particularly vulnerable, with very low EER values for B3 and B4 ($6.886\%$ and $6.533\%$) and the highest dECE values for both systems ($0.557$ and $0.621$). Speaker 1221 also remains highly vulnerable across all systems, with low EER values and high dECE values.

% The trend also remains and here also EER and dECE do not always agree at the speaker level. For B4, speaker 2961 has an EER of $30.407\%$, which suggests that it is close to being protected, but its dECE is $0.274$, ranking it among the speakers with the highest information leakage. A similar mismatch appears for B5, where speaker 2961 again has an EER of $30.407\%$, but its dECE remains relatively high at $0.229$.

% \begin{table}[H]
% \centering
% \textbf{test\_f}
% \scriptsize
% \setlength{\tabcolsep}{4pt}

% \begin{tabular}{lcccccccccccc}
% \toprule
% \textbf{Spk.}
% & \textbf{B3 dECE} & \textbf{DR} & \textbf{B3 EER} & \textbf{ER}
% & \textbf{B4 dECE} & \textbf{DR} & \textbf{B4 EER} & \textbf{ER}
% & \textbf{B5 dECE} & \textbf{DR} & \textbf{B5 EER} & \textbf{ER} \\
% \midrule
% 121  & 0.013 & 3  & 54.539 & 2  & 0.280 & 13 & 22.873 & 13 & 0.110 & 7  & 45.103 & 1  \\
% 237  & 0.331 & 11 & 18.700 & 10 & 0.144 & 7  & 33.357 & 4  & 0.062 & 3  & 40.469 & 4  \\
% 1284 & 0.292 & 8  & 23.527 & 8  & 0.078 & 2  & 40.028 & 2  & 0.060 & 2  & 40.028 & 5  \\
% 1580 & 0.239 & 7  & 25.622 & 7  & 0.138 & 6  & 33.357 & 4  & 0.164 & 10 & 30.161 & 10 \\
% \textcolor{red}{1995} & 0.394 & 12 & \textcolor{red}{14.368} & 13 & 0.206 & 9  & \textcolor{red}{26.173} & 11 & 0.323 & 15 & \textcolor{red}{21.914} & 15 \\
% 2961 & 0.295 & 9  & 17.767 & 11 & 0.274 & 12 & 30.407 & 8  & 0.229 & 12 & 30.407 & 9  \\
% \textcolor{red}{3570} & 0.557 & 16 & \textcolor{red}{6.886}  & 16 & 0.621 & 16 & \textcolor{red}{6.533}  & 16 & 0.264 & 14 & \textcolor{red}{25.905} & 13 \\
% 4446 & 0.000 & 1  & 84.766 & 1  & 0.237 & 11 & 25.066 & 12 & 0.258 & 13 & 23.420 & 14 \\
% 4970 & 0.402 & 13 & 15.505 & 12 & 0.095 & 3  & 34.353 & 3  & 0.059 & 1  & 43.884 & 3  \\
% 4992 & 0.304 & 10 & 21.374 & 9  & 0.290 & 14 & 21.087 & 14 & 0.063 & 4  & 36.309 & 6  \\
% 5142 & 0.004 & 2  & 48.685 & 4  & 0.132 & 5  & 33.357 & 4  & 0.148 & 9  & 29.596 & 12 \\
% 5683 & 0.147 & 5  & 31.896 & 6  & 0.129 & 4  & 31.896 & 7  & 0.140 & 8  & 29.814 & 11 \\
% 6829 & 0.454 & 14 & 12.518 & 14 & 0.173 & 8  & 28.592 & 9  & 0.174 & 11 & 31.152 & 8  \\
% 8463 & 0.019 & 4  & 51.078 & 3  & 0.209 & 10 & 27.658 & 10 & 0.086 & 6  & 36.135 & 7  \\
% 8555 & 0.171 & 6  & 34.050 & 5  & 0.052 & 1  & 41.399 & 1  & 0.064 & 5  & 44.825 & 2  \\
% \textcolor{red}{1221} & 0.519 & 15 & \textcolor{red}{12.288} & 15 & 0.370 & 15 & \textcolor{red}{14.046} & 15 & 0.458 & 16 & \textcolor{red}{9.341}  & 16 \\
% \midrule
% \textbf{Global}
% & \textbf{0.125} & -- & \textbf{28.682} & --
% & \textbf{0.142} & -- & \textbf{30.109} & --
% & \textbf{0.107} & -- & \textbf{32.483} & -- \\
% \bottomrule
% \end{tabular}

% \vspace{-4pt}
% \caption{Per-speaker ZEBRA dECE, EER, and corresponding ranks for B3, B4, and B5 on the test female trial set. DR corresponds to the dECE speaker rank, while ER corresponds to the EER speaker rank. Red marks speakers whose EER is below $30\%$ across all systems.}
% \label{tab:test_f_zebra_eer_ranks}
% \end{table}

% \vspace{-4pt}
% Lastly, for global dECE, Table~\ref{tab:test_m_zebra_eer_ranks} shows that once again B5 performs best overall, with the lowest dECE value of $0.116$ and the highest global EER of $32.770\%$. B4 follows with a dECE of $0.135$ and an EER of $29.619\%$, while B3 performs worse, with a higher dECE of $0.210$ and a lower EER of $26.278\%$.

% At the speaker level once again there are several speakers that remain vulnerable. Speakers 1320, 4077, 7021, and 8224 have EER values below $30\%$ for B3, B4, and B5, indicating insufficient protection under the EER criterion and high dECE values, especially for B3 and B4.

% Again inconsistencies are noticed between speaker specific EER and dECE. In more detail, in B5 speaker 61 has an EER of $30.244\%$, which suggests that the speaker is close to protected according to the EER threshold, but their dECE value is $0.156$, indicating measurable residual information leakage. This and the other aforementioned cases, suggest that even when EER suggests acceptable privacy protection, ZEBRA may still reveal remaining speaker information.

% \begin{table}[H]
% \centering
% \textbf{test\_m}
% \scriptsize
% \setlength{\tabcolsep}{4pt}

% \begin{tabular}{lcccccccccccc}
% \toprule
% \textbf{Spk.}
% & \textbf{B3 dECE} & \textbf{DR} & \textbf{B3 EER} & \textbf{ER}
% & \textbf{B4 dECE} & \textbf{DR} & \textbf{B4 EER} & \textbf{ER}
% & \textbf{B5 dECE} & \textbf{DR} & \textbf{B5 EER} & \textbf{ER} \\
% \midrule
% 61   & 0.335 & 7  & 21.294 & 7  & 0.086 & 4  & 37.815 & 5  & 0.156 & 7  & 30.244 & 7  \\
% 260  & 0.118 & 2  & 38.608 & 2  & 0.057 & 2  & 43.388 & 1  & 0.108 & 3  & 31.787 & 6  \\
% 908  & 0.392 & 10 & 17.216 & 10 & 0.115 & 6  & 37.860 & 4  & 0.365 & 13 & 17.761 & 13 \\
% 1089 & 0.243 & 4  & 25.034 & 4  & 0.169 & 7  & 35.020 & 7  & 0.243 & 9  & 19.973 & 11 \\
% \textcolor{red}{1320} & 0.245 & 5  & \textcolor{red}{24.470} & 5  & 0.382 & 11 & \textcolor{red}{14.668} & 11 & 0.202 & 8  & \textcolor{red}{26.244} & 9  \\
% 2830 & 0.032 & 1  & 44.397 & 1  & 0.201 & 9  & 23.972 & 9  & 0.117 & 6  & 35.177 & 3  \\
% \textcolor{red}{4077} & 0.465 & 12 & \textcolor{red}{11.788} & 12 & 0.512 & 12 & \textcolor{red}{11.788} & 12 & 0.338 & 11 & \textcolor{red}{23.510} & 10 \\
% 5105 & 0.409 & 11 & 12.842 & 11 & 0.108 & 5  & 35.929 & 6  & 0.037 & 2  & 45.691 & 2  \\
% 6930 & 0.547 & 13 & 8.149  & 13 & 0.077 & 3  & 38.462 & 3  & 0.109 & 4  & 34.668 & 5  \\
% \textcolor{red}{7021} & 0.388 & 9  & \textcolor{red}{18.632} & 9  & 0.375 & 10 & \textcolor{red}{18.356} & 10 & 0.347 & 12 & \textcolor{red}{18.908} & 12 \\
% 7127 & 0.268 & 6  & 23.665 & 6  & 0.186 & 8  & 33.885 & 8  & 0.111 & 5  & 34.977 & 4  \\
% 7176 & 0.196 & 3  & 25.614 & 3  & 0.035 & 1  & 41.052 & 2  & 0.002 & 1  & 51.367 & 1  \\
% \textcolor{red}{8224} & 0.382 & 8  & \textcolor{red}{20.607} & 8  & 0.564 & 13 & \textcolor{red}{7.114} & 13 & 0.272 & 10 & \textcolor{red}{27.053} & 8  \\
% \midrule
% \textbf{Global}
% & \textbf{0.210} & -- & \textbf{26.278} & --
% & \textbf{0.135} & -- & \textbf{29.619} & --
% & \textbf{0.116} & -- & \textbf{32.770} & -- \\
% \bottomrule
% \end{tabular}

% \vspace{-4pt}
% \caption{Per-speaker ZEBRA dECE, EER, and corresponding ranks for B3, B4, and B5 on the test male trial set. DR corresponds to the dECE speaker rank, while ER corresponds to the EER speaker rank. Red marks speakers whose EER is below $30\%$ across all systems.}
% \label{tab:test_m_zebra_eer_ranks}

% \end{table}
% Finally, \texttt{max\_abs\_LLR} scores are seen in Table~\ref{tab:max_abs_llr_scores}. These values indicate the strongest evidence produced by the calibrated LLR scores. Lower \texttt{max\_abs\_LLR} values suggest that the system produces less extreme evidence for speaker identity, while higher values indicate stronger residual speaker evidence. This is important because EER provides a global summary of verification performance, whereas ZEBRA allows the analysis to move closer to the level of individual information leakage.

% Overall, the \texttt{max\_abs\_LLR} values range from $1.415$ to $2.389$. The lowest value is observed for B4 on the test-m set ($1.415$, category B), while the highest value is observed for B5 on the test-f set ($2.389$, category C). Across the four trial sets, B4 most often falls into category B, except on dev-m where it reaches category C. B3 is mostly assigned to category C, while B5 alternates between categories B and C depending on the trial set.

% \begin{table}[H]
% \centering
% \small
% \setlength{\tabcolsep}{6pt}
% \begin{tabular}{llcc}
% \hline
% Trial set & System & \texttt{max\_abs\_LLR} & Category \\
% \hline
%   & B3 & \color{red}2.219 & C \\
% dev-f  & B4 & 1.617 & B \\
%   & B5 & 1.793 & B \\
% \hline
%   & B3 & 2.039 & C \\
% dev-m  & B4 & \color{red}2.340 & C \\
%   & B5 & 1.696 & B \\
% \hline
%  & B3 & 1.979 & B \\
% test-f & B4 & 1.787 & B \\
%  & B5 & \color{red}2.389 & C \\
% \hline
%  & B3 & \color{red}2.169 & C \\
% test-m & B4 & 1.415 & B \\
% & B5 & 2.022 & C \\
% \hline
% \end{tabular}
% \caption{Maximum absolute log-likelihood-ratio values across trial sets and anonymization systems \& qualitative leakage category assigned to each maximum absolute LLR value.}
% \label{tab:max_abs_llr_scores}
% \end{table}

\subsection{Results for EER \& ZEBRA}
\label{subsec:results_for_eer_and_zebra}

Table~\ref{tab:dev_f_zebra_eer_ranks} reports per-speaker ZEBRA, \texttt{dECE} 
and EER results for the LibriSpeech development female trial set. As seen, the ASRBN pipeline achieves the best global privacy performance with the lowest \texttt{dECE} ($0.073$) and the highest EER ($34.231\%$). The NAC system follows with a \texttt{dECE} of $0.088$ and an EER of $33.949\%$, while the STTTS baseline performs the worst with a dECE of $0.136$ and an EER of $28.265\%$. 

\begin{table}[H]
\centering
\textbf{dev\_f}\\[2pt]
\scriptsize
\setlength{\tabcolsep}{3pt}
\begin{tabular}{lcccccccccccc}
\toprule
\textbf{Spk.}
& \textbf{B3 dECE} & \textbf{DR} & \textbf{B3 EER} & \textbf{ER}
& \textbf{B4 dECE} & \textbf{DR} & \textbf{B4 EER} & \textbf{ER}
& \textbf{B5 dECE} & \textbf{DR} & \textbf{B5 EER} & \textbf{ER} \\
\midrule
84   & 0.253 & 8  & 22.873 & 8  & 0.000 & 1  & 61.105 & 1  
     & 0.005 & 1  & 69.142 & 1  \\
1462 & 0.343 & 11 & 18.492 & 11 & 0.143 & 9  & 31.718 & 6  
     & 0.194 & 12 & 22.708 & 13 \\
\textcolor{red}{1673} 
     & 0.452 & 15 & \textcolor{red}{10.078} & 15 
     & 0.388 & 15 & \textcolor{red}{18.651} & 15 
     & 0.412 & 14 & \textcolor{red}{14.515} & 14 \\
1988 & 0.020 & 3  & 46.715 & 3  & 0.008 & 2  & 55.476 & 2  
     & 0.032 & 2  & 48.853 & 2  \\
1993 & 0.378 & 12 & 14.824 & 13 & 0.072 & 4  & 34.073 & 5  
     & 0.047 & 4  & 37.035 & 7  \\
\textcolor{red}{2035} 
     & 0.202 & 6  & \textcolor{red}{29.979} & 5  
     & 0.141 & 8  & \textcolor{red}{29.979} & 9  
     & 0.050 & 6  & \textcolor{red}{26.498} & 12 \\
2277 & 0.210 & 7  & 27.349 & 7  & 0.084 & 5  & 38.332 & 4  
     & 0.149 & 10 & 30.148 & 9  \\
\textcolor{red}{2412} 
     & 0.401 & 14 & \textcolor{red}{17.083} & 12 
     & 0.243 & 13 & \textcolor{red}{29.271} & 10 
     & 0.491 & 15 & \textcolor{red}{12.188} & 15 \\
3536 & 0.286 & 9  & 21.215 & 10 & 0.201 & 11 & 30.278 & 8  
     & 0.094 & 8  & 39.342 & 4  \\
5338 & 0.121 & 4  & 36.354 & 4  & 0.124 & 7  & 31.412 & 7  
     & 0.034 & 3  & 39.082 & 6  \\
5895 & 0.311 & 10 & 21.988 & 9  & 0.168 & 10 & 29.134 & 11 
     & 0.048 & 5  & 39.463 & 3  \\
\textcolor{red}{6313} 
     & 0.386 & 13 & \textcolor{red}{14.721} & 14 
     & 0.087 & 6  & \textcolor{red}{19.949} & 13 
     & 0.166 & 11 & \textcolor{red}{29.391} & 10 \\
\textcolor{red}{6319} 
     & 0.176 & 5  & \textcolor{red}{29.391} & 6  
     & 0.318 & 14 & \textcolor{red}{19.949} & 13 
     & 0.196 & 13 & \textcolor{red}{29.391} & 10 \\
6345 & 0.016 & 2  & 48.535 & 2  & 0.052 & 3  & 43.258 & 3  
     & 0.075 & 7  & 39.289 & 5  \\
8842 & 0.004 & 1  & 54.738 & 1  & 0.217 & 12 & 23.790 & 12 
     & 0.135 & 9  & 30.845 & 8  \\
\midrule
\textbf{Global}
& \textbf{0.136} & -- & \textbf{28.265} & --
& \textbf{0.088} & -- & \textbf{33.949} & --
& \textbf{0.073} & -- & \textbf{34.231} & -- \\
\bottomrule
\end{tabular}
\caption{Per-speaker ZEBRA dECE, EER, and performance ranks for individual speakers for B3, B4, B5 on \texttt{dev-f}, with ranking of 1 meaning best privacy preservation. DR=dECE speaker rank, ER=EER speaker rank. Red: Speakers with EER below $30\%$ across all systems.}
\label{tab:dev_f_zebra_eer_ranks}
\end{table}

At the speaker level though, substantial variability between speakers and systems can be observed. More specifically, speaker \textbf{84} seems strongly protected under NAC and ASRBN (EERs of $61.105\%$ and $69.142\%$ and \texttt{dECE} of $0$ in NAC), while speaker \textbf{1673} remains very vulnerable across all systems with constant EERs under $20\%$ and also ranking last under \texttt{dECE} across baselines. Another interesting finding is that EER and \texttt{dECE} do not always agree in speaker rankings: Specifically under NAC, speaker \textbf{6313} has a low EER of $19.949\%$ (suggesting big privacy risk), but a low \texttt{dECE} of $0.087$ suggesting otherwise. Speaker \textbf{3536} has an EER of $30.278\%$ suggesting protection but a higher \texttt{dECE} of $0.201$. Lastly, speaker \textbf{2412} shows a similar mismatch with an EER close to $30\%$ but a moderate \texttt{dECE} of $0.243$.


Table~\ref{tab:dev_m_zebra_eer_ranks} shows that NAC and ASRBN perform similarly on \texttt{dev-m}, with a \texttt{dECE} of $0.118$ for both and an EER of $30.588\%$ and $30.740\%$ respectively. STTTS performs worse with a \texttt{dECE} of $0.196$ and an EER of $23.788\%$. Speakers \textbf{652}, \textbf{2803}, \textbf{2902}, \textbf{5536}, \textbf{7976}, and \textbf{8297} remain vulnerable across all three systems making this trial set the one with the most vulnerable individuals. Speaker \textbf{2803} is 
particularly exposed, with EER values of $9.174\%$, $18.946\%$, and $17.102\%$ and quite high \texttt{dECE} across all systems. Again, an EER and \texttt{dECE} diverge is noted for some speakers: For example, under NAC, speaker \textbf{6241} has a \texttt{dECE} of $0.145$ indicating leakage while their EER of $31.659\%$ suggests adequate protection.

\begin{table}[H]
\centering
\textbf{dev\_m}\\[2pt]
\scriptsize
\setlength{\tabcolsep}{3pt}
\begin{tabular}{lcccccccccccc}
\toprule
\textbf{Spk.}
& \textbf{B3 dECE} & \textbf{DR} & \textbf{B3 EER} & \textbf{ER}
& \textbf{B4 dECE} & \textbf{DR} & \textbf{B4 EER} & \textbf{ER}
& \textbf{B5 dECE} & \textbf{DR} & \textbf{B5 EER} & \textbf{ER} \\
\midrule
174  & 0.122 & 1  & 36.129 & 1  & 0.163 & 8  & 31.241 & 7  
     & 0.150 & 7  & 30.538 & 7  \\
251  & 0.124 & 2  & 33.333 & 3  & 0.118 & 6  & 31.528 & 6  
     & 0.105 & 5  & 37.274 & 4  \\
\textcolor{red}{652}  
     & 0.581 & 14 & \textcolor{red}{8.011}  & 13 
     & 0.174 & 9  & \textcolor{red}{25.802} & 9  
     & 0.406 & 14 & \textcolor{red}{10.000} & 14 \\
1272 & 0.167 & 7  & 25.847 & 10 & 0.105 & 5  & 29.622 & 8  
     & 0.051 & 3  & 39.617 & 2  \\
2428 & 0.214 & 9  & 27.283 & 6  & 0.048 & 2  & 40.924 & 3  
     & 0.051 & 2  & 41.197 & 1  \\
\textcolor{red}{2803} 
     & 0.501 & 12 & \textcolor{red}{9.174}  & 12 
     & 0.386 & 14 & \textcolor{red}{18.946} & 13 
     & 0.399 & 13 & \textcolor{red}{17.102} & 13 \\
\textcolor{red}{2902} 
     & 0.576 & 13 & \textcolor{red}{7.130}  & 14 
     & 0.190 & 10 & \textcolor{red}{24.596} & 10 
     & 0.236 & 10 & \textcolor{red}{29.393} & 8  \\
3752 & 0.127 & 3  & 29.719 & 5  & 0.061 & 3  & 35.722 & 4  
     & 0.065 & 4  & 38.668 & 3  \\
\textcolor{red}{5536} 
     & 0.359 & 11 & \textcolor{red}{15.401} & 11 
     & 0.385 & 13 & \textcolor{red}{12.816} & 14 
     & 0.177 & 9  & \textcolor{red}{28.543} & 10 \\
5694 & 0.135 & 4  & 30.182 & 4  & 0.063 & 4  & 42.385 & 2  
     & 0.110 & 6  & 35.444 & 6  \\
6241 & 0.137 & 5  & 33.868 & 2  & 0.145 & 7  & 31.659 & 5  
     & 0.311 & 12 & 22.209 & 12 \\
6295 & 0.145 & 6  & 27.283 & 6  & 0.036 & 1  & 45.435 & 1  
     & 0.036 & 1  & 36.521 & 5  \\
\textcolor{red}{7976} 
     & 0.204 & 8  & \textcolor{red}{26.645} & 8  
     & 0.203 & 11 & \textcolor{red}{23.300} & 11 
     & 0.175 & 8  & \textcolor{red}{29.331} & 9  \\
\textcolor{red}{8297} 
     & 0.216 & 10 & \textcolor{red}{26.019} & 9  
     & 0.300 & 12 & \textcolor{red}{19.651} & 12 
     & 0.254 & 11 & \textcolor{red}{25.526} & 11 \\
\midrule
\textbf{Global}
& \textbf{0.196} & -- & \textbf{23.788} & --
& \textbf{0.118} & -- & \textbf{30.588} & --
& \textbf{0.118} & -- & \textbf{30.740} & -- \\
\bottomrule
\end{tabular}
\caption{Per-speaker ZEBRA dECE, EER, and performance ranks for individual speakers for B3, B4, B5 on \texttt{dev-m}, with ranking of 1 meaning best privacy preservation. DR=dECE speaker rank, ER=EER speaker rank. Red: Speakers with EER below $30\%$ across all systems.}
\label{tab:dev_m_zebra_eer_ranks}
\end{table}

Table~\ref{tab:test_f_zebra_eer_ranks} shows ASRBN performing best on  \texttt{test-f} with \texttt{dECE} of $0.107$ and EER of $32.483\%$, followed by STTTS (\texttt{dECE} of $0.125$, EER of $28.682\%$) and NAC (\texttt{dECE} $0.142$, EER $30.109\%$). Speakers \textbf{1995}, \textbf{3570}, and \textbf{1221} remain insufficiently protected across all systems. Speaker \textbf{3570} is particularly vulnerable, with EER values of $6.886\%$ and $6.533\%$ under STTTS and ASRBN and the highest \texttt{dECE} values for both ($0.557$ and $0.621$). Speaker \textbf{2961} illustrates EER-\texttt{dECE} divergence: Their EER of $30.407\%$ under NAC suggests borderline protection, while its \texttt{dECE} of $0.274$ ranks them among the most 
exposed speakers.

\begin{table}[H]
\centering
\textbf{test\_f}\\[2pt]
\scriptsize
\setlength{\tabcolsep}{3pt}
\begin{tabular}{lcccccccccccc}
\toprule
\textbf{Spk.}
& \textbf{B3 dECE} & \textbf{DR} & \textbf{B3 EER} & \textbf{ER}
& \textbf{B4 dECE} & \textbf{DR} & \textbf{B4 EER} & \textbf{ER}
& \textbf{B5 dECE} & \textbf{DR} & \textbf{B5 EER} & \textbf{ER} \\
\midrule
121  & 0.013 & 3  & 54.539 & 2  & 0.280 & 13 & 22.873 & 13 
     & 0.110 & 7  & 45.103 & 1  \\
237  & 0.331 & 11 & 18.700 & 10 & 0.144 & 7  & 33.357 & 4  
     & 0.062 & 3  & 40.469 & 4  \\
1284 & 0.292 & 8  & 23.527 & 8  & 0.078 & 2  & 40.028 & 2  
     & 0.060 & 2  & 40.028 & 5  \\
1580 & 0.239 & 7  & 25.622 & 7  & 0.138 & 6  & 33.357 & 4  
     & 0.164 & 10 & 30.161 & 10 \\
\textcolor{red}{1995} 
     & 0.394 & 12 & \textcolor{red}{14.368} & 13 
     & 0.206 & 9  & \textcolor{red}{26.173} & 11 
     & 0.323 & 15 & \textcolor{red}{21.914} & 15 \\
2961 & 0.295 & 9  & 17.767 & 11 & 0.274 & 12 & 30.407 & 8  
     & 0.229 & 12 & 30.407 & 9  \\
\textcolor{red}{3570} 
     & 0.557 & 16 & \textcolor{red}{6.886}  & 16 
     & 0.621 & 16 & \textcolor{red}{6.533}  & 16 
     & 0.264 & 14 & \textcolor{red}{25.905} & 13 \\
4446 & 0.000 & 1  & 84.766 & 1  & 0.237 & 11 & 25.066 & 12 
     & 0.258 & 13 & 23.420 & 14 \\
4970 & 0.402 & 13 & 15.505 & 12 & 0.095 & 3  & 34.353 & 3  
     & 0.059 & 1  & 43.884 & 3  \\
4992 & 0.304 & 10 & 21.374 & 9  & 0.290 & 14 & 21.087 & 14 
     & 0.063 & 4  & 36.309 & 6  \\
5142 & 0.004 & 2  & 48.685 & 4  & 0.132 & 5  & 33.357 & 4  
     & 0.148 & 9  & 29.596 & 12 \\
5683 & 0.147 & 5  & 31.896 & 6  & 0.129 & 4  & 31.896 & 7  
     & 0.140 & 8  & 29.814 & 11 \\
6829 & 0.454 & 14 & 12.518 & 14 & 0.173 & 8  & 28.592 & 9  
     & 0.174 & 11 & 31.152 & 8  \\
8463 & 0.019 & 4  & 51.078 & 3  & 0.209 & 10 & 27.658 & 10 
     & 0.086 & 6  & 36.135 & 7  \\
8555 & 0.171 & 6  & 34.050 & 5  & 0.052 & 1  & 41.399 & 1  
     & 0.064 & 5  & 44.825 & 2  \\
\textcolor{red}{1221} 
     & 0.519 & 15 & \textcolor{red}{12.288} & 15 
     & 0.370 & 15 & \textcolor{red}{14.046} & 15 
     & 0.458 & 16 & \textcolor{red}{9.341}  & 16 \\
\midrule
\textbf{Global}
& \textbf{0.125} & -- & \textbf{28.682} & --
& \textbf{0.142} & -- & \textbf{30.109} & --
& \textbf{0.107} & -- & \textbf{32.483} & -- \\
\bottomrule
\end{tabular}
\caption{Per-speaker ZEBRA dECE, EER, and performance ranks for individual speakers for B3, B4, B5 on \texttt{test-f}, with ranking of 1 meaning best privacy preservation. DR=dECE speaker rank, ER=EER speaker rank. Red: Speakers with EER below $30\%$ across all systems.}
\label{tab:test_f_zebra_eer_ranks}
\end{table}

On Table \ref{tab:test_m_zebra_eer_ranks} \texttt{test-m}, ASRBN again performs best with a \textbf{dECE} of $0.116$ and an EER of $32.770\%$, followed by NAC (\texttt{dECE} $0.135$, EER $29.619\%$) and STTTS (\texttt{dECE} $0.210$, EER $26.278\%$). Here the individuals \textbf{1320}, \textbf{4077}, \textbf{7021}, and \textbf{8224} remain insufficiently protected across all systems, with particularly high \texttt{dECE} values under STTTS and NAC. Speaker \textbf{61} illustrates a further EER-\texttt{dECE} mismatch: Their EER of $30.244\%$ under ASRBN suggests borderline protection, while their \texttt{dECE} of $0.156$ indicates moderate leakage, confirming that EER alone can mask remaining speaker information.

\begin{table}[H]
\centering
\textbf{test\_m}\\[2pt]
\scriptsize
\setlength{\tabcolsep}{3pt}
\begin{tabular}{lcccccccccccc}
\toprule
\textbf{Spk.}
& \textbf{B3 dECE} & \textbf{DR} & \textbf{B3 EER} & \textbf{ER}
& \textbf{B4 dECE} & \textbf{DR} & \textbf{B4 EER} & \textbf{ER}
& \textbf{B5 dECE} & \textbf{DR} & \textbf{B5 EER} & \textbf{ER} \\
\midrule
61   & 0.335 & 7  & 21.294 & 7  & 0.086 & 4  & 37.815 & 5  
     & 0.156 & 7  & 30.244 & 7  \\
260  & 0.118 & 2  & 38.608 & 2  & 0.057 & 2  & 43.388 & 1  
     & 0.108 & 3  & 31.787 & 6  \\
908  & 0.392 & 10 & 17.216 & 10 & 0.115 & 6  & 37.860 & 4  
     & 0.365 & 13 & 17.761 & 13 \\
1089 & 0.243 & 4  & 25.034 & 4  & 0.169 & 7  & 35.020 & 7  
     & 0.243 & 9  & 19.973 & 11 \\
\textcolor{red}{1320} 
     & 0.245 & 5  & \textcolor{red}{24.470} & 5  
     & 0.382 & 11 & \textcolor{red}{14.668} & 11 
     & 0.202 & 8  & \textcolor{red}{26.244} & 9  \\
2830 & 0.032 & 1  & 44.397 & 1  & 0.201 & 9  & 23.972 & 9  
     & 0.117 & 6  & 35.177 & 3  \\
\textcolor{red}{4077} 
     & 0.465 & 12 & \textcolor{red}{11.788} & 12 
     & 0.512 & 12 & \textcolor{red}{11.788} & 12 
     & 0.338 & 11 & \textcolor{red}{23.510} & 10 \\
5105 & 0.409 & 11 & 12.842 & 11 & 0.108 & 5  & 35.929 & 6  
     & 0.037 & 2  & 45.691 & 2  \\
6930 & 0.547 & 13 & 8.149  & 13 & 0.077 & 3  & 38.462 & 3  
     & 0.109 & 4  & 34.668 & 5  \\
\textcolor{red}{7021} 
     & 0.388 & 9  & \textcolor{red}{18.632} & 9  
     & 0.375 & 10 & \textcolor{red}{18.356} & 10 
     & 0.347 & 12 & \textcolor{red}{18.908} & 12 \\
7127 & 0.268 & 6  & 23.665 & 6  & 0.186 & 8  & 33.885 & 8  
     & 0.111 & 5  & 34.977 & 4  \\
7176 & 0.196 & 3  & 25.614 & 3  & 0.035 & 1  & 41.052 & 2  
     & 0.002 & 1  & 51.367 & 1  \\
\textcolor{red}{8224} 
     & 0.382 & 8  & \textcolor{red}{20.607} & 8  
     & 0.564 & 13 & \textcolor{red}{7.114}  & 13 
     & 0.272 & 10 & \textcolor{red}{27.053} & 8  \\
\midrule
\textbf{Global}
& \textbf{0.210} & -- & \textbf{26.278} & --
& \textbf{0.135} & -- & \textbf{29.619} & --
& \textbf{0.116} & -- & \textbf{32.770} & -- \\
\bottomrule
\end{tabular}
\caption{Per-speaker ZEBRA dECE, EER, and performance ranks for individual speakers for B3, B4, B5 on \texttt{test-m}, with ranking of 1 meaning best privacy preservation. DR=dECE speaker rank, ER=EER speaker rank. Red: Speakers with EER below $30\%$ across all systems.}
\label{tab:test_m_zebra_eer_ranks}
\end{table}

Lastly, Table~\ref{tab:max_abs_llr_scores} reports \texttt{max\_abs\_LLR} values across all conditions. Values range from $1.415$ to $2.389$. The lowest one is observed for NAC on \texttt{test-m} ($1.415$, category B) and the highest for ASRBN on \texttt{test-f} ($2.389$, category C). NAC most frequently falls into category B, STTTS predominantly into category C, while ASRBN alternates between B and C across trial sets. The results indicate that strong 
global performance does not preclude worst-case individual evidence leakage.

The corresponding ZEBRA ECE curves for all three systems and trial partitions are visualized in Appendix~\ref{app:zebra-curves}.
\begin{table}[H]
\centering
\small
\setlength{\tabcolsep}{6pt}
\begin{tabular}{llcc}
\hline
Trial set & System & \texttt{max\_abs\_LLR} & Category \\
\hline
          & B3 & \color{red}2.219 & C \\
\texttt{dev-f}     & B4 & 1.617 & B \\
          & B5 & 1.793 & B \\
\hline
          & B3 & 2.039 & C \\
\texttt{dev-m}     & B4 & \color{red}2.340 & C \\
          & B5 & 1.696 & B \\
\hline
          & B3 & 1.979 & B \\
\texttt{test-f}    & B4 & 1.787 & B \\
          & B5 & \color{red}2.389 & C \\
\hline
          & B3 & \color{red}2.169 & C \\
\texttt{test-m}    & B4 & 1.415 & B \\
          & B5 & 2.022 & C \\
\hline
\end{tabular}
\caption{Maximum absolute LLR values and qualitative leakage category 
across trial sets and anonymization systems.}
\label{tab:max_abs_llr_scores}
\end{table}

\subsection{Results for Linkability \& Singling Out}
\label{subsec:results_linkability_and_singling_out}

The Linkability results in Table~\ref{tab:linkability_all_results} and Figure~\ref{fig:linkability_plots} follow the same trend reported by \cite{vauquier25_interspeech}. That is, Linkability decreases as $N'$ (number of speakers in
\begin{figure}[H]
    \centering
    \includegraphics[width=\textwidth]{linkability_plots.pdf}
    \caption{Linkability as a function of the number of enrollment 
    speakers $N'$ for different conversation lengths $L$, for all three systems STTTS, NAC and ASRBN.}
    \label{fig:linkability_plots}
\end{figure}
 the enrollment pool from whom the attacker has to choose from) increases, while the metric also scores higher with longer conversation lengths $L$ (meaning the number of utterances used to average the anonymized speaker's embedding profile).

\begin{table}[H]
\centering
\scriptsize
\setlength{\tabcolsep}{3pt}
\renewcommand{\arraystretch}{0.90}
\begin{tabular}{llcccc}
\hline
$L$ & $N'$ & Original & ASRBN & STTTS & NAC \\
\hline
1 & 5  & $0.9690 \pm 0.0150$ & $0.6198 \pm 0.0762$ 
        & $0.6698 \pm 0.0420$ & $0.7069 \pm 0.0475$ \\
  & 10 & $0.9526 \pm 0.0132$ & $0.4776 \pm 0.0572$ 
        & $0.5690 \pm 0.0386$ & $0.5819 \pm 0.0566$ \\
  & 15 & $0.9422 \pm 0.0166$ & $0.4155 \pm 0.0418$ 
        & $0.4897 \pm 0.0691$ & $0.5302 \pm 0.0396$ \\
  & 20 & $0.9431 \pm 0.0190$ & $0.3741 \pm 0.0503$ 
        & $0.4388 \pm 0.0563$ & $0.4733 \pm 0.0521$ \\
  & 25 & $0.9319 \pm 0.0221$ & $0.3405 \pm 0.0452$ 
        & $0.4181 \pm 0.0602$ & $0.4190 \pm 0.0619$ \\
  & 30 & $0.9284 \pm 0.0251$ & $0.3267 \pm 0.0501$ 
        & $0.3897 \pm 0.0546$ & $0.4448 \pm 0.0682$ \\
  & 40 & $0.9241 \pm 0.0184$ & $0.2784 \pm 0.0476$ 
        & $0.3690 \pm 0.0618$ & $0.3638 \pm 0.0478$ \\
  & 50 & $0.9103 \pm 0.0208$ & $0.2595 \pm 0.0588$ 
        & $0.3190 \pm 0.0686$ & $0.3457 \pm 0.0464$ \\
\hline
3 & 5  & $0.9767 \pm 0.0099$ & $0.7388 \pm 0.0515$ 
        & $0.7810 \pm 0.0405$ & $0.8500 \pm 0.0408$ \\
  & 10 & $0.9672 \pm 0.0132$ & $0.6466 \pm 0.0699$ 
        & $0.6784 \pm 0.0540$ & $0.7440 \pm 0.0580$ \\
  & 15 & $0.9560 \pm 0.0207$ & $0.5793 \pm 0.0417$ 
        & $0.6440 \pm 0.0420$ & $0.7147 \pm 0.0441$ \\
  & 20 & $0.9474 \pm 0.0228$ & $0.5371 \pm 0.0575$ 
        & $0.5853 \pm 0.0474$ & $0.6810 \pm 0.0458$ \\
  & 25 & $0.9388 \pm 0.0200$ & $0.5224 \pm 0.0554$ 
        & $0.5767 \pm 0.0376$ & $0.6474 \pm 0.0464$ \\
  & 30 & $0.9362 \pm 0.0238$ & $0.5034 \pm 0.0403$ 
        & $0.5526 \pm 0.0380$ & $0.6509 \pm 0.0393$ \\
  & 40 & $0.9302 \pm 0.0221$ & $0.4690 \pm 0.0596$ 
        & $0.5224 \pm 0.0482$ & $0.5784 \pm 0.0524$ \\
  & 50 & $0.9284 \pm 0.0175$ & $0.4310 \pm 0.0488$ 
        & $0.4957 \pm 0.0534$ & $0.5534 \pm 0.0493$ \\
\hline
10 & 5  & $0.9741 \pm 0.0128$ & $0.8388 \pm 0.0350$ 
         & $0.8586 \pm 0.0265$ & $0.9026 \pm 0.0337$ \\
   & 10 & $0.9612 \pm 0.0143$ & $0.7483 \pm 0.0461$ 
         & $0.7560 \pm 0.0290$ & $0.8681 \pm 0.0310$ \\
   & 15 & $0.9586 \pm 0.0127$ & $0.7207 \pm 0.0339$ 
         & $0.7086 \pm 0.0373$ & $0.8190 \pm 0.0387$ \\
   & 20 & $0.9526 \pm 0.0171$ & $0.6810 \pm 0.0458$ 
         & $0.6828 \pm 0.0246$ & $0.7905 \pm 0.0394$ \\
   & 25 & $0.9491 \pm 0.0176$ & $0.6698 \pm 0.0359$ 
         & $0.6586 \pm 0.0352$ & $0.7733 \pm 0.0390$ \\
   & 30 & $0.9422 \pm 0.0166$ & $0.6414 \pm 0.0297$ 
         & $0.6284 \pm 0.0347$ & $0.7733 \pm 0.0315$ \\
   & 40 & $0.9379 \pm 0.0138$ & $0.6009 \pm 0.0457$ 
         & $0.6250 \pm 0.0288$ & $0.7474 \pm 0.0460$ \\
   & 50 & $0.9293 \pm 0.0163$ & $0.5776 \pm 0.0384$ 
         & $0.6009 \pm 0.0274$ & $0.7147 \pm 0.0364$ \\
\hline
30 & 5  & $0.9670 \pm 0.0134$ & $0.8193 \pm 0.0410$ 
         & $0.8159 \pm 0.0366$ & $0.9159 \pm 0.0314$ \\
   & 10 & $0.9614 \pm 0.0178$ & $0.7330 \pm 0.0448$ 
         & $0.7318 \pm 0.0334$ & $0.8670 \pm 0.0375$ \\
   & 15 & $0.9545 \pm 0.0190$ & $0.6773 \pm 0.0480$ 
         & $0.6898 \pm 0.0361$ & $0.8284 \pm 0.0233$ \\
   & 20 & $0.9398 \pm 0.0230$ & $0.6614 \pm 0.0336$ 
         & $0.6705 \pm 0.0384$ & $0.7989 \pm 0.0307$ \\
   & 25 & $0.9227 \pm 0.0182$ & $0.6375 \pm 0.0309$ 
         & $0.6466 \pm 0.0264$ & $0.7852 \pm 0.0325$ \\
   & 30 & $0.9239 \pm 0.0194$ & $0.6216 \pm 0.0281$ 
         & $0.6239 \pm 0.0254$ & $0.7727 \pm 0.0321$ \\
   & 40 & $0.9216 \pm 0.0183$ & $0.5875 \pm 0.0290$ 
         & $0.6068 \pm 0.0250$ & $0.7455 \pm 0.0310$ \\
   & 50 & $0.9068 \pm 0.0142$ & $0.5750 \pm 0.0217$ 
         & $0.5898 \pm 0.0183$ & $0.7318 \pm 0.0223$ \\
\hline
\end{tabular}
\caption{Linkability results for conversation lengths $L$ and 
enrollment candidate sizes $N'$. Values are mean $\pm$ std across 
random runs. Lower values indicate lower linkability risk.}
\label{tab:linkability_all_results}
\end{table}

As seen from Table \ref{tab:linkability_all_results} and further illustrated in Figure~\ref{fig:linkability_plots}, the original condition (meaning the original matching of speakers' utterances with the enrollment profiles) remains highly linkable throughout the different experimental conditions. For example, for $L=1$, Linkability decreases only from $0.9690$ to $0.9103$ as $N'$ grows from 5 to 50. Compared to the original setting, anonymization systems substantially reduce Linkability. At $L=1$ and $N'=50$, scores drop to $0.2595$ for ASRBN, $0.3190$ for STTTS, and $0.3457$ for NAC.

Moreover, the effect of conversation length is clearly visible both in the table and the figure diagram. For ASRBN at $N'=50$, Linkability rises from $0.2595$ at $L=1$ to $0.4310$ at $L=3$, then notes $0.5776$ at $L=10$, and $0.5750$ at $L=30$. This trend confirms that aggregating more utterances increases the stability of speaker-specific information. In the figure, the same pattern is observed, with the anonymized systems showing substantially higher scores for $L=10$ than for $L=1$.

ASRBN consistently produces the lowest Linkability scores, while NAC shows the highest residual risk, especially at longer conversation lengths. At $L=30$ and $N'=50$, Linkability scores are $0.5750$ for ASRBN, $0.5898$ for STTTS, and $0.7318$ for NAC. An interesting notice is that there is a reversal from the EER and ZEBRA rankings where NAC appeared as the second best system.

Singling Out results in Table~\ref{tab:singling_out_all_results} and Figure~\ref{fig:singling_out_plots} follow the same pattern. The original condition shows consistently high isolation risk across all settings, remaining between $0.8716$ and $0.9110$ for $L=1$ and above $0.92$ for $L=3$ and $L=10$.

\begin{figure}[H]
    \centering
    \includegraphics[width=\textwidth]{singling_out_plots.pdf}
    \caption{Singling Out score as a function of speaker population 
    size $N$ for different conversation lengths $L$. The dotted line 
    indicates the random predicate $\exp(-1) \approx 0.3679$.}
    \label{fig:singling_out_plots}
\end{figure}

\begin{table}[H]
\centering
\small
\renewcommand{\arraystretch}{1.02}
\resizebox{\textwidth}{!}{%
\begin{tabular}{llcccc}
\hline
$L$ & $N$ & Original & ASRBN & STTTS & NAC \\
\hline
1 & 20 & $0.9110 \pm 0.0159$ & $0.3790 \pm 0.0337$ 
        & $0.4020 \pm 0.0286$ & $0.3880 \pm 0.0196$ \\
  & 30 & $0.8820 \pm 0.0232$ & $0.3760 \pm 0.0273$ 
        & $0.4020 \pm 0.0251$ & $0.4000 \pm 0.0165$ \\
  & 40 & $0.8895 \pm 0.0186$ & $0.3910 \pm 0.0288$ 
        & $0.3945 \pm 0.0138$ & $0.3860 \pm 0.0222$ \\
  & 50 & $0.8716 \pm 0.0124$ & $0.3676 \pm 0.0105$ 
        & $0.4000 \pm 0.0083$ & $0.3928 \pm 0.0179$ \\
\hline
3 & 20 & $0.9413 \pm 0.0378$ & $0.4642 \pm 0.0333$ 
        & $0.5368 \pm 0.0404$ & $0.5052 \pm 0.0863$ \\
  & 30 & $0.9370 \pm 0.0231$ & $0.4264 \pm 0.0278$ 
        & $0.5009 \pm 0.0438$ & $0.4693 \pm 0.0369$ \\
  & 40 & $0.9293 \pm 0.0319$ & $0.3940 \pm 0.0535$ 
        & $0.5120 \pm 0.0258$ & $0.4320 \pm 0.0143$ \\
  & 50 & $0.9292 \pm 0.0133$ & $0.3978 \pm 0.0242$ 
        & $0.4798 \pm 0.0167$ & $0.4320 \pm 0.0333$ \\
\hline
10 & 20 & $0.9533 \pm 0.0194$ & $0.5533 \pm 0.0287$ 
         & $0.6033 \pm 0.0662$ & $0.7267 \pm 0.0573$ \\
   & 30 & $0.9756 \pm 0.0178$ & $0.4867 \pm 0.0285$ 
         & $0.5911 \pm 0.0447$ & $0.6289 \pm 0.0736$ \\
   & 40 & $0.9233 \pm 0.0153$ & $0.4717 \pm 0.0427$ 
         & $0.5783 \pm 0.0239$ & $0.6083 \pm 0.0444$ \\
\hline
\end{tabular}%
}
\caption{Singling Out results for conversation lengths $L$ and speaker 
population sizes $N$. Values are mean $\pm$ std across random runs. 
Random-predicate baseline: $\exp(-1) \approx 0.3679$.}
\label{tab:singling_out_all_results}
\end{table}

All anonymization systems substantially seem to reduce Singling Out compared to the original utterances. More specifically, at $L=1$ and $N=50$, scores drop to $0.3676$ for ASRBN, $0.4000$ for STTTS, and $0.3928$ for NAC, approaching the random baseline of $\exp(-1) \approx 0.3679$.  As $L$ increases, scores rise above the random baseline for all systems: for ASRBN at $N=20$, Singling Out increases from $0.3790$ at $L=1$ to $0.4642$ at $L=3$ and $0.5533$ at $L=10$. ASRBN again provides the strongest protection, while NAC shows the highest residual risk at longer conversation lengths. At $L=10$ and $N=20$, scores are $0.5533$ for ASRBN, $0.6033$ for STTTS, and $0.7267$ for NAC, further supporting the Linkability finding that NAC 
preserves more residual speaker information than EER and ZEBRA suggest.

% Lower values indicate lower residual re-identification risk. For Linkability, the chance level depends on the number of enrollment candidates and is equal to $1/N'$. For Singling Out, the trivial random-predicate baseline is $\exp(-1) \approx 0.3679$.

% Across all systems, the evaluation was performed on 58 common speakers between the pooled enrollment and trial partitions. The pooled trial set contained 80 speakers in total, while the intersection with the enrollment set resulted in 58 usable speakers. The number of available trial utterances per common speaker ranged from 14 to 73, with a mean of 40.43 utterances. This limited the number of speakers available for longer conversation lengths, especially for the Singling Out evaluation at $L=30$.

% The Linkability results are reported in Table~\ref{tab:linkability_all_results}. Overall, the results follow the same trend as reported by ~\cite{vauquier25_interspeech}, meaning that as the number of enrollment candidates $N'$ increases Linkability decreases, while it generally increases when more trial utterances are available for the construction of the test speaker representation.

% The original, non-anonymized condition consistently produced the highest Linkability scores across all values of $L$ and $N'$. Even when the enrollment candidate set increased from $N'=5$ to $N'=50$, the original condition remained highly linkable. For example, for $L=1$, Linkability decreased only from $0.9690$ to $0.9103$, showing that the original speaker embeddings remain strongly identifiable even among a large number of candidates. On the other side, all anonymization systems substantially reduced Linkability compared with the original condition. In ASRBN (best performing system) for $L=1$ and $N'=50$, the Linkability score decreased from $0.9103$ in the original condition to $0.2595$, for STTTS it went to $0.3190$, and for NAC to $0.3457$.

% The effect of conversation length is clearly visible across anonymization systems. For all anonymized conditions, Linkability is lower for $L=1$ than for longer values of $L$, indicating that short trial segments provide less stable speaker information for linkage. For example, for ASRBN at $N'=50$, Linkability increases from $0.2595$ at $L=1$ to $0.4310$ at $L=3$, $0.5776$ at $L=10$, and $0.5750$ at $L=30$. A similar trend is observed for STTTS and NAC. This confirms that aggregating more utterances increases the stability of speaker specific information and makes linkage easier.

% \begin{figure}[H]
%     \centering
%     \includegraphics[width=\textwidth]{linkability_plots.pdf}
%     \caption{Linkability as a function of the number of enrollment speakers $N'$ for different conversation lengths $L$.}
%     \label{fig:linkability_plots}
% \end{figure}

% Figure~\ref{fig:linkability_plots} further illustrates the same pattern observed in Table~\ref{tab:linkability_all_results}. Linkability decreases as the number of enrollment candidates $N'$ increases, since the attacker must identify the correct speaker among a larger candidate set. At the same time, Linkability increases with conversation length $L$, with the anonymized systems showing substantially higher scores for $L=10$ than for $L=1$.

% Across systems, ASRBN generally produces the lowest Linkability scores and therefore the strongest protection according to this metric. STTTS usually follows closely, while NAC often shows the highest residual Linkability among the anonymized systems, especially for longer conversation lengths. For example, at $L=30$ and $N'=50$, the Linkability scores are $0.5750$ for ASRBN, $0.5898$ for STTTS, and $0.7318$ for NAC. This suggests that, under the Linkability metric, NAC preserves more speaker-linking information than the other two anonymization systems in longer-conversation settings, situation that differs from its behavior under EER and ZEBRA, where NAC appeared to perform as the second-best privacy-preserving system. 

% %This discrepancy highlights that Linkability captures a different aspect of residual privacy risk, namely the ability to associate trial and enrollment speaker representations across datasets, rather than only verification performance or score calibration behavior.
% %%%%%%%%% TABLE LINKABILITY
% \begin{table}[H]
% \centering
% \scriptsize
% \setlength{\tabcolsep}{3pt}
% \renewcommand{\arraystretch}{0.90}

% \begin{tabular}{llcccc}
% \hline
% $L$ & $N'$ & Original & ASRBN & STTTS & NAC \\
% \hline
% 1 & 5  & $0.9690 \pm 0.0150$ & $0.6198 \pm 0.0762$ & $0.6698 \pm 0.0420$ & $0.7069 \pm 0.0475$ \\
%   & 10 & $0.9526 \pm 0.0132$ & $0.4776 \pm 0.0572$ & $0.5690 \pm 0.0386$ & $0.5819 \pm 0.0566$ \\
%   & 15 & $0.9422 \pm 0.0166$ & $0.4155 \pm 0.0418$ & $0.4897 \pm 0.0691$ & $0.5302 \pm 0.0396$ \\
%   & 20 & $0.9431 \pm 0.0190$ & $0.3741 \pm 0.0503$ & $0.4388 \pm 0.0563$ & $0.4733 \pm 0.0521$ \\
%   & 25 & $0.9319 \pm 0.0221$ & $0.3405 \pm 0.0452$ & $0.4181 \pm 0.0602$ & $0.4190 \pm 0.0619$ \\
%   & 30 & $0.9284 \pm 0.0251$ & $0.3267 \pm 0.0501$ & $0.3897 \pm 0.0546$ & $0.4448 \pm 0.0682$ \\
%   & 40 & $0.9241 \pm 0.0184$ & $0.2784 \pm 0.0476$ & $0.3690 \pm 0.0618$ & $0.3638 \pm 0.0478$ \\
%   & 50 & $0.9103 \pm 0.0208$ & $0.2595 \pm 0.0588$ & $0.3190 \pm 0.0686$ & $0.3457 \pm 0.0464$ \\

% \hline
% 3 & 5  & $0.9767 \pm 0.0099$ & $0.7388 \pm 0.0515$ & $0.7810 \pm 0.0405$ & $0.8500 \pm 0.0408$ \\
%   & 10 & $0.9672 \pm 0.0132$ & $0.6466 \pm 0.0699$ & $0.6784 \pm 0.0540$ & $0.7440 \pm 0.0580$ \\
%   & 15 & $0.9560 \pm 0.0207$ & $0.5793 \pm 0.0417$ & $0.6440 \pm 0.0420$ & $0.7147 \pm 0.0441$ \\
%   & 20 & $0.9474 \pm 0.0228$ & $0.5371 \pm 0.0575$ & $0.5853 \pm 0.0474$ & $0.6810 \pm 0.0458$ \\
%   & 25 & $0.9388 \pm 0.0200$ & $0.5224 \pm 0.0554$ & $0.5767 \pm 0.0376$ & $0.6474 \pm 0.0464$ \\
%   & 30 & $0.9362 \pm 0.0238$ & $0.5034 \pm 0.0403$ & $0.5526 \pm 0.0380$ & $0.6509 \pm 0.0393$ \\
%   & 40 & $0.9302 \pm 0.0221$ & $0.4690 \pm 0.0596$ & $0.5224 \pm 0.0482$ & $0.5784 \pm 0.0524$ \\
%   & 50 & $0.9284 \pm 0.0175$ & $0.4310 \pm 0.0488$ & $0.4957 \pm 0.0534$ & $0.5534 \pm 0.0493$ \\

% \hline
% 10 & 5  & $0.9741 \pm 0.0128$ & $0.8388 \pm 0.0350$ & $0.8586 \pm 0.0265$ & $0.9026 \pm 0.0337$ \\
%    & 10 & $0.9612 \pm 0.0143$ & $0.7483 \pm 0.0461$ & $0.7560 \pm 0.0290$ & $0.8681 \pm 0.0310$ \\
%    & 15 & $0.9586 \pm 0.0127$ & $0.7207 \pm 0.0339$ & $0.7086 \pm 0.0373$ & $0.8190 \pm 0.0387$ \\
%    & 20 & $0.9526 \pm 0.0171$ & $0.6810 \pm 0.0458$ & $0.6828 \pm 0.0246$ & $0.7905 \pm 0.0394$ \\
%    & 25 & $0.9491 \pm 0.0176$ & $0.6698 \pm 0.0359$ & $0.6586 \pm 0.0352$ & $0.7733 \pm 0.0390$ \\
%    & 30 & $0.9422 \pm 0.0166$ & $0.6414 \pm 0.0297$ & $0.6284 \pm 0.0347$ & $0.7733 \pm 0.0315$ \\
%    & 40 & $0.9379 \pm 0.0138$ & $0.6009 \pm 0.0457$ & $0.6250 \pm 0.0288$ & $0.7474 \pm 0.0460$ \\
%    & 50 & $0.9293 \pm 0.0163$ & $0.5776 \pm 0.0384$ & $0.6009 \pm 0.0274$ & $0.7147 \pm 0.0364$ \\

% \hline
% 30 & 5  & $0.9670 \pm 0.0134$ & $0.8193 \pm 0.0410$ & $0.8159 \pm 0.0366$ & $0.9159 \pm 0.0314$ \\
%    & 10 & $0.9614 \pm 0.0178$ & $0.7330 \pm 0.0448$ & $0.7318 \pm 0.0334$ & $0.8670 \pm 0.0375$ \\
%    & 15 & $0.9545 \pm 0.0190$ & $0.6773 \pm 0.0480$ & $0.6898 \pm 0.0361$ & $0.8284 \pm 0.0233$ \\
%    & 20 & $0.9398 \pm 0.0230$ & $0.6614 \pm 0.0336$ & $0.6705 \pm 0.0384$ & $0.7989 \pm 0.0307$ \\
%    & 25 & $0.9227 \pm 0.0182$ & $0.6375 \pm 0.0309$ & $0.6466 \pm 0.0264$ & $0.7852 \pm 0.0325$ \\
%    & 30 & $0.9239 \pm 0.0194$ & $0.6216 \pm 0.0281$ & $0.6239 \pm 0.0254$ & $0.7727 \pm 0.0321$ \\
%    & 40 & $0.9216 \pm 0.0183$ & $0.5875 \pm 0.0290$ & $0.6068 \pm 0.0250$ & $0.7455 \pm 0.0310$ \\
%    & 50 & $0.9068 \pm 0.0142$ & $0.5750 \pm 0.0217$ & $0.5898 \pm 0.0183$ & $0.7318 \pm 0.0223$ \\

% \hline
% \end{tabular}

% \caption{Linkability results for the evaluated conversation lengths $L$ and enrollment candidate sizes $N'$. Values are reported as mean $\pm$ standard deviation across random runs. Lower values indicate lower linkability risk.}
% \label{tab:linkability_all_results}
% \end{table}

% As reported in Table~\ref{tab:singling_out_all_results} and shown in Figure~\ref{fig:singling_out_plots}, as far as Singling Out is concerned, the original, non-anonymized condition shows consistently high singling-out risk across all evaluated speaker population sizes and conversation lengths. For $L=1$, the original condition remains between $0.8716$ and $0.9110$, while for $L=3$ and $L=10$ it remains above $0.92$ in almost all settings. On the other hand, all anonymization systems substantially reduce Singling Out compared with the original condition. For example, on  $L=1$, where the anonymized systems are close to the random-predicate baseline $\exp(-1) \approx 0.3679$, shown as the dotted grey line in Figure~\ref{fig:singling_out_plots}. Also, at $N=50$, ASRBN obtains a score of $0.3676$, STTTS obtains a score of $0.4000$, and NAC obtains a score of $0.3928$. This suggests that, in the shortest conversation setting, anonymization reduces the isolation risk.

% The effect of conversation length is also visible. As $L$ increases, the anonymized Singling Out scores move farther away from the random baseline. For example, for ASRBN at $N=20$, Singling Out increases from $0.3790$ at $L=1$ to $0.4642$ at $L=3$ and $0.5533$ at $L=10$. A similar trend is observed for STTTS and NAC. This once again follows the results reported by \cite{vauquier25_interspeech}, where longer conversation lengths lead to higher re-identification risk.

% %%%%%%%

% \begin{figure}[H]
%     \centering
%     \includegraphics[width=\textwidth]{singling_out_plots.pdf}
%     \caption{Singling Out score as a function of the speaker population size $N$ for different conversation lengths $L$. The dotted horizontal line indicates the random predicate $\exp(-1) \approx 0.3679$.}
%     \label{fig:singling_out_plots}
% \end{figure}
% \vspace{-1em}
% Across systems, ASRBN generally shows the lowest Singling Out scores and therefore the strongest protection according to this metric. STTTS is usually slightly higher, while NAC shows the highest residual singling-out risk for longer conversation lengths. This is especially visible for $L=10$, where at $N=20$ the Singling Out scores are $0.5533$ for ASRBN, $0.6033$ for STTTS, and $0.7267$ for NAC. Therefore, similarly to the Linkability results, the Singling Out metric suggests that NAC preserves more residual speaker residual information.

% %%%%%% SINGLING OUT TABLE
% \begin{table}[H]
% \centering
% \small
% \renewcommand{\arraystretch}{1.02}

% \resizebox{\textwidth}{!}{%
% \begin{tabular}{llcccc}
% \hline
% $L$ & $N$ & Original & ASRBN & STTTS & NAC \\
% \hline
% 1 & 20 & $0.9110 \pm 0.0159$ & $0.3790 \pm 0.0337$ & $0.4020 \pm 0.0286$ & $0.3880 \pm 0.0196$ \\
%   & 30 & $0.8820 \pm 0.0232$ & $0.3760 \pm 0.0273$ & $0.4020 \pm 0.0251$ & $0.4000 \pm 0.0165$ \\
%   & 40 & $0.8895 \pm 0.0186$ & $0.3910 \pm 0.0288$ & $0.3945 \pm 0.0138$ & $0.3860 \pm 0.0222$ \\
%   & 50 & $0.8716 \pm 0.0124$ & $0.3676 \pm 0.0105$ & $0.4000 \pm 0.0083$ & $0.3928 \pm 0.0179$ \\

% \hline
% 3 & 20 & $0.9413 \pm 0.0378$ & $0.4642 \pm 0.0333$ & $0.5368 \pm 0.0404$ & $0.5052 \pm 0.0863$ \\
%   & 30 & $0.9370 \pm 0.0231$ & $0.4264 \pm 0.0278$ & $0.5009 \pm 0.0438$ & $0.4693 \pm 0.0369$ \\
%   & 40 & $0.9293 \pm 0.0319$ & $0.3940 \pm 0.0535$ & $0.5120 \pm 0.0258$ & $0.4320 \pm 0.0143$ \\
%   & 50 & $0.9292 \pm 0.0133$ & $0.3978 \pm 0.0242$ & $0.4798 \pm 0.0167$ & $0.4320 \pm 0.0333$ \\

% \hline
% 10 & 20 & $0.9533 \pm 0.0194$ & $0.5533 \pm 0.0287$ & $0.6033 \pm 0.0662$ & $0.7267 \pm 0.0573$ \\
%    & 30 & $0.9756 \pm 0.0178$ & $0.4867 \pm 0.0285$ & $0.5911 \pm 0.0447$ & $0.6289 \pm 0.0736$ \\
%    & 40 & $0.9233 \pm 0.0153$ & $0.4717 \pm 0.0427$ & $0.5783 \pm 0.0239$ & $0.6083 \pm 0.0444$ \\

% \hline
% \end{tabular}%
% }

% \caption{Singling Out results for the evaluated conversation lengths $L$ and speaker population sizes $N$. Values are reported as mean $\pm$ standard deviation across random runs. The random-predicate baseline is $\exp(-1) \approx 0.3679$.}
% \label{tab:singling_out_all_results}
% \end{table}


\subsection{Phonetic, Acoustic, and Linguistic Factors in Speaker Re-identification}
\label{subsec:phonetic_acoustic_linguistic_factors}

This section reports the results of the phonetic and acoustic feature analysis which are presented in four steps:

\begin{enumerate}
    \item Target minus non-target feature space separation is used to test whether each feature space assigns higher similarity to same-speaker trials than to different speaker trials (Table \ref{tab:feature_space_target_non_separation}).
    \item Trial-level correlations with ASV scores (Table \ref{tab:feature_space_similarity_asv}).
    \item Predictive models are used to test how much ASV-score variation can be explained by linguistic and phonetic/acoustic features (Table \ref{tab:linguistic_model_comparison_all}).
    \item Speaker-level correlations with individual EER are reported to explore which specific features may be associated with speaker vulnerability (Table \ref{tab:negative_preservation_feature_eer}).
\end{enumerate}

Initially, in Table~\ref{tab:feature_space_target_non_separation} the clearest evidence that the phonetic feature spaces preserve speaker-discriminative structure is provided. There, larger positive values indicate that the feature space assigns higher similarity to the same-speaker trials than to different-speaker trials. Across all trial sets and systems, rhythm, phone duration and trajectory features show consistently positive target minus non-target separation. Overall, phone duration and trajectory features show the strongest separations, for example in the \texttt{test-m} condition phone duration separation reaches $0.436$ for STTTS, $0.442$ for NAC and $0.406$ for ASRBN, while trajectory separation reaches $0.421$, $0.455$ and $0.415$, respectively. Rhythm also shows consistent separation but slightly lower than the aforementioned collective features. 

In contrast, openSMILE (which includes the combination of the 89 extracted features in the library) and ASV-embedding similarity show much weaker separation. openSMILE values remain close to zero, with the largest value being $0.075$ for \texttt{test-f} STTTS, while ASV-embedding similarity is also close to zero or slightly negative in most cases.
This suggests that rhythm, phone duration and trajectory, retain more same-speaker structure than the broader acoustic or embedding-based similarity measures. This means that after anonymization, same-speaker trial pairs remain more similar than different-speaker trial pairs in the phonetic feature spaces of rhythm, phone durations and trajectory features. However, this does not necessarily mean that these features alone are sufficient for re-identification, but it rather shows that they contain speaker-discriminative information that may provide the attacker with cues.

%%%%%%%
\begin{table}[H]
\centering
\small
\begin{tabular}{llccccc}
\hline
Trial set & System & Rhythm & Phone dur. & Trajectory & openSMILE & ASV emb. \\
\hline
 & STTTS & 0.272 & \color{darkgreen}0.381 & 0.345 & 0.028 & -0.003 \\
\texttt{dev-f}  & NAC   & 0.288 & 0.359 & \color{darkgreen}0.400 & 0.004 & 0.007 \\
  & ASRBN & 0.223 & 0.381 & \color{darkgreen}0.406 & 0.022 & 0.002 \\
\hline
 & STTTS & 0.316 & \color{darkgreen}0.380 & 0.373 & 0.005 & 0.016 \\
\texttt{dev-m}  & NAC   & 0.342 & \color{darkgreen}0.362 & 0.359 & 0.044 & 0.007 \\
  & ASRBN & 0.275 & 0.369 & \color{darkgreen}0.381 & 0.028 & -0.002 \\
\hline
& STTTS & 0.326 & \color{darkgreen}0.348 & 0.309 & 0.075 & -0.006 \\
\texttt{test-f} & NAC   & 0.334 & 0.343 & \color{darkgreen}0.360 & 0.051 & -0.000 \\
 & ASRBN & 0.276 & 0.347 & \color{darkgreen}0.391 & 0.050 & -0.005 \\

\hline
 & STTTS & 0.356 & \color{darkgreen}0.436 & 0.421 & 0.049 & 0.021 \\
\texttt{test-m} & NAC   & 0.366 & 0.442 & \color{darkgreen}0.455 & 0.060 & -0.009 \\
 & ASRBN & 0.316 & 0.406 & \color{darkgreen}0.415 & 0.031 & -0.002 \\
\hline
\end{tabular}
\caption{Target minus non-target feature-space cosine similarity separation across systems and trial sets. ASV emb. denotes ASV embedding similarity.}
\label{tab:feature_space_target_non_separation}
\end{table}
%%%%%%%

In Table~\ref{tab:feature_space_similarity_asv} the Spearman correlations between feature-space similarity and ASV scores are reported. The correlations are generally positive but modest. This means that higher phonetic similarity is often associated with higher ASV scores, but the relationship is not strong enough to fully explain the ASV system's behaviour.
Phone duration similarity shows the most consistent positive association with ASV scores, especially for ASRBN, where it reaches $0.221$ on \texttt{dev-f}, $0.224$ on \texttt{dev-m} and $0.218$ on \texttt{test-m}. Rhythm also shows moderate positive correlations in several conditions, such as $0.185$ for \texttt{dev-m} ASRBN and $0.170$ for \texttt{test-f} NAC.

%%%%%%%%%%
\begin{table}[H]
\centering
\small
\begin{tabular}{llccccc}
\hline
Trial set & System & Rhythm & Phone dur. & Trajectory & openSMILE & ASV emb. \\
\hline
  & STTTS & 0.090 & \color{darkgreen}0.128 & 0.065 & 0.009 & 0.102 \\
\texttt{dev-f} & ASRBN & 0.155 & \color{darkgreen}0.221 & 0.062 & -0.002 & 0.051 \\
  & NAC   & 0.149 & \color{darkgreen}0.203 & 0.047 & -0.043 & -0.023 \\
\hline
  & STTTS & \color{darkgreen}0.165 & 0.147 & 0.079 & -0.008 & 0.203 \\
\texttt{dev-m} & ASRBN & 0.185 & \color{darkgreen}0.224 & 0.074 & 0.046 & 0.066 \\
  & NAC   & 0.167 & \color{darkgreen}0.175 & 0.051 & 0.003 & 0.017 \\
\hline
 & STTTS & \color{darkgreen}0.111 & 0.102 & 0.059 & 0.046 & 0.071 \\
\texttt{test-f} & ASRBN & 0.146 & \color{darkgreen}0.148 & 0.044 & -0.035 & -0.126 \\
& NAC   & \color{darkgreen}0.170 & 0.142 & 0.053 & -0.020 & -0.049 \\
\hline
 & STTTS & 0.037 & 0.067 & \color{darkgreen}0.155 & -0.035 & 0.202 \\
\texttt{test-m} & ASRBN & 0.173 & \color{darkgreen}0.218 & 0.036 & 0.000 & 0.021 \\
 & NAC   & \color{darkgreen}0.163 & 0.147 & 0.003 & -0.024 & -0.090 \\
\hline
\end{tabular}
\caption{Spearman correlations between feature-space similarity and ASV score across all trials. ASV emb. denotes ASV embedding similarity.}
\label{tab:feature_space_similarity_asv}
\end{table}

Trajectory correlations with ASV scores are weaker and less consistent, except for \texttt{test-m} STTTS, where trajectory reaches $0.155$. openSMILE correlations are close to zero or negative in most conditions, suggesting that the openSMILE similarity measure used here does not strongly track ASV scores. ASV-embedding similarity is also inconsistent, with low positive values for some STTTS conditions but negative values for NAC in \texttt{dev-f}, \texttt{test-f} and \texttt{test-m}. 

Overall, this table suggests that phonetic similarity is related to ASV scoring, particularly through phone duration and rhythm, but only partially. Phone duration preservation may be contributing to the way target and non-target trials are matched by the ASV system. However, the correlations remain modest, meaning that phone duration similarity is not sufficient on its own to explain ASV score behaviour. So the stronger evidence for residual phonetic speaker structure comes from the target minus non-target separation in Table~\ref{tab:feature_space_target_non_separation}, while Table~\ref{tab:feature_space_similarity_asv} shows that this structure only weakly transfers to ASV similarity score prediction.

Next, the predictive contribution of linguistic, phonetic/acoustic and their interaction feature groups is indicated in Table~\ref{tab:linguistic_model_comparison_all}. The $R^2$ values show that phonetic/acoustic features alone explain only a limited proportion of the ASV-score variance. Then, the phonetic/acoustic-only $R^2$ values, ranging from $0.024$ to $0.077$, indicate that these features are not sufficient on their own to model the ASV score. Text-only features generally explain more variance than phonetic/acoustic-only features for NAC and ASRBN, especially in the female trial sets. For example, in \texttt{dev-f} ASRBN, the text-only model reaches $R^2=0.123$, while the phonetic/acoustic-only model reaches $R^2=0.061$.

The combined models consistently outperform the individual feature groups. This indicates that text-based and phonetic/acoustic information contribute complementary information to ASV-score prediction. For example, in \texttt{test-m} ASRBN, the combined model reaches $R^2=0.176$, compared with $0.119$ for text-only and $0.057$ for phonetic/acoustic-only. A similar pattern is observed for NAC, where the combined model reaches $0.160$ on \texttt{test-m}, compared with $0.121$ for text-only and $0.041$ for phonetic/acoustic-only. However, adding interaction terms produces only very small improvements over the combined model. This suggests that the interaction between text similarity and phonetic/acoustic similarity does not improve the ASV-score prediction beyond the additive combination of both feature groups. 

Last but not least, the best AUC column shows that the same feature groups are more successful at distinguishing target from non-target trials than at predicting the exact ASV score. All AUC values are above $0.80$, with the highest value observed for \texttt{test-m} STTTS ($0.868$).

Therefore, the feature groups contain enough information to support target/non-target separation within this trial-level evaluation setup, even though they explain only a limited amount of the ASV-score variance. As a result, this table supports the interpretation that phonetic/acoustic features are not sufficient to fully predict ASV scores, but they contribute useful information when combined with text features and support target/non-target classification within the evaluated trial pairs.
%%%%%%%

\begin{table}[H]
\centering
\small
\resizebox{\textwidth}{!}{%
\begin{tabular}{llccccc}
\hline
Trial set & System & Text only & Phon./Acous. only & Combined & + Int. & Best AUC \\
\hline
  & STTTS & 0.031 & 0.025 & 0.051 & \color{darkgreen}0.052 & 0.819 \\
\texttt{dev-f}  & NAC   & 0.100 & 0.053 & 0.137 & \color{darkgreen}0.138 & 0.826 \\
  & ASRBN & 0.123 & 0.061 & 0.163 & \color{darkgreen}0.165 & 0.815 \\

\hline
  & STTTS & 0.041 & 0.069 & \color{darkgreen}0.108 & \color{darkgreen}0.108 & 0.838 \\
\texttt{dev-m}  & NAC   & 0.103 & 0.049 & \color{darkgreen}0.155 & \color{darkgreen}0.155 & 0.818 \\
  & ASRBN & 0.068 & 0.077 & 0.145 & \color{darkgreen}0.147 & 0.813 \\
\hline
 & STTTS & 0.016 & 0.024 & 0.041 & \color{darkgreen}0.044 & 0.801 \\
\texttt{test-f} & NAC   & 0.105 & 0.043 & 0.142 & \color{darkgreen}0.144 & 0.824 \\
 & ASRBN & 0.115 & 0.051 & 0.156 & \color{darkgreen}0.157 & 0.813 \\
\hline
 & STTTS & 0.016 & 0.069 & 0.087 & \color{darkgreen}0.088 & 0.868 \\
\texttt{test-m} & NAC   & 0.121 & 0.041 & 0.160 & \color{darkgreen}0.163 & 0.841 \\
 & ASRBN & 0.119 & 0.057 & 0.176 & \color{darkgreen}0.176 & 0.825 \\
\hline
\end{tabular}%
}
\caption{Five-fold trial-level cross-validated comparison of linguistic and phonetic/acoustic feature groups across anonymization systems and trial sets. The Text only, Phon./Acous. only, Combined, and + Int. columns report $R^2$ for \textbf{ASV score prediction}. Combined denotes the model using both text and phonetic/acoustic features. + Int. denotes the model with text, phonetic/acoustic features, and interaction terms. Best AUC reports the highest ROC-AUC obtained for target/non-target classification.}
\label{tab:linguistic_model_comparison_all}
\end{table}

Table~\ref{tab:negative_preservation_feature_eer} focuses on phonetic preservation features that are negatively correlated with individual speaker EER. Spearman correlation is used because the relationship between phonetic preservation features and EER is not assumed to be linear. The purpose of the table is to identify phonetic properties whose preservation after anonymization may be associated with residual re-identification. However, given the small number of speakers per condition ($n=13$--$16$), the absence of inferential significance testing, and the lack of correction for multiple comparisons, the reported coefficients should be interpreted only as descriptive and exploratory point estimates, rather than as statistically confirmed associations or reliable differences across systems.

%%%%%%%
\setlength{\LTcapwidth}{\linewidth}

\begingroup
\scriptsize
\setlength{\tabcolsep}{3pt}
\renewcommand{\arraystretch}{1.08}

\begin{longtable}{p{1.2cm}p{1.2cm}p{1.3cm}p{7.2cm}cc}
\hline
Trial set & System & Group & Feature & $n$ & $\rho$ \\
\hline
\endfirsthead

\hline
Trial set & System & Group & Feature & $n$ & $\rho$ \\
\hline
\endhead

\hline
\endfoot

\hline
\\
\caption{Phonetic preservation features negatively correlated (Spearman) with speaker-level EER, by system and trial set. Negative correlations indicate that stronger preservation is associated with lower EER.}
\label{tab:negative_preservation_feature_eer}\\[0.8em]
\endlastfoot

\hline
\texttt{dev-f} & STTTS & Rhythm & Vocalic nPVI preservation & 15 & -0.164 \\

\hline
\texttt{dev-f} & NAC & Phone d. & Phone-duration pattern preservation & 15 & -0.557 \\
 & NAC & Traj. & F3 preservation & 15 & -0.471 \\
 & NAC & Phone d. & Phone-duration pattern preservation & 15 & -0.243 \\
 & NAC & Rhythm &  Vocalic rPVI preservation & 15 & -0.196 \\
 & NAC & Rhythm & Mean consonant duration preservation & 15 & -0.161 \\

\hline
\texttt{dev-f} & ASRBN & Rhythm & Vocalic $\Delta V$ preservation & 15 & -0.357 \\
 & ASRBN & Traj. & F1 preservation & 15 & -0.314 \\
 & ASRBN & Traj. & F3 preservation & 15 & -0.296 \\
 & ASRBN & Traj. & F0 preservation & 15 & -0.154 \\

\hline
\texttt{dev-m} & STTTS & Rhythm & Percentage of consonantal material preservation & 14 & -0.510 \\
 & STTTS & Traj. & F2 preservation & 14 & -0.317 \\

\hline
\texttt{dev-m} & NAC & Traj. & F2 preservation & 14 & -0.116 \\

\hline
\texttt{dev-m} & ASRBN & Rhythm & Vocalic $\Delta V$ preservation & 14 & -0.371 \\
 & ASRBN & Rhythm & Consonantal $\Delta C$ preservation & 14 & -0.270 \\
 & ASRBN & Rhythm & Percentage of consonantal material preservation & 14 & -0.116 \\
 & ASRBN & Rhythm & Percentage of vocalic material preservation & 14 & -0.116 \\

\hline
\texttt{test-f} & STTTS & Rhythm & Vocalic nPVI preservation & 16 & -0.512 \\
 & STTTS & Traj. & F2 preservation & 16 & -0.253 \\
 
\hline
\texttt{test-f} & NAC & Traj. & F0 preservation  & 16 & -0.227 \\
 & NAC & Rhythm & Consonantal nPVI preservation & 16 & -0.139 \\

\hline
\texttt{test-f} & ASRBN & Rhythm & Consonantal nPVI preservation & 16 & -0.468 \\
 & ASRBN & Rhythm & Percentage of vocalic material preservation & 16 & -0.288 \\
 & ASRBN & Rhythm & Percentage of consonantal material preservation & 16 & -0.288 \\

\hline
\texttt{test-m} & STTTS & Traj. & F1 trajectory preservation & 13 & -0.747 \\
 & STTTS & Rhythm & Consonantal rPVI preservation & 13 & -0.401 \\
 & STTTS & Rhythm & Vocalic nPVI preservation & 13 & -0.286 \\
 & STTTS & Rhythm & Vocalic Varco preservation & 13 & -0.247 \\
 & STTTS & Rhythm & Consonantal nPVI preservation & 13 & -0.247 \\

\hline
\texttt{test-m }& NAC & Phone d. & Phone-duration pattern preservation & 13 & -0.527 \\
 & NAC & Rhythm & Vocalic rPVI preservation  & 13 & -0.473 \\
 & NAC & Traj. & F0  preservation & 13 & -0.198 \\
 & NAC & Rhythm & Vocalic nPVI preservation & 13 & -0.159 \\
 & NAC & Rhythm & Percentage of consonantal material preservation& 13 & -0.148 \\

\hline
\texttt{test-m} & ASRBN & Rhythm & Consonantal rPVI preservation & 13 & -0.269 \\
 & ASRBN & Traj. & F2  preservation & 13 & -0.264 \\
 & ASRBN & Traj. & F3 preservation & 13 & -0.236 \\
 & ASRBN & Rhythm & Vocalic nPVI preservation & 13 & -0.187 \\
 & ASRBN & Rhythm & Consonantal Varco preservation & 13 & -0.176 \\

\end{longtable}
\endgroup

The reported features are preservation measures, for example, a rhythm preservation feature indicates that a given rhythm property remains similar between the original and anonymized speech. Similarly, a trajectory preservation feature indicates that the temporal shape of a pitch or formant contour remains similar after anonymization. Only the negative Spearman correlations are reported because they indicate cases where higher preservation of a phonetic feature is associated with lower speaker EER. In simple terms, the more strongly this feature is preserved for a speaker after anonymization, the easier that speaker may be for the attacker to re-identify.

Several types of phonetic features appear in the table. Rhythm features describe the temporal organization of speech. \textit{Vocalic} and \textit{consonantal percentages} measure how much of an utterance is occupied by vowel-like and consonant-like intervals. Mean vowel and consonant durations describe the average duration of these intervals. $\Delta V$ and $\Delta C$ refer to the standard deviation of vocalic and consonantal interval durations, respectively. \textit{Varco-V} and \textit{Varco-C} are normalized variability measures, computed by normalizing $\Delta V$ and $\Delta C$ by the corresponding mean vowel or consonant duration. Finally, \textit{rPVI} (raw) and \textit{nPVI}  (normalized) measure pairwise variability between successive vocalic or consonantal intervals. Phone-duration features describe the duration pattern of individual phones, and phone-duration pattern preservation means that the relative timing of phones in the anonymized speech remains similar to the original speech. Lastly, trajectory features describe the temporal evolution of acoustic contours. F0 refers to the pitch contour, while F1, F2 and F3 refer to the first three formant trajectories. These formants are related to the resonant structure of speech and are influenced by articulation and vocal tract characteristics. 

The strongest negative association in the table appears for \texttt{test-m} STTTS, where F1 trajectory preservation has a correlation with EER of $\rho=-0.747$. This indicates that, in this condition, speakers whose F1 trajectory is more preserved after anonymization tend to have substantially lower EER. Therefore, F1 trajectory preservation may act as a residual cue for the attacker in this subset. Another strong negative association is observed for \texttt{dev-f} and \texttt{test-m} NAC, where phone-duration pattern preservation has $\rho=-0.557$ and $\rho=-0.527$, respectively. This suggests that preserving the original phone-duration pattern may make some speakers more vulnerable under NAC.

Rhythm preservation features also appear very frequently in the table. For example, \texttt{test-f} STTTS shows a negative correlation for vocalic nPVI preservation ($\rho=-0.512$), while in ASRBN it shows negative correlations for consonantal nPVI preservation ($\rho=-0.468$) and percentage-based rhythm preservation. Rhythm-related features are the most frequent category overall, appearing across almost all systems and trial sets. This suggests that timing structure, especially the relative organization of vocalic and consonantal intervals, may carry residual speaker-specific information after anonymization.

Trajectory features are the second most frequent category. They are especially important in \texttt{test-m} STTTS, where F1 trajectory preservation shows the strongest negative correlation in the table, and in \texttt{dev-f} ASRBN and \texttt{test-m} ASRBN, where formant trajectory preservation also appears. This indicates that the preservation of pitch or formant contour shapes may also contribute to re-identification vulnerability in some conditions. Phone-duration preservation appears less frequently than rhythm and trajectory features, but for NAC phone-duration pattern preservation is strongly negatively correlated with EER in both \texttt{dev-f} and \texttt{test-m}.

The speaker-level heatmaps in Appendix \ref{app:feature-heatmaps} provide a visual complement to Table \ref{tab:negative_preservation_feature_eer}, showing that some vulnerable speakers (top rows, lowest EER) exhibit above-average preservation in phone-duration and trajectory feature spaces. This trend, nevertheless, is not constant across all speakers and systems. Speaker \textbf{6241} in B3 \texttt{dev-m}, for instance, retains high preservation across nearly all feature spaces despite a mid-range EER, while speaker \textbf{84} in NAC \texttt{dev-f} is among the best-protected speakers overall (EER = 61.1\%) yet shows strongly above-average preservation in phone-duration and rhythm, and then speaker \textbf{6241} in ASRBN \texttt{dev-m} shows a similar mismatch in phone-duration and rhythm. 
On the other hand there are speakers with quite low EER scores like \textbf{7021} in ASRBN \texttt{test-m} or \textbf{8297} in STTTS, which do not show consistently high preservation values across the examined feature spaces.
These cases indicate that the link between phonetic preservation and vulnerability is a general tendency rather than a strict rule.
% The results also show that the attacker-helpful features are not identical across systems or trial sets. For STTTS, strong negative associations appear in both rhythm and trajectory features. For NAC, phone-duration pattern preservation is particularly important in dev-f and test-m ASRBN, the negative preservation features are mostly rhythm-related, such as vocalic $\Delta V$, consonantal $\Delta C$, consonantal nPVI, consonantal rPVI and percentage-based rhythm preservation. This variation suggests that each anonymization system may preserve different types of speaker-specific phonetic information.
However, once again, these results should be interpreted carefully, because the number of available speakers per condition is small. Therefore, the correlations are exploratory and neither the table nor the heatmaps prove that these features cause re-identification but they only indicate that phonetic preservation and speaker vulnerability tend to co-occur more often than not.
% but instead show that, for some systems and subsets, speakers with higher preservation of certain phonetic features also tend to have lower EER. 
% Overall, the table suggests that rhythm preservation is the most consistently observed potential vulnerability factor, while trajectory and phone-duration preservation may provide particularly strong attacker-helpful cues in specific conditions.

%%%%%%%
\section{Discussion}
\label{sec:discussion}
This thesis is formed on the basis of four questions. First, it asks in which level the three selected VPC 2024 baseline systems, 
STTTS, NAC, and ASRBN, preserve privacy across individual speakers based on the official adopted by VPC metric, EER, and whether their 
differences become visible (\textbf{RQ1}, Section \ref{subsubsec:baseline_performance_and_eer}). Second, it examines whether alternative privacy metrics proposed in literature, 
i.e. ZEBRA, Linkability \& Singling Out, provide a more detailed view of anonymization performance than EER alone (\textbf{RQ2}, Sections \ref{subsubsec:what_zebra_adds} and \ref{subsubsec:linkabilty_and_singling_out_a_different_threat_metric}). 
Third, it investigates whether speaker-level vulnerability after anonymization is associated with factors such as gender, 
phoneme duration, pitch, rhythm characteristics, formant trajectories and linguistic content (\textbf{RQ3}, Section \ref{subsubsec:phonetic_linguistic_discussion}). 
Finally, it asks what these findings reveal about the remaining weaknesses of current anonymization systems and what should be taken into 
account in order to make voice anonymization systems and their evaluation more robust in the future (\textbf{RQ4}, partially in Section \ref{subsubsec:phonetic_linguistic_discussion} and mainly in Section \ref{subsubsec:implications_for_future_systems}).

The results presented in the previous section address these questions. 
Taken together, they show that privacy performance after voice anonymization cannot be fully understood through a single global score 
and even though the systems differ in their overall behaviour (ASRBN often achieving the strongest global privacy performance and STTTS the weakest) 
the important finding is that some speakers remain comparatively easy to re-identify across systems with different architectures,
while others appear to be protected more effectively, confirming the literature findings \citep{williams_anonymizing_2024,gaznepoglu2025sayexploitinglinguisticcontent} 
and further supporting the statement that global EER can hide individual-level differences. 
The alternative metrics refine rather than simply reproduce this picture: ZEBRA shows that speakers with similar EER values may still differ in the amount of residual identity evidence available to the attacker 
based on \texttt{dECE}, while Linkability \& Singling Out shift the interpretation to the possible linkage or isolation of anonymized speech samples. 
Lastly, the feature-level analyses add interpretation to the metrics' results by showing that residual recognizability
seems to be present and related to the preservation of specific phonetic/acoustic and linguistic properties, though this relationship does not seem to be consistent
across systems and trial sets, thus strong claims should be avoided.  

\subsection{Baseline Performance and EER as an Incomplete Privacy Metric}
\label{subsubsec:baseline_performance_and_eer}
The results in Section~\ref{subsec:results_for_eer_and_zebra} show that global EER, although convenient, can give an incomplete picture of 
anonymization effectiveness. Across all trial sets and systems, global EER values fall within or close to the VPC privacy thresholds. 
ASRBN achieves the strongest global results, with EER values of 34.231\% on \texttt{dev-f} and 32.770\% on \texttt{test-m}, while STTTS performs the
weakest, with 28.265\% and 26.278\% on the same sets. At the population level, these results suggest meaningful privacy protection.

However, the speaker-level EER analysis reveals persistent individual vulnerabilities. On \texttt{dev-f}, speaker \textbf{1673} obtains low EER 
values across all systems, with 10.078\%, 18.651\%, and 14.515\% under STTTS, NAC, and ASRBN respectively. Speaker \textbf{2412} shows a similar 
pattern, remaining below 30\% EER for all systems. On \texttt{dev-m}, speakers \textbf{652}, \textbf{2803}, and \textbf{5536} are also consistently 
vulnerable, with speaker \textbf{2803} reaching EER values as low as 9.174\%, 18.946\%, and 17.102\%. Similar cases appear in the test sets, 
including speakers \textbf{3570} and \textbf{1221} on \texttt{test-f}, and speakers \textbf{1320}, \textbf{4077}, \textbf{7021}, and \textbf{8224} on 
\texttt{test-m}. The fact that these speakers remain vulnerable across fundamentally different architectures suggests that the failures are not only system-specific, 
but may also reflect persistent speaker-dependent properties. This supports previous findings by \cite{williams_anonymizing_2024}, who identified 
speakers that remain difficult to anonymize across systems, and by \cite{gaznepoglu2025sayexploitinglinguisticcontent}, who showed that per-speaker 
EER can be very low even when global EER appears acceptable.

From a privacy-law perspective, such individual failures are important. Under Recital 26 of the GDPR, information is anonymous only when the data 
subject is no longer identifiable, not merely when most speakers are protected \citep{noauthor_recital_nodate}. Global EER cannot capture these 
selective failures because it summarizes performance across speakers. This motivates the use of complementary metrics that also expose individual-level 
privacy leakage.

\subsection{What ZEBRA Adds: Its Divergence from EER}
\label{subsubsec:what_zebra_adds}
The ZEBRA framework with its complementary measures of privacy disclosure operates on the evidence produced by the ASV system rather
than on a single operating point threshold (as described in Section \ref{subsubsec:zebra_framework}). With the \texttt{dECE} quantifying how much identity evidence in bits
the anonymized speech provides to the adversary, and the \texttt{max\_abs\_LLR} capturing the strongest individual piece of evidence observed among all trial scores 
this framework exposes dimensions of privacy that the EER cannot capture.

At the global level, the ordering of systems under \texttt{dECE} is consistent with their ordering under EER. 
ASRBN achieves the lowest global \texttt{dECE} across three of the four trial sets, corresponding to its generally highest global EER,
while STTTS consistently produces the highest global \texttt{dECE}, corresponding to its lowest global EER. 
This might suggest that ZEBRA and EER are acting similarly, but the speaker-level analysis reveals that this 
consistency does not extend to the individual level, and that the two metrics capture meaningfully different aspects of privacy disclosure.

Several speakers show mismatches between speaker-level EER and \texttt{dECE}. For example, on \texttt{dev-f} under NAC, speaker \textbf{6313}
has a low EER of 19.949\%, suggesting weak protection, but a relatively low \texttt{dECE} of 0.087, indicating limited overall evidence leakage. 
In contrast, speaker \textbf{3536} has an EER of 30.278\%, which appears more acceptable, but a higher \texttt{dECE} of 0.201. 
Speaker \textbf{2412} shows a similar pattern, with an EER close to 30\% but a \texttt{dECE} of 0.243. 
Similar cases appear on \texttt{test-f} under NAC, where speaker \textbf{2961} has an EER of 30.407\% but a \texttt{dECE} of 0.274, 
and on \texttt{test-m} under ASRBN, where speaker \textbf{61} has an EER of 30.244\% but a \texttt{dECE} of 0.156. These cases are only some examples
that reflect the different nature of the two metrics. EER measures only a point but it does not describe the full score distribution or measure the identity leakage, unlike \texttt{dECE}.
Moreover, the inherent inclusion of \texttt{max\_abs\_LLR} in the framework highlights the worst-case leakage by capturing high-risk individual trials that are masked by EER.


\subsection{Linkability \& Singling Out: A Different Threat Assessment Metric}
\label{subsubsec:linkabilty_and_singling_out_a_different_threat_metric}

The Linkability \& Singling Out results in Section \ref{subsec:results_linkability_and_singling_out} provide a different view of privacy risk, closer to the legal notion of anonymization. 
Unlike EER and ZEBRA, which are based on ASV similarity scores, Linkability evaluates whether anonymized trial embeddings can still be associated with the correct enrollment speaker within a candidate pool 
and Singling Out evaluates whether an attacker can isolate one individual from a set of anonymized representations, therefore capturing residual speaker information that may not appear as strong ASV 
performance, but still represent a privacy risk under the GDPR framework.

Here, the effect of conversation length seems of high importance across all systems. As $L$ increases, Linkability also increases, 
showing that more speech material produces more stable speaker representations. For example, for NAC at $N'=50$, Linkability rises from $0.3457$ at $L=1$ to $0.7318$ at $L=30$. 
The Singling Out results follow the same tendency with systems close to the random predicate baseline at $L=1$, but clearly exceed it at longer conversation lengths. 
This supports the idea that repeated speech collection creates a realistic privacy risk, because an attacker can aggregate anonymized utterances over time.

The main finding is that system ranking changes when these metrics are used. While NAC ranks second under EER and ZEBRA, it shows the highest residual risk under Linkability 
and Singling Out across most conditions. For example, at $L=30$ and $N'=50$, NAC reaches a Linkability score of $0.7318$, compared with $0.5750$ for 
ASRBN and $0.5898$ for STTTS. Similarly, at $L=10$ and $N=20$, NAC reaches a Singling Out score of $0.7267$, compared with $0.5533$ for ASRBN and $0.6033$ for STTTS. 
This reversal shows that a system can perform well under verification-based metrics while still preserving information that allows utterances from the same speaker to be linked over time.

This behaviour could be connected to NAC's token-based architecture. In this approach, semantic tokens are extracted from the source utterance, while acoustic prompts from pseudospeakers condition the 
generation of acoustic tokens for the anonymized waveform, with the aim to suppress speaker information \citep{panariello_speaker_2024}. 
However, the higher Linkability scores observed here suggest that some stable speaker-specific regularities may remain in the generated token sequences. 
When multiple utterances are aggregated, these residual regularities become more stable, which could explain why NAC is more vulnerable under Linkability, especially for longer conversation lengths.

ASRBN performs best under both Linkability \& Singling Out, which is consistent with the role of vector quantization in its architecture. 
As ASR bottleneck features are not inherently speaker-independent, they can still leak speaker identity \citep{champion_anonymizing_2024}, so
by constraining the ASR bottleneck representation to a limited set of discrete prototypes, vector quantization reduces the amount of speaker-specific variation preserved in the linguistic features.
This may explain why ASRBN remains less linkable than NAC and STTTS when several utterances are aggregated. 

Overall, these results show that privacy evaluation should not rely only on EER. Linkability \& Singling Out capture the accumulation of speaker evidence across repeated utterances, which is highly relevant for real-world deployments. 
Their inclusion in future evaluation pipelines, as suggested by \cite{tomashenko_third_2026}, would make system comparison more aligned with both GDPR oriented privacy requirements and practical attack scenarios.

\subsection{The Role of Gender and Phonetic/Linguistic Features as Residual Identity Cues}
\label{subsubsec:phonetic_linguistic_discussion}

The phonetic, acoustic, and linguistic analysis in Section~\ref{subsec:phonetic_acoustic_linguistic_factors} shows that anonymization does not fully remove speaker-discriminative structure from the speech signal.
As shown in Table \ref{tab:feature_space_target_non_separation}, rhythm, phone-duration, and pitch/formant trajectory features consistently produce positive target minus non-target separation across all systems and trial sets, with the 
strongest separations observed for phone-duration and trajectory features.

The trial-level correlations in Table \ref{tab:feature_space_similarity_asv} further support this interpretation, with phone-duration similarity showing
the most consistent positive association with ASV scores, especially for ASRBN, where it reaches $0.221$ on \texttt{dev-f}, $0.224$ on \texttt{dev-m}, and $0.218$ on \texttt{test-m}.
Rhythm also correlates positively in several conditions, for example $0.185$ for ASRBN on \texttt{dev-m} and $0.173$ on \texttt{test-m}. 
These correlations are modest, so phonetic preservation should not be treated as the only factor driving ASV scores. Rather, they indicate that timing and phonetic structure partially overlap with the information exploited by the attacker.

To continue at the speaker level, Table \ref{tab:negative_preservation_feature_eer} shows that stronger preservation of specific phonetic features is often associated with lower speaker-level EER, meaning
that speakers whose anonymized speech preserves these features more strongly tend to be more vulnerable. The descriptive patterns differ by system, with different feature categories appearing more prominently across conditions.

For STTTS, rhythm and trajectory features are prominent. For example, on \texttt{test-m}, $F_1$ trajectory preservation has a strong negative correlation with EER ($\rho=-0.747$), while consonantal rPVI preservation also correlates negatively ($\rho=-0.401$).
This is plausible given the architecture of STTTS, which retains original phonetic transcriptions and phone durations unmodified, and only randomly perturbs pitch and energy at the phone level \cite{tomashenko2024voiceprivacy2024challengeevaluation}.
It is interesting how in \cite{meyer_prosody_2023} unmodified prosody cloning (duration, pitch, and energy) does not by itself compromise privacy, provided the substituted speaker embedding is sufficiently dissimilar from the original. 
This suggests the possibility that the descriptive associations observed here are related to specific formant and rhythm properties that remain after B3's perturbation scheme, rather than to prosody preservation in general.

For NAC, phone-duration preservation provides some of the largest negative descriptive correlations with speaker-level EER on \texttt{dev-f}. 
Also, \texttt{test-m} has a correlation with EER of ($\rho=-0.527$). This could be possibly interpreted based on the fact that NAC is not explicitly designed to modify phone durations \citep{panariello_speaker_2024}. 
The results therefore suggest that some segmental timing regularities may survive the token-based transformation and become useful when comparing utterances from the same speaker.

For ASRBN, rhythm-related features appear more frequently among the reported negative descriptive correlations. For example, vocalic $\Delta V$ preservation correlates with lower EER on \texttt{dev-f} ($\rho=-0.357$) and \texttt{dev-m} ($\rho=-0.371$), while consonantal nPVI preservation reaches $\rho=-0.468$ on \texttt{test-f}. 
This fits the motivation of ASRBN only in a cautious sense, as ASR bottleneck features may still contain speaker information \citep{champion_anonymizing_2024}. 

The linguistic analysis in Table \ref{tab:linguistic_model_comparison_all} adds another layer to these interpretations. Text-only features seem to explain more ASV-score variance than phonetic/acoustic features alone for NAC and ASRBN in several trial sets. 
This trend seems constant with the findings of \cite{gaznepoglu2025sayexploitinglinguisticcontent}, \cite{franzreb_content_2026} and \cite{franzreb_private_2025} who show that similar linguistic content in enrollment and trial embeddings can be a strong residual identity cue.
For example, on \texttt{dev-f}, ASRBN reaches $R^2=0.123$ with text-only features compared with $0.061$ for phonetic/acoustic features, and NAC reaches $0.100$ compared with $0.053$. 
For STTTS, the contribution is more balanced, and phonetic/acoustic features are sometimes stronger, for example on \texttt{dev-m} and \texttt{test-m}, mainly because STTTS retains the prosody of the speaker in the anonymized version. Across all systems, the combined model outperforms either feature group alone.
The interaction terms add only small improvements, suggesting that linguistic and phonetic/acoustic information contribute mostly additively and not in an interactive way to re-identification, not supporting the claim that \textit{what is said} is relevant to \textit{how it is said}.

The gender-partitioned results further show that vulnerability is present in both female and male speaker sets, but it is not distributed identically across them. At the global level, the male partitions generally show lower EER values than the 
female partitions, especially for STTTS. For example, STTTS decreases from $28.265\%$ on \texttt{dev-f} to $23.788\%$ on \texttt{dev-m}, and from $28.682\%$ on \texttt{test-f} to $26.278\%$ on \texttt{test-m}. 
A similar tendency appears for NAC and ASRBN on the development sets, where both systems obtain higher EER on \texttt{dev-f} than on \texttt{dev-m}. 
The speaker-level results also show that several male speakers remain below the $30\%$ EER threshold across all systems, such as speakers \textbf{652}, \textbf{2803}, and \textbf{5536} on \texttt{dev-m}, and speakers \textbf{1320}, \textbf{4077}, \textbf{7021}, and \textbf{8224} on \texttt{test-m}. 
However, this pattern is not exclusive to male speakers, since female speakers such as \textbf{1673}, \textbf{2412}, \textbf{3570}, and \textbf{1221} also remain consistently vulnerable across systems. These results should not be interpreted as showing a simple gender effect, but rather as indicating 
that gender-partitioned evaluation can reveal differences in how vulnerability is distributed across speaker groups.

\subsection{Implications for Future Anonymization Systems and their Evaluation}
\label{subsubsec:implications_for_future_systems}

The results suggest that future anonymization systems should not aim solely on increasing global EER, but should also adopt the introduction of a broader set of metrics. 
EER still remains useful for comparison with previous VPC results, but the findings of this thesis show that it can hide individual-level vulnerability and metric-dependent changes in system ranking. 
Linkability \& Singling Out and the ZEBRA framework also need to be considered as they incorporate metrics which are based on the legal and forensic frameworks respectively,
and therefore correspond to different aspects of privacy protection. Especially in real-life scenarios where an attacker may collect several anonymized utterances from the same person over time, these metrics provide a more
realistic assessment of privacy risk than EER alone.

Linguistic content can also function as a leakage channel in this experimental setup, yet remains a largely untouched topic by the challenge organisers. The current systems anonymize only \emph{how} something is said, not \emph{what} is
said. Future VPC occurrences should aim for testing other types of datasets with data that imitate real-life scenarios as \cite{franzreb_content_2026} suggest. Moreover, stronger privacy guarantees may need to consider content-level interventions, or at least explicitly report the probability of this leakage rather
than treating it as a solely corpus artefact that the system failed to conceal.

Finally, the results indicate that anonymization should be tested across speaker groups like \cite{williams_anonymizing_2024} and not only through average scores. 
Some speakers remain vulnerable across all three systems, which suggests that future work should include larger-scale and targeted analyses of hard to anonymize 
speakers and use these cases to further guide system design. As far as the scope of this thesis is concerned, B5 or ASRBN appears to be the most robust system overall according to 
all the tested metrics and their individual-level analyses. This is in line with \citet{champion_anonymizing_2024} that constraining ASR bottleneck features with vector 
quantization substantially reduces speaker linkability since a sufficiently small prototype dictionary forces the network to encode a more
speaker-invariant linguistic representation. He showed also that F0 transformations alone provide little privacy gain when the main leakage source lies in the linguistic bottleneck rather
than in pitch, further also supporting the advocators of linguistic leakage. Given ASRBN's comparatively strong and consistent performance across all metrics used in this thesis, future work should prioritize refining this 
quantization-based approach, for instance by extending it to rhythm and duration representations (which emerged as candidate residual factors in the exploratory analysis in Table \ref{tab:negative_preservation_feature_eer}), rather than moving away from it. 
So, a more robust anonymization pipeline would therefore combine a wider, real-life based and naturalistic evaluation dataset with an ASRBN style bottleneck quantization approach extended to more strongly suppress rhythmic cues,
evaluated with methods that reveal both population-level and speaker-level privacy risks. This would be a promising direction for future experimentation.

\section{Conclusions and Future Work}
\label{sec:conclusions}

This thesis examines alternative metrics for assessing the performance of voice anonymization systems, highlights the shortcomings of global EER, tests different metrics available and more suitable for this task and 
lastly investigates if the residual speaker vulnerability can be associated with phonetic/acoustic features and their combination with linguistic information. 
Three VPC 2024 baseline systems, STTTS, NAC, and ASRBN, were evaluated using global and speaker-level EER, ZEBRA, Linkability, and Singling Out, together with the feature-level analysis to provide explainability for the obtained results.

The results showed that global privacy performance can conceal substantial differences between individuals. Although ASRBN generally achieved the strongest privacy protection and STTTS the weakest, 
several speakers remained vulnerable across all three systems. ZEBRA further showed that speakers with similar EER values may differ in the amount of identity evidence disclosed. 
Linkability \& Singling Out revealed that anonymized speech became easier to associate and isolate when several utterances were aggregated. In particular, NAC showed greater residual risk under these metrics than its EER results suggested.

The feature analysis indicated that rhythm, phone-duration patterns, and pitch and formant trajectories preserve speaker-discriminative structure after anonymization. 
Linguistic similarity also contributed to ASV-score prediction, especially when combined with phonetic and acoustic information. However, these findings were not fully consistent across systems and partitions and should therefore be interpreted as exploratory, given also the limited number of speakers in the trial sets.

Overall, the findings demonstrate that voice anonymization should not be assessed through global EER alone given that there are evaluation frameworks available that combine population-level and speaker-level measures and which are evidence-based, correspond to real-life scenarios or are legally grounded. 

Future work should repeat the application of metrics and the feature analyses using larger speaker populations, different languages, taking also into account the latest VPC of 2026 \citep{voiceprivacy2026}, and more naturalistic datasets containing spontaneous and conversational speech that reflect real-life application scenarios. 
This would help determine whether the observed rhythm, duration, trajectory, and linguistic effects generalize beyond LibriSpeech.

A recently proposed metric that would be especially relevant for future evaluation is Similarity Rank Disclosure (SRD) \citep{chandra2026evaluatingvoiceanonymisationusing}. SRD measures privacy disclosure directly from feature representations by ranking the matching speaker among a set of reference speakers. 
It is independent of classifier thresholds and reports both average and worst-case disclosure in information bits. 
It can also be applied to different representations, including speaker embeddings, fundamental frequency, and phone-based features which is useful because it provides a multidimensional view of privacy, which in this work I tried to do with the feature-level analyses. 
Future work could therefore compare its findings with EER, ZEBRA, Linkability \& Singling Out, test them after performing anonymization on larger, more naturalistic datasets and see if the results align with the trends noticed in this thesis. 

  
%%%%%%%%%%%%%%%%%%%%%%%%%%%%%%%CONFIG%%%%%%%%%%%%%%%%%%%%%%%%%%%%%%%%%%%%%%%
% \section{Formatting Examples}

% \subsection{Lists}\label{listen}

% Lists look like this:

% \begin{itemize}
% \item This 
% \item is 
% \item  a 
% \item list 
% \end{itemize}


% \subsection{Enumerations}

% Lists with numbers are called enumerations:

% \begin{enumerate}
% \item One
% \item Two
% \item Three 
% \item Many
% \end{enumerate}

% In contrast to the list in Section \ref{listen}, the enumeration has numbers.


% \section{Figures, Tables etc.}

% \subsection{Tables}

% Tables can be put in the running text, as shown below.

% \begin{tabular}{clr}
% \toprule
% Column 1 & Column 2 & Column 3\\
% \midrule
% Column 1 & Column 2 & Column 3 \\
% centered & aligned left &  aligned right \\
% \bottomrule
% \end{tabular}

% But tables should be placed as float objects, which gives them an automatic number and caption and places them at reasonable places (see Table \ref{tabelle}).

% \begin{table}
% \begin{center}
% \begin{tabular}{clr}
% \toprule
% Yet & Another & Table \\
% \midrule
% Column 1 & Column 2 & Column 3 \\
% centered & aligned left &  aligned right \\
% \bottomrule
% \end{tabular}
% \caption{Isn't this a nice table?}\label{tabelle}
% \end{center}
% \end{table}

% \subsection{Figures}

% Pictures can be included using the package \texttt{epsfig}, as can be seen in Figure \ref{bild}. Similar to tables used in the \texttt{table}-environment, figures are placed automatically.

% \begin{figure}
% \begin{center}
% \epsfig{file=example.pdf,width=0.5\textwidth}
% \caption{This figure is red} \label{bild}
% \end{center}
% \end{figure}

%%%%%%%%%%%%%

% Bibliographie
\bibliographystyle{abbrvnat}
\bibliography{references}

\clearpage
\appendix
\section{Appendix}
\label{app:appendix}

\subsection{ZEBRA Curves on Partition Sets}
\label{app:zebra-curves}

This section presents additional ZEBRA ECE curves across systems and partition sets. The black curve indicates the zero-evidence reference, while the colored curve shows the calibrated ECE profile. The legend summarizes \texttt{dECE}, maximum absolute LLR, and the ZEBRA evidence category.

\begin{figure}[H]
    \centering

    \begin{subfigure}{0.47\textwidth}
        \centering
        \includegraphics[width=\textwidth]{appendix/ZEBRA-zebra_dev_sttts_f.png}
        \caption{STTTS, dev-f}
        \label{fig:zebra-dev-sttts-f}
    \end{subfigure}
    \hfill
    \begin{subfigure}{0.47\textwidth}
        \centering
        \includegraphics[width=\textwidth]{appendix/ZEBRA-zebra_dev_sttts_m.png}
        \caption{STTTS, dev-m}
        \label{fig:zebra-dev-sttts-m}
    \end{subfigure}

    \caption{ZEBRA ECE curves for STTTS on the development partitions.}
    \label{fig:zebra-dev-sttts}
\end{figure}

%%%%%%%%%

\begin{figure}[H]
    \centering

    \begin{subfigure}{0.47\textwidth}
        \centering
        \includegraphics[width=\textwidth]{appendix/ZEBRA-zebra_test_sttts_f.png}
        \caption{STTTS, test-f}
        \label{fig:zebra-test-sttts-f}
    \end{subfigure}
    \hfill
    \begin{subfigure}{0.47\textwidth}
        \centering
        \includegraphics[width=\textwidth]{appendix/ZEBRA-zebra_test_sttts_m.png}
        \caption{STTTS, test-m}
        \label{fig:zebra-test-sttts-m}
    \end{subfigure}

    \caption{ZEBRA ECE curves for STTTS on the test partitions.}
    \label{fig:zebra-test-sttts}
\end{figure}

%%%%%%%%%

\begin{figure}[H]
    \centering

    \begin{subfigure}{0.47\textwidth}
        \centering
        \includegraphics[width=\textwidth]{appendix/ZEBRA-zebra_nac_dev_f.png}
        \caption{NAC, dev-f}
        \label{fig:zebra-dev-nac-f}
    \end{subfigure}
    \hfill
    \begin{subfigure}{0.47\textwidth}
        \centering
        \includegraphics[width=\textwidth]{appendix/ZEBRA-zebra_dev_nac_m.png}
        \caption{NAC, dev-m}
        \label{fig:zebra-dev-nac-m}
    \end{subfigure}

    \caption{ZEBRA ECE curves for NAC on the development partitions.}
    \label{fig:zebra-dev-nac}
\end{figure}

%%%%%%%

\begin{figure}[H]
    \centering

    \begin{subfigure}{0.47\textwidth}
        \centering
        \includegraphics[width=\textwidth]{appendix/ZEBRA-zebra_test_nac_f.png}
        \caption{NAC, test-f}
        \label{fig:zebra-test-nac-f}
    \end{subfigure}
    \hfill
    \begin{subfigure}{0.47\textwidth}
        \centering
        \includegraphics[width=\textwidth]{appendix/ZEBRA-zebra_test_nac_m.png}
        \caption{NAC, test-m}
        \label{fig:zebra-test-nac-m}
    \end{subfigure}

    \caption{ZEBRA ECE curves for NAC on the test partitions.}
    \label{fig:zebra-test-nac}
\end{figure}

%%%%%%%
\begin{figure}[H]
    \centering

    \begin{subfigure}{0.47\textwidth}
        \centering
        \includegraphics[width=\textwidth]{appendix/ZEBRA-zebra_dev_asrbn_f.png}
        \caption{ASRBN, dev-f}
        \label{fig:zebra-dev-asrbn-f}
    \end{subfigure}
    \hfill
    \begin{subfigure}{0.47\textwidth}
        \centering
        \includegraphics[width=\textwidth]{appendix/ZEBRA-zebra_dev_asrbn_m.png}
        \caption{ASRBN, dev-m}
        \label{fig:zebra-dev-asrbn-m}
    \end{subfigure}

    \caption{ZEBRA ECE curves for ASRBN on the development partitions.}
    \label{fig:zebra-dev-asrbn}
\end{figure}

%%%%%%%
\begin{figure}[H]
    \centering

    \begin{subfigure}{0.47\textwidth}
        \centering
        \includegraphics[width=\textwidth]{appendix/ZEBRA-zebra_test_asrbn_f.png}
        \caption{ASRBN, test-f}
        \label{fig:zebra-test-asrbn-f}
    \end{subfigure}
    \hfill
    \begin{subfigure}{0.47\textwidth}
        \centering
        \includegraphics[width=\textwidth]{appendix/ZEBRA-zebra_test_asrbn_m.png}
        \caption{ASRBN, test-m}
        \label{fig:zebra-test-asrbn-m}
    \end{subfigure}

    \caption{ZEBRA ECE curves for ASRBN on the test partitions.}
    \label{fig:zebra-test-asrbn}
\end{figure}
%%%%%%%%%%%%%%%%%%%%%%%%%%%%%%%%%%%%%%%%%%%%%%%%%%%%%%%%%%%%%%%

\subsection{Speaker-Level Feature Preservation Heatmaps}
\label{app:feature-heatmaps}

This section presents additional speaker-level heatmaps of feature preservation after anonymization. Speakers are sorted by EER, with lower values indicating higher vulnerability. Columns show standardized preservation scores across feature spaces, where warmer colors indicate above-average preservation and cooler colors indicate below-average preservation.

\begin{figure}[H]
    \centering

    \begin{subfigure}{0.95\textwidth}
        \centering
        \includegraphics[height=0.40\textheight, keepaspectratio]{appendix/speaker_feature_heatmap_sorted_by_B3_EER_dev_f.png}
        \caption{B3, dev-f}
        \label{fig:heatmap-b3-dev-f}
    \end{subfigure}

    \vspace{0.3cm}

    \begin{subfigure}{0.95\textwidth}
        \centering
        \includegraphics[height=0.40\textheight, keepaspectratio]{appendix/speaker_feature_heatmap_sorted_by_B3_EER_dev_m.png}
        \caption{B3, dev-m}
        \label{fig:heatmap-b3-dev-m}
    \end{subfigure}

    \caption{Speaker-level preservation features sorted by B3 EER on the development partitions.}
    \label{fig:heatmap-b3-dev}
\end{figure}

%%%%%%%%


\begin{figure}[H]
    \centering

    \begin{subfigure}{0.95\textwidth}
        \centering
        \includegraphics[height=0.40\textheight, keepaspectratio]{appendix/speaker_feature_heatmap_sorted_by_B3_EER_test_f.png}
        \caption{B3, test-f}
        \label{fig:heatmap-b3-test-f}
    \end{subfigure}

    \vspace{0.3cm}

    \begin{subfigure}{0.95\textwidth}
        \centering
        \includegraphics[height=0.40\textheight, keepaspectratio]{appendix/speaker_feature_heatmap_sorted_by_B3_EER_test_m.png}
        \caption{B3, test-m}
        \label{fig:heatmap-b3-test-m}
    \end{subfigure}

    \caption{Speaker-level preservation features sorted by B3 EER on the test partitions.}
    \label{fig:heatmap-b3-test}
\end{figure}

%%%%%%%%%%%%%
\begin{figure}[H]
    \centering

    \begin{subfigure}{0.95\textwidth}
        \centering
        \includegraphics[height=0.40\textheight, keepaspectratio]{appendix/speaker_feature_heatmap_sorted_by_B4_EER_dev_f.png}
        \caption{B4, dev-f}
        \label{fig:heatmap-b4-dev-f}
    \end{subfigure}

    \vspace{0.3cm}

    \begin{subfigure}{0.95\textwidth}
        \centering
        \includegraphics[height=0.40\textheight, keepaspectratio]{appendix/speaker_feature_heatmap_sorted_by_B4_EER_dev_m.png}
        \caption{B4, dev-m}
        \label{fig:heatmap-b4-dev-m}
    \end{subfigure}

    \caption{Speaker-level preservation features sorted by B4 EER on the development partitions.}
    \label{fig:heatmap-b4-dev}
\end{figure}

%%%%%%%%%%
\begin{figure}[H]
    \centering

    \begin{subfigure}{0.95\textwidth}
        \centering
        \includegraphics[height=0.40\textheight, keepaspectratio]{appendix/speaker_feature_heatmap_sorted_by_B4_EER_test_f.png}
        \caption{B4, test-f}
        \label{fig:heatmap-b4-test-f}
    \end{subfigure}

    \vspace{0.3cm}

    \begin{subfigure}{0.95\textwidth}
        \centering
        \includegraphics[height=0.40\textheight, keepaspectratio]{appendix/speaker_feature_heatmap_sorted_by_B4_EER_test_m.png}
        \caption{B4, test-m}
        \label{fig:heatmap-b4-test-m}
    \end{subfigure}

    \caption{Speaker-level preservation features sorted by B4 EER on the test partitions.}
    \label{fig:heatmap-b4-test}
\end{figure}

 
%%%%%%%%%%
\begin{figure}[H]
    \centering

    \begin{subfigure}{0.95\textwidth}
        \centering
        \includegraphics[height=0.40\textheight, keepaspectratio]{appendix/speaker_feature_heatmap_sorted_by_B5_EER_dev_f.png}
        \caption{B5, dev-f}
        \label{fig:heatmap-b5-dev-f}
    \end{subfigure}

    \vspace{0.3cm}

    \begin{subfigure}{0.95\textwidth}
        \centering
        \includegraphics[height=0.40\textheight, keepaspectratio]{appendix/speaker_feature_heatmap_sorted_by_B5_EER_dev_m.png}
        \caption{B5, dev-m}
        \label{fig:heatmap-b5-dev-m}
    \end{subfigure}

    \caption{Speaker-level preservation features sorted by B5 EER on the development partitions.}
    \label{fig:heatmap-b5-dev}
\end{figure}

%%%%%%%%%%
\begin{figure}[H]
    \centering

    \begin{subfigure}{0.95\textwidth}
        \centering
        \includegraphics[height=0.40\textheight, keepaspectratio]{appendix/speaker_feature_heatmap_sorted_by_B5_EER_test_f.png}
        \caption{B5, test-f}
        \label{fig:heatmap-b5-test-f}
    \end{subfigure}

    \vspace{0.3cm}

    \begin{subfigure}{0.95\textwidth}
        \centering
        \includegraphics[height=0.40\textheight, keepaspectratio]{appendix/speaker_feature_heatmap_sorted_by_B5_EER_test_m.png}
        \caption{B5, test-m}
        \label{fig:heatmap-b5-test-m}
    \end{subfigure}

    \caption{Speaker-level preservation features sorted by B5 EER on the test partitions.}
    \label{fig:heatmap-b5-test}
\end{figure}

\subsection{Use of AI-Tools}
\label{app:ai-tools}

The author of this thesis has used AI tools, i.e. ChatGPT  versions GPT5.4 and GPT5.5, in order to assist in the code support when errors were encountered in the implementation of the
evaluation metrics and in order to secure their compliance with their accompanied papers. AI was also used for the analysis part of the thesis, again with support in the implementation of the analysis code and the encountered errors. 
The author has also used this tool in order to compile the tables seen in this thesis and in some cases in order to refine the text. All the other parts of the thesis, including the writing, the background work, the related work, the analysis and the interpretation of the results, were done solely by the author.


\end{document}
